\documentclass[11pt,oneside]{article}

\usepackage[a4paper,margin=26mm,headheight=14pt]{geometry}
\usepackage[T1]{fontenc}
\usepackage[utf8]{inputenc}
\usepackage{lmodern}
%% Unicode characters used inside \lean{...} citations.
\DeclareUnicodeCharacter{2192}{\ensuremath{\rightarrow}}
\DeclareUnicodeCharacter{2295}{\ensuremath{\oplus}}
\DeclareUnicodeCharacter{2090}{\ensuremath{{}_{\mathrm a}}}
\DeclareUnicodeCharacter{2243}{\ensuremath{\simeq}}
\DeclareUnicodeCharacter{2264}{\ensuremath{\le}}
\DeclareUnicodeCharacter{2115}{\ensuremath{\mathbb N}}
\DeclareUnicodeCharacter{222A}{\ensuremath{\cup}}
\DeclareUnicodeCharacter{2208}{\ensuremath{\in}}
\DeclareUnicodeCharacter{2200}{\ensuremath{\forall}}
\DeclareUnicodeCharacter{21A5}{\ensuremath{\uparrow}}
\DeclareUnicodeCharacter{00D7}{\ensuremath{\times}}

\usepackage{microtype}
\usepackage{amsmath,amssymb,amsthm,mathtools}
\usepackage{booktabs,longtable,array,tabularx}
\usepackage{enumitem}
\usepackage{xcolor}
\usepackage{hyperref}
\usepackage{cleveref}
\usepackage{listings}
\usepackage{fancyhdr}

\definecolor{leanblue}{HTML}{0B4F8A}
\definecolor{darkgreen}{HTML}{246B45}
\definecolor{brick}{HTML}{8B2E2E}

\hypersetup{
  colorlinks=true,
  linkcolor=leanblue,
  citecolor=darkgreen,
  urlcolor=leanblue,
  pdftitle={The group configuration theorem for algebraically closed fields: a formalization blueprint},
  pdfauthor={acl\_geom project}
}

\pagestyle{fancy}
\fancyhf{}
\lhead{Group configuration theorem for ACF}
\rhead{Lean/Mathlib blueprint}
\cfoot{\thepage}

\newtheorem{theorem}{Theorem}[section]
\newtheorem{lemma}[theorem]{Lemma}
\newtheorem{proposition}[theorem]{Proposition}
\newtheorem{corollary}[theorem]{Corollary}
\theoremstyle{definition}
\newtheorem{definition}[theorem]{Definition}
\newtheorem{construction}[theorem]{Construction}
\newtheorem{convention}[theorem]{Convention}
\theoremstyle{remark}
\newtheorem{remark}[theorem]{Remark}
\newtheorem{warning}[theorem]{Warning}
\newtheorem{leannote}[theorem]{Lean note}

\newcommand{\acl}{\operatorname{acl}}
\newcommand{\dcl}{\operatorname{dcl}}
\newcommand{\trdeg}{\operatorname{trdeg}}
\newcommand{\Aut}{\operatorname{Aut}}
\newcommand{\rk}{\operatorname{r}}
\newcommand{\Loc}{\operatorname{Loc}}
\newcommand{\FOD}{\operatorname{Fod}}
\newcommand{\Cb}{\operatorname{Cb}}
\newcommand{\Germ}{\operatorname{Germ}}
\newcommand{\Gen}{\operatorname{gen}}
\newcommand{\Fix}{\operatorname{Fix}}
\newcommand{\code}[1]{\ulcorner #1\urcorner}
\newcommand{\indep}{\mathrel{\raisebox{0.15ex}{$\downarrow$}\hspace{-0.63em}\raisebox{0.15ex}{$\smile$}}}
\newcommand{\ind}[1]{\mathrel{\indep_{#1}}}
\newcommand{\interalg}{\sim}
\newcommand{\lean}[1]{\texttt{\color{leanblue}#1}}
\newcommand{\Kv}{\mathbf{K}}

\lstset{
  basicstyle=\ttfamily\small,
  columns=fullflexible,
  keepspaces=true,
  breaklines=true,
  frame=single,
  framerule=0.3pt,
  xleftmargin=4pt,
  xrightmargin=4pt
}

\title{The group configuration theorem for algebraically closed fields\\[1ex]
\large A self-contained formalization blueprint}
\author{acl\_geom project}
\date{October 2026 (draft)}

\begin{document}
\maketitle

\begin{abstract}
This document is a complete, self-contained proof of Hrushovski's group configuration
theorem specialised to algebraically closed fields, in the form needed by the
reconstruction blueprint \texttt{sources/blueprint.tex}: from six tuples forming a
group configuration (three ``group'' points of rank $n$ and three ``space'' points
of rank $m\le n$, lines given by algebraic closure, with the non-degeneracy condition
that each group point is determined by the germ of the correspondence it induces)
one obtains, after adjoining finitely many independent generic parameters, a
connected group $G$ of rank $n$ and a transitive faithful $G$-set $S$ containing the
generic points of a variety of dimension $m$, such that the configuration is pointwise
interalgebraic with a generic configuration $(g,h,hg^{-1};\, g\cdot s,\, s,\, h\cdot s)$
of $(G,S)$.

The proof follows Hrushovski's argument in the exposition of Pillay and Bays
(germs of definable functions, canonical bases, the Hrushovski--Weil theorem),
but every model-theoretic notion is replaced by an explicit field-theoretic one:
types by generic points of irreducible varieties, forking by transcendence degree,
canonical bases by smallest fields of definition, definable functions by
quasi-rational maps, and imaginaries by explicit codes.  The group is produced as
a group of germs of quasi-rational self-maps of a variety, and is shown to be
\emph{constructible} (definable) in the sense of a constructible set modulo a
constructible equivalence relation with constructible operations.  Passing from a
constructible group to an algebraic group (Weil's theorem) is outside the scope of
this document and is not needed for the ``up to interalgebraicity'' statement.
\end{abstract}

\tableofcontents

\section{Scope, conventions, and the shape of the main theorem}
\label{sec:scope}

\subsection{What is proved}

Throughout, $k\subseteq K$ are algebraically closed fields and $K$ has infinite
transcendence degree over $k$ (\cref{conv:ud}).  All tuples are finite tuples of
elements of $K$; $\acl$ is relative algebraic closure in $K$ over $k$, and the
rank $\rk(a/B)$ of a tuple $a$ over a set $B$ is the transcendence degree of
$k(B,a)$ over $k(B)$ (\cref{sec:ud}).

The object of study is a \emph{group configuration} (\cref{def:gc}): six tuples
$a,b,c,x,y,z$ with the four \emph{lines}
\[
  \{a,b,c\},\qquad \{a,x,y\},\qquad \{b,y,z\},\qquad \{c,z,x\},
\]
such that (i) each group point is algebraic over the two others on the group line,
(ii) on each space line the two space points are interalgebraic over the group point,
(iii) every non-collinear triple is independent, and (iv) each group point is
interalgebraic with the \emph{canonical base} of the correspondence it defines
between the two space points on its line (the ``no smaller field of definition''
condition; in the reconstruction blueprint it is clause (iv) of $\Psi$).  Group
points have a common rank $n\ge 1$, space points a common rank $m$, $1\le m\le n$.

\begin{theorem}[main theorem, informal; precise form in \cref{thm:main}]
\label{thm:main-informal}
Let $(a,b,c,x,y,z)$ be a group configuration over $k$.  Then there are finitely many
elements $e$ of $K$, independent from $abcxyz$ over $k$, such that, writing
$k'=\acl(k,e)$ for the new base field, there exist
\begin{itemize}[nosep]
\item a group $G$ of germs of quasi-rational self-maps of an irreducible $k'$-variety
  $V$ of dimension $m$, connected, of rank $n$, whose elements of rank $n$ are exactly
  the germs $f_1^{-1}f_2$ for independent members $f_1,f_2$ of a generic set $\mathcal F$
  of invertible germs from $V$ to a second variety $V'$ (a single $\Gamma$-orbit, hence
  canonically parametrised by the generic points of one variety of codes);
\item a transitive faithful $G$-set $S$ containing the generic points of $V$;
\item generic independent $g,h\in G$ and a generic $s\in V(K)$,
\end{itemize}
such that, with $x'=g\cdot s$, $y'=s$, $z'=h\cdot s$, $a'=g$, $b'=h$, $c'=hg^{-1}$,
each of $a,b,c,x,y,z$ is interalgebraic over $k'$ with the corresponding primed
element (group elements being represented by their canonical codes, which are tuples
of $K$).  Moreover, for a generic point $y_0$ of $V$ over $k'(abcxyz)$ and a further
tuple $t$ of fresh generic points (chosen after $y_0$) with $(e,t)$ independent from
$abcxyz$, $G$ is \emph{constructible over $k'(t)$} (\cref{def:constructible-group}) and
$S$ is a constructible $G$-set over $k'(t,y_0)$, namely $G/H$ for the constructible
stabiliser $H$ of $\iota(y_0)$ (\cref{thm:constructible}).
\end{theorem}

In the application to the reconstruction blueprint (\cref{sec:application}),
$n=2$, $m=1$, and the configuration is the $\Psi$-witness of the geometric
relation $Q$; the theorem supplies the two-dimensional group acting on a curve whose
classification (blueprint Lemma \texttt{affine-action}) yields the affine coordinates
of Lemma \texttt{affine-grid-extraction}.  The classification itself is not part of
this document.

\subsection{What is not proved here}

\begin{itemize}
\item \textbf{Weil's theorem.}  We do not prove that a constructible group over an
  algebraically closed field is (isomorphic to) an algebraic group in the sense of
  algebraic geometry, nor do we construct a variety structure on $G$.  The main
  theorem is an ``up to interalgebraicity'' statement and every quantity in it is
  defined in terms of fields, ideals of polynomials and constructible sets.
  \Cref{sec:constructibility} records the precise definable-group structure that
  is obtained, which is the input any later algebraization would start from.
\item \textbf{Classification} of the resulting group (dimension two acting on a
  curve, blueprint Lemma \texttt{affine-action}).
\item \textbf{Elimination of imaginaries in full generality.}  Only the codes that
  are actually needed are constructed explicitly (\cref{sec:codes}); the
  constructible group of \cref{sec:constructibility} is presented as a quotient of a
  constructible set by a constructible equivalence relation, and we do not prove
  that such a quotient can be identified with a constructible subset of some $K^N$.
\end{itemize}

\subsection{How the proof is organised}

\begin{enumerate}
\item \Cref{sec:ud}: the universal domain, rank calculus, automorphisms, loci.
\item \Cref{sec:fod}: smallest fields of definition of ideals and the resulting
  \emph{canonical bases} $\Cb(a/L)$; their Galois and independence properties.
\item \Cref{sec:codes}: definable closure as perfect closure; explicit codes of
  finite sets of tuples.
\item \Cref{sec:germs}: germs of quasi-rational maps between irreducible varieties,
  their evaluation at generic points, composition, inverses, and the dictionary
  between a triple $(b,x,y)$ with $y\in\dcl(b,x)$ and a germ whose code is
  $\Cb(xy/b)$.
\item \Cref{sec:hw}: the Hrushovski--Weil theorem: a generically transitive generic
  set of invertible germs (a single orbit under the automorphism group) closed under
  inverse and generic composition generates a connected group of germs whose elements
  of maximal rank are exactly that set, and a transitive faithful homogeneous space.
\item \Cref{sec:bij}: the version for families of bijections between two varieties
  (composites $f_{b_1}^{-1}f_{b_2}$), with the rank criterion for the independence
  hypothesis.
\item \Cref{sec:gc}: group configurations, their invariance under base change and
  interalgebraic replacement, and the converse (configurations of a homogeneous space).
\item \Cref{sec:proof}: the three steps of Hrushovski's proof (making the lines
  definable, the independence hypothesis, identification) and the main theorem.
\item \Cref{sec:constructibility}: the group and the space are constructible
  (definable) over the extended base.
\item \Cref{sec:descent}: descent of the extra parameters into $k$ for first-order
  consequences.
\item \Cref{sec:application}: the $\Psi$-witness is a group configuration; what the
  theorem gives for the reconstruction blueprint.
\item \Cref{sec:lean}: proposed Lean architecture and the Mathlib/AclGeom API used.
\end{enumerate}

\subsection{Conventions}

\begin{convention}[orientation]\label{conv:orientation}
We use the orientation of Bays' notes, which coincides with the letters of the
reconstruction blueprint's $\Psi$-witness: the group line is $\{a,b,c\}$ with
\emph{$b=c\ast a$}, and the space lines read
\[
   x=a\ast y,\qquad z=b\ast y,\qquad z=c\ast x .
\]
Thus $a$ carries $y$ to $x$, $c$ carries $x$ to $z$, and $b=c\ast a$ carries $y$ to
$z$.  The independent triple of the configuration is $(a,b,y)$.  In the
$\Psi$-witness of the reconstruction blueprint, $(a,b,c,x,y,z)=(A,B,C,X,Y,Z)$:
$X\le A\vee Y$, $Z\le B\vee Y$, $Z\le C\vee X$ (\cref{sec:application}).
\end{convention}

\begin{convention}[citations to Lean]\label{conv:lean}
Statements that are already available in Mathlib (revision
\lean{d13f23b7}, tag \lean{v4.34.1}) or in the project library \lean{AclGeom} are
marked with their declaration names in \lean{teletype}.  Such statements are
\emph{not} reproved here.  Every other lemma is proved in full.  The Lean names
used are listed in \cref{sec:lean}; when a cited name has a slightly different
form from the statement in the text (for instance a lattice formulation versus a
tuple formulation), the required bridge is stated explicitly there.
\end{convention}

\begin{convention}[characteristic]\label{conv:char}
The characteristic $p\ge 0$ of $k$ is arbitrary.  In characteristic $p>0$, the
definable closure of a field is its perfect closure and ``definable functions'' are
quasi-rational (rational after extracting $p$-power roots).  No separability
hypothesis is ever used.  All constructions are uniform in $p$; in characteristic
$0$ every ``quasi-rational'' map is rational.
\end{convention}

\section{The universal domain}
\label{sec:ud}

\subsection{Setting and closure operations}

\begin{convention}[the universal domain]\label{conv:ud}
Fix algebraically closed fields $k\subseteq K$ such that $\trdeg(K/k)$ is infinite.
A \emph{tuple} is a finite sequence of elements of $K$; tuples are concatenated
freely ($ab$ denotes the concatenation of $a$ and $b$) and a finite set of tuples
is identified with the concatenation of its members.  For a set $A\subseteq K$ we
write $k(A)$ for the subfield generated by $k\cup A$.  A subfield $L\subseteq K$ is
\emph{small} if $L\supseteq k$ and $\trdeg(L/k)<\infty$.
\end{convention}

\begin{leannote}
In Lean the ambient field is a type \lean{K} with \lean{[Field K] [IsAlgClosed K]}
and a base \lean{k} with \lean{[Field k] [IsAlgClosed k] [Algebra k K]}; the
hypothesis on the transcendence degree is \lean{Cardinal.aleph0 ≤ Algebra.trdeg k K}.
Tuples are functions \lean{Fin n → K}.  In the reconstruction blueprint the ambient
field is only assumed to have $\trdeg\ge5$; \cref{rem:finite-trdeg} and
\cref{rem:psi-hyp}(c) explain how the present results are applied there (by passing to
an ambient extension of infinite transcendence degree).
\end{leannote}

\begin{definition}[closures, rank, independence]\label{def:closures}
Let $A,B,C\subseteq K$ be sets (in practice finite tuples or small fields).
\begin{enumerate}[nosep]
\item $\acl(A)$ is the relative algebraic closure of $k(A)$ in $K$ (\lean{AclGeom.racl}).
  It is an algebraically closed subfield of $K$ containing $k$.
\item $\dcl(A)$ is the \emph{perfect closure} of $k(A)$ in $K$: in characteristic
  $0$, $\dcl(A)=k(A)$; in characteristic $p>0$,
  $\dcl(A)=\{x\in K:\ x^{p^e}\in k(A)\text{ for some }e\ge 0\}$.
\item $\rk(A/C)=\trdeg\bigl(k(C,A)/k(C)\bigr)\in\mathbb N\cup\{\infty\}$; it is finite
  when $A$ is finite.  We write $\rk(A)=\rk(A/\emptyset)$.
\item $A\ind{C}B$ (``$A$ is independent from $B$ over $C$'') means
  $\rk(a/C\cup B)=\rk(a/C)$ for every finite tuple $a$ from $A$.
\item $A$ and $B$ are \emph{interalgebraic}, $A\interalg B$, if $\acl(A)=\acl(B)$; they
  are \emph{interalgebraic over $C$} if $\acl(CA)=\acl(CB)$; they are
  \emph{interdefinable over $C$} if $\dcl(CA)=\dcl(CB)$.
\end{enumerate}
\end{definition}

\begin{lemma}[rank calculus]\label{lem:rank}
For finite tuples $a,a',b,d,e$ and sets $B,C,D$ (independence of sets is as in
\cref{def:closures}(4): $B\ind{C}D$ iff $b\ind{C}D$ for every finite tuple $b$ from $B$):
\begin{enumerate}[label=\textup{(R\arabic*)},start=0]
\item\label{R:sub} (finite character, sub-tuples) $\rk(a/C)=\rk(a/C_0)$ for some finite
  $C_0\subseteq C$, and $a\subseteq\acl(C)$ iff $a\subseteq\acl(C_0)$ for some finite
  $C_0\subseteq C$; if $ab\ind{C}D$ then $a\ind{C}D$; $a\ind{C}D$ iff $a\ind{C}d$ for
  every finite tuple $d$ from $D$.
\item\label{R:add} (additivity) $\rk(ab/C)=\rk(a/C)+\rk(b/Ca)$.
\item\label{R:acl} $\rk(a/C)=\rk(a/\acl(C))$, and $\rk(a/C)=0$ iff $a\subseteq\acl(C)$.
\item\label{R:mono} if $C\subseteq C'$ then $\rk(a/C')\le\rk(a/C)$.
\item\label{R:sym} $a\ind{C}b$ iff $b\ind{C}a$.
\item\label{R:trans} $a\ind{C}bd$ iff $a\ind{C}b$ and $a\ind{Cb}d$.
\item\label{R:aclinv} if $\acl(Ca')=\acl(Ca)$ then $a\ind{C}B \iff a'\ind{C}B$; and
  $a\ind{C}B\iff a\ind{C}\acl(CB)\iff \acl(Ca)\ind{C}\acl(CB)$ (independence of the
  closures, in the sense of \cref{def:closures}(4)).
\item\label{R:meet} if $B\ind{C}D$ then $\acl(CB)\cap\acl(CD)=\acl(C)$.
\item\label{R:selfind} if $B\ind{C}B$ then $B\subseteq\acl(C)$.
\item\label{R:swap} if $b\ind{C}D$ and $a\ind{C}bD$ then $ab\ind{C}D$; if
  $a\ind{C}b$ and $ab\ind{C}D$ then $a\ind{C}bD$.
\item\label{R:basechange} if $e\ind{}abC$ then $a\ind{C}b\iff a\ind{Ce}b$, and
  $a\subseteq\acl(C)\iff a\subseteq\acl(Ce)$.
\end{enumerate}
\end{lemma}

\begin{proof}
\ref{R:sub}: $k(C,a)$ is generated over $k(C)$ by finitely many elements, and a
transcendence basis of $k(C,a)/k(C)$ together with the finitely many algebraic
relations among $a$ and that basis involve only finitely many elements $C_0$ of $C$;
hence $\rk(a/C)=\rk(a/C_0)$, and likewise an algebraic dependence of each $a_i$ on $C$
involves finitely many elements of $C$.  Sub-tuples: $\rk(a/CD)\le\rk(a/C)$ and
$\rk(b/CaD)\le\rk(b/Ca)$ by \ref{R:mono}, and the sum of the two is
$\rk(ab/CD)=\rk(ab/C)$ by \ref{R:add}, so both are equalities.  The last clause is
\ref{R:sub} for $\rk(a/CD)=\rk(a/CD_0)$ with $D_0$ finite.
\ref{R:add} is the tower law for transcendence degree
(\lean{trdeg\_add\_eq}, or \lean{lift\_trdeg\_add\_eq} for the universe-lifted form)
applied to $k(C)\subseteq k(C,a)\subseteq k(C,a,b)$.
\ref{R:acl}: $k(C,a)$ and $\acl(C)(a)$ have the same transcendence degree over $k(C)$
and over $\acl(C)$ respectively because $\acl(C)/k(C)$ is algebraic; the second clause
is the definition of algebraic closure.  \ref{R:mono} is monotonicity of transcendence
degree in the base.  \ref{R:sym}: by \ref{R:add}, $\rk(ab/C)=\rk(a/C)+\rk(b/Ca)
=\rk(b/C)+\rk(a/Cb)$, so $\rk(b/Ca)=\rk(b/C)$ iff $\rk(a/Cb)=\rk(a/C)$.
\ref{R:trans}: $\rk(a/Cbd)\le\rk(a/Cb)\le\rk(a/C)$ with equality throughout iff
both inequalities are equalities.  \ref{R:aclinv}: the first clause follows from \ref{R:acl} and \ref{R:add}
applied to $\rk(aa'/C\cdot)$ (for each finite $b$ from $B$, by \ref{R:sub}); for the
second, $\rk(a/CB)=\rk(a/\acl(CB))$ by \ref{R:acl}, and for a finite tuple $b'$ from
$\acl(CB)$ there is a finite $b$ from $B$ with $b'\subseteq\acl(Cb)$ (\ref{R:sub}), so
$a\ind{C}b'$ follows from $a\ind{C}b$ and the first clause.  \ref{R:meet}: let $d$ be
a finite tuple from $\acl(CB)\cap\acl(CD)$.  Choose finite $b$ from $B$ with
$d\subseteq\acl(Cb)$ (\ref{R:sub}); then $b\ind{C}D$ and \ref{R:aclinv} give $d\ind{C}D$,
so $\rk(d/C)=\rk(d/CD)=0$ and $d\subseteq\acl(C)$.  \ref{R:selfind}: for $b$ from $B$,
$\rk(b/C)=\rk(b/CB)=0$.
\ref{R:swap}: first claim: $\rk(ab/CD)=\rk(b/CD)+\rk(a/CbD)=\rk(b/C)+\rk(a/C)=
\rk(ab/C)$ using \ref{R:trans} for $a$.  Second claim: $\rk(a/CbD)=\rk(ab/CD)-\rk(b/CD)
\ge\rk(ab/C)-\rk(b/C)=\rk(a/C)$, and the reverse inequality is \ref{R:mono}.
\ref{R:basechange}: $e\ind{}abC$ gives $e\ind{C}ab$ (\ref{R:trans}), hence $ab\ind{C}e$
(\ref{R:sym}) and then $a\ind{Cb}e$, $a\ind{C}e$ by \ref{R:trans}; these say
$\rk(a/Ceb)=\rk(a/Cb)$ and $\rk(a/Ce)=\rk(a/C)$, which gives the first
equivalence; the second is the case $b=\emptyset$ with \ref{R:acl}.
\end{proof}

\begin{leannote}
The relative closure $\acl$ is \lean{AclGeom.racl}; the needed facts are
\lean{AclGeom.mem\_racl\_iff}, \lean{AclGeom.racl\_exchange},
\lean{AclGeom.mem\_racl\_adjoin\_base\_iff} (change of base field),
\lean{AclGeom.mem\_racl\_of\_mem\_racl\_insert} (the non-degenerate case $|B|=|D|=1$,
$b\notin\acl(D)$, of \ref{R:meet})
and \lean{AclGeom.algebraicIndependent\_iff\_forall\_notMem\_racl}.  Ranks of closed
subfields are available as \lean{AclGeom.RankEq} with the matroid bridge
\lean{AclGeom.rankEq\_iff\_eRk}; the tuple-level rank $\rk(a/C)$ is
\lean{Algebra.trdeg (IntermediateField.adjoin k C) (IntermediateField.adjoin k (C ∪ a))}
or, equivalently, the matroid rank of $a$ relative to $C$ in
\lean{AlgebraicIndependent.matroid k K}.  Lemma~\ref{lem:rank} should be proved once
for the matroid rank (submodularity and the closure axioms are available in Mathlib's
\lean{Matroid} API) and transported.
\end{leannote}

\subsection{Automorphisms}

Let $\Gamma=\Aut(K/k)$ and, for a subfield $L\subseteq K$, $\Gamma_L=\Aut(K/L)$.
Every $\sigma\in\Gamma$ acts on tuples, sets, subfields and polynomial ideals
coordinatewise.

\begin{lemma}[fresh elements]\label{lem:fresh}
For every small subfield $L$ and every $n$ there is a tuple $t\in K^n$ with
$\rk(t/L)=n$.
\end{lemma}
\begin{proof}
$\trdeg(K/L)=\trdeg(K/k)$ is infinite since $\trdeg(L/k)$ is finite
(\lean{trdeg\_add\_eq}); a transcendence basis of $K$ over $L$ has at least $n$
elements.  (\lean{AclGeom.fresh\_four\_of\_five\_le\_trdeg} is the special case used in
the reconstruction blueprint.)
\end{proof}

\begin{lemma}[extension of isomorphisms; conjugacy]\label{lem:conj}
Let $L\subseteq K$ be any subfield and let $a,a'\in K^n$ be tuples with the same
vanishing ideal over $L$:
$\{f\in L[X_1,\dots,X_n]: f(a)=0\}=\{f: f(a')=0\}$.  Then there is
$\sigma\in\Gamma_L$ with $\sigma(a)=a'$.
\end{lemma}
\begin{proof}
This is \lean{AclGeom.exists\_ambientAutomorphism\_of\_idealOf\_eq}, applied with base
field $L$: the two vanishing ideals agree, so $L(a)\cong L(a')$ over $L$ by
$a\mapsto a'$; both are finitely generated over $L$, so $\trdeg(K/L(a))=\trdeg(K/L(a'))$,
and the isomorphism extends along transcendence bases and then to the algebraic
closure $K$ (\lean{IsAlgClosure.equivOfEquiv}).
\end{proof}

\begin{lemma}[equivariance]\label{lem:equiv}
For $\sigma\in\Gamma$: $\sigma(\acl A)=\acl(\sigma A)$, $\sigma(\dcl A)=\dcl(\sigma A)$,
$\rk(\sigma a/\sigma C)=\rk(a/C)$, and $a\ind{C}B\iff\sigma a\ind{\sigma C}\sigma B$.
\end{lemma}
\begin{proof}
$\sigma$ is a field automorphism fixing $k$; all notions are defined from the field
structure and $k$.
\end{proof}

\begin{lemma}[finite orbits]\label{lem:orbit}
Let $L\subseteq K$ be any subfield containing $k$ and $x\in K$.  The orbit
$\Gamma_L\cdot x$ is finite iff $x\in\acl(L)$ (the relative algebraic closure of $L$ in
$K$).  If $x$ is transcendental over $L$, then $\Gamma_L\cdot x$ is the set of all
elements of $K$ transcendental over $L$.
\end{lemma}
\begin{proof}
If $x$ is algebraic over $L$, every $\sigma x$ is a root of the minimal polynomial of
$x$ over $L$, so the orbit is finite.  If $x$ is transcendental over $L$, then for any
$x'$ transcendental over $L$ the vanishing ideals of $x$ and $x'$ over $L$ are both
zero, and \cref{lem:conj} gives $\sigma\in\Gamma_L$ with $\sigma x=x'$; the elements
$x,x^2,x^3,\dots$ are pairwise distinct and transcendental over $L$, so the orbit is
infinite.
\end{proof}

\subsection{Loci, generic points, dimension}

\begin{definition}[loci]\label{def:locus}
Let $a\in K^n$ and let $L\subseteq K$ be a subfield.  The \emph{vanishing ideal} of $a$
over $L$ is $I_L(a)=\{f\in L[X]: f(a)=0\}$, a prime ideal of $L[X]=L[X_1,\dots,X_n]$
(\lean{AclGeom.idealOf}, \lean{AclGeom.idealOf\_isPrime}; equivalently
\lean{MvPolynomial.vanishingIdeal L \{a\}}).  If $L$ is algebraically closed, the
\emph{locus ideal} of $a$ over $L$ is the extended ideal
\[
  P_L(a)=I_L(a)\cdot K[X]\subseteq K[X].
\]
An ideal $I\subseteq K[X]$ is \emph{defined over} a subfield $F\subseteq K$ if
$I=(I\cap F[X])\cdot K[X]$.  A tuple $a'\in K^n$ is a \emph{generic point of a prime $P$
over an algebraically closed subfield $M$} if $P_M(a')=P$.
\end{definition}

\begin{lemma}[extension and contraction]\label{lem:extcontr}
Let $F\subseteq M\subseteq K$ be subfields and $J\subseteq F[X]$ an ideal.
\begin{enumerate}[nosep]
\item (change of coefficients) $(J\cdot K[X])\cap M[X]=J\cdot M[X]$ and
  $(J\cdot K[X])\cap F[X]=J$.
\item (dropping variables) If $J\subseteq M[Y,Z]$ is an ideal in two groups of variables,
  then $(J\cdot K[Y,Z])\cap K[Y]=(J\cap M[Y])\cdot K[Y]$; in particular an ideal of
  $K[Y,Z]$ defined over $M$ contracts to an ideal of $K[Y]$ defined over $M$.
\end{enumerate}
\end{lemma}
\begin{proof}
(1) Choose an $M$-basis $(\beta_i)$ of $K$ containing $1$.  Then $K[X]=\bigoplus_i M[X]\beta_i$
as $M[X]$-modules and $J\cdot K[X]=(J\cdot M[X])\cdot K[X]=\bigoplus_i (J\cdot M[X])\beta_i$;
intersecting with $M[X]$ (the summand of $\beta=1$) gives $J\cdot M[X]$.  The second
claim is the case $M=F$.  (\lean{MvPolynomial.ker\_map} and
\lean{Ideal.comap\_map\_eq\_self\_of\_faithfullyFlat} are the Mathlib ingredients:
$K[X]$ is faithfully flat over $M[X]$.)
(2) With the same basis, $K[Y,Z]=\bigoplus_i M[Y,Z]\beta_i$ and
$K[Y]=\bigoplus_i M[Y]\beta_i$; so $J\cdot K[Y,Z]=\bigoplus_i J\beta_i$ and an element
$\sum_i h_i\beta_i$ ($h_i\in J$) lies in $K[Y]$ iff every $h_i\in M[Y]$, i.e.\ iff every
$h_i\in J\cap M[Y]$.
\end{proof}

\begin{lemma}[locus ideals are prime]\label{lem:locusprime}
If $L$ is algebraically closed then $P_L(a)$ is a prime ideal of $K[X]$, and
$P_L(a)\cap L[X]=I_L(a)$.  More generally, for algebraically closed $L\subseteq M\subseteq K$
and a prime $I\subseteq L[X]$, the extension $I\cdot M[X]$ is prime.
\end{lemma}
\begin{proof}
\lean{AclGeom.isPrime\_map\_of\_isAlgClosed} (the extension of a prime ideal from an
algebraically closed field along any field extension is prime) and
\cref{lem:extcontr}.
\end{proof}

\begin{lemma}[existence of generic points]\label{lem:genexists}
Let $M$ be a small algebraically closed subfield and $P\subseteq K[X]$ a prime ideal
defined over $M$.  Then $P$ has a generic point over $M$.  Any two generic points of
$P$ over $M$ are conjugate under $\Gamma_M$, and every $\Gamma_M$-conjugate of a generic
point is generic.  Moreover (generic over a larger base implies generic over a smaller
one): if $F\subseteq M$ is algebraically closed, $P$ is defined over $F$, and
$P_M(a')=P$, then $P_F(a')=P$.
\end{lemma}
\begin{proof}
$Q=P\cap M[X]$ is prime and $P=Q\cdot K[X]$.  The domain $M[X]/Q$ is finitely
generated over $M$; its fraction field $E$ has finite transcendence degree $d$ over $M$.
Choose a transcendence basis $t_1,\dots,t_d$ of $E/M$ and elements $s_1,\dots,s_d\in K$
with $\rk(s/M)=d$ (\cref{lem:fresh}); the $M$-embedding $M(t)\to K$, $t_i\mapsto s_i$,
extends to $E$ since $E/M(t)$ is algebraic and $K$ is algebraically closed
(\lean{AclGeom.adjoinTranscendentalAlgHom},
\lean{AclGeom.exists\_extension\_of\_isAlgebraic}, or \lean{IsAlgClosed.lift}).  The image
$a'$ of the residue classes of $X_1,\dots,X_n$ satisfies $I_M(a')=Q$, i.e.\ $P_M(a')=P$.
If $a',a''$ are two generic points then $I_M(a')=I_M(a'')=Q$ and \cref{lem:conj} applies;
conversely $I_M(\sigma a')=\sigma(I_M(a'))=I_M(a')$ for $\sigma\in\Gamma_M$.
For the last clause, $P_F(a')\subseteq P_M(a')=P$ since $I_F(a')\subseteq I_M(a')$;
conversely $P\cap F[X]\subseteq P\cap M[X]=I_M(a')$ (\cref{lem:extcontr}), so
$P\cap F[X]\subseteq I_M(a')\cap F[X]=I_F(a')$ and
$P=(P\cap F[X])K[X]\subseteq I_F(a')K[X]=P_F(a')$.
\end{proof}

\begin{lemma}[dimension drop]\label{lem:dimdrop}
Let $M$ be any field and $Q\subseteq Q'$ prime ideals of $M[X]$.  Then
$\trdeg_M\bigl(\operatorname{Frac}(M[X]/Q')\bigr)\le\trdeg_M\bigl(\operatorname{Frac}(M[X]/Q)\bigr)$,
with strict inequality if $Q\subsetneq Q'$.
\end{lemma}
\begin{proof}
Write $A=M[X]/Q$, $A'=M[X]/Q'$, $\pi:A\to A'$ the surjection, $d'=\trdeg_M A'$ and
$d=\trdeg_M A$ (transcendence degrees of fraction fields).  Preimages of algebraically
independent elements are algebraically independent, so $d'\le d$.  Now let
$Q\subsetneq Q'$ and suppose $d'=d$.
Choose $t_1,\dots,t_d\in A'$ algebraically independent over $M$ and lifts $s_i\in A$;
then $s_1,\dots,s_d$ form a transcendence basis of $\operatorname{Frac}A$ over $M$.
Pick $f\in Q'\setminus Q$; its class $\bar f\in A$ is nonzero and $\pi(\bar f)=0$.
Since $\bar f$ is algebraic over $M(s)$ there is a nonzero $g\in M[S_1,\dots,S_d,T]$
with $g(s,\bar f)=0$; take $g$ of least $T$-degree.  Its constant term $g_0(S)$ is
nonzero: otherwise $g=T\cdot g_1$, and $\bar f g_1(s,\bar f)=0$ in the domain $A$ with
$\bar f\ne 0$ gives $g_1(s,\bar f)=0$ with smaller $T$-degree.  Applying $\pi$,
$0=\pi(g(s,\bar f))=g(t,0)=g_0(t)$, contradicting the algebraic independence of $t$.
\end{proof}

\begin{lemma}[dimension of a locus]\label{lem:dim}
Let $L\subseteq M$ be small algebraically closed subfields and $a\in K^n$.
\begin{enumerate}[nosep]
\item If $a'$ is a generic point of $P_L(a)$ over $M$ then $\rk(a'/M)=\rk(a/L)$.
\item (locus under an independent base extension) $P_M(a)=P_L(a)$ iff $a\ind{L}M$.
\item $P_L(a)=P_L(a')$ iff $a$ and $a'$ are $\Gamma_L$-conjugate.
\end{enumerate}
Consequently the \emph{dimension} $\dim P:=\rk(a'/M)$ of a prime $P\subseteq K[X]$
defined over a small algebraically closed $M$, computed with any generic point $a'$
over any such $M$, is well defined, and $\dim P_L(a)=\rk(a/L)$.
\end{lemma}
\begin{proof}
(1) By \cref{lem:locusprime,lem:extcontr}, $I_M(a')=P_L(a)\cap M[X]=I_L(a)\cdot M[X]$.
This says that $a'$ is a generic point over $M$ of the extension of the prime
$I_L(a)\subseteq L[X]$ to $M$; the theorem
\lean{AclGeom.trdeg\_adjoin\_of\_isGeneric} (with base $L$, extension $M$ and ideal
$I_L(a)$) gives $\rk(a'/M)=\trdeg_M M(a')=\trdeg_L L(a')$, and
\lean{AclGeom.isGeneric\_baseChange} gives $I_L(a')=I_L(a)$, so by \cref{lem:conj} $a'$
and $a$ are $\Gamma_L$-conjugate and $\trdeg_L L(a')=\rk(a/L)$.

(2) Always $P_L(a)\subseteq P_M(a)$ since $I_L(a)\subseteq I_M(a)$.  Both are prime
ideals; by (1) applied to a generic point $a'$ of $P_L(a)$ over $M$ and to $a$ itself,
$\trdeg_M\operatorname{Frac}(M[X]/(P_L(a)\cap M[X]))=\rk(a/L)$ and
$\trdeg_M\operatorname{Frac}(M[X]/(P_M(a)\cap M[X]))=\rk(a/M)$.  If $a\ind{L}M$ the
two numbers are equal and \cref{lem:dimdrop} forces $P_L(a)\cap M[X]=P_M(a)\cap M[X]$,
hence $P_L(a)=P_M(a)$ by \cref{lem:extcontr}.  Conversely if $P_L(a)=P_M(a)$ the two
numbers agree, i.e.\ $\rk(a/M)=\rk(a/L)$.

(3) is \cref{lem:conj} with the base $L$ together with \cref{lem:extcontr}
($P_L(a)=P_L(a')$ iff $I_L(a)=I_L(a')$).  The final statement follows from (1)
applied over $\acl(M\cup M')$ for two choices $M,M'$ of base.
\end{proof}

\begin{lemma}[conjugates independent from a given tuple]\label{lem:fresh-conj}
Let $M$ be a small algebraically closed subfield, $a$ a tuple and $w$ any tuple.  Then
there is $\sigma\in\Gamma_M$ with $\sigma a\ind{M}w$.  Consequently, if $\Phi$ is a
$\Gamma_M$-invariant property of tuples and $\Phi(a)$ holds, there is $a'$ with
$\Phi(a')$ and $a'\ind{M}w$.
\end{lemma}
\begin{proof}
Let $M'=\acl(M,w)$ and let $a'$ be a generic point of $P_M(a)$ over $M'$
(\cref{lem:genexists}).  By \cref{lem:dim}(1), $\rk(a'/M')=\rk(a/M)$; by the last
clause of \cref{lem:genexists}, $P_M(a')=P_M(a)$, so $a'=\sigma a$ for some
$\sigma\in\Gamma_M$ (\cref{lem:dim}(3)) and $\rk(a'/M)=\rk(a/M)=\rk(a'/M')$, i.e.\
$a'\ind{M}w$ (\ref{R:acl}).
\end{proof}

\begin{remark}[the infinite transcendence degree is essential]\label{rem:finite-trdeg}
The hypothesis $\trdeg(K/k)=\infty$ of \cref{conv:ud} is used in two ways.  Explicitly,
the proofs adjoin fresh generic points: the parameter tuple $e$ of \cref{thm:main} has
rank $5n$ (\cref{rem:params}), \cref{thm:constructible} adds a generic point $y_0$ and the
sampling tuple $t=(t_0,t_1)$, and the auxiliary lemmas (\cref{lem:genexists},
\cref{lem:fresh}, \cref{lem:fresh-conj}) adjoin further generic points in a nested way.
Implicitly, and more seriously, the definitions quantify over all \emph{small}
algebraically closed subfields: if $\trdeg(K/k)$ were finite, $K$ itself would be small,
a variety of positive dimension would have no generic point over $K$, and the pair
$(K,f)$ would be a representative of a germ vacuously, equivalent to every other
representative; the germ quotient of \cref{def:germ} would collapse.  So the results of
this document are \emph{not} claimed for ambient fields of finite transcendence degree,
and the parameter count $5n$ does not control the fields quantified over in the
definitions.  For an algebraically closed $K$ of finite transcendence degree over $k$
one passes to an algebraically closed extension $K'$ of infinite transcendence degree,
applies the theorems there, and then needs two separate descents: the parameter
descent of \cref{sec:descent}, which is proved here and moves the adjoined parameters
into $k$, and a bounded first-order (model completeness) argument, not proved here,
which returns existential witnesses from $K'$ to $K$ (\cref{rem:psi-hyp}(c)).
\end{remark}

\section{Fields of definition and canonical bases}
\label{sec:fod}

Canonical bases of types are, in algebraically closed fields, the smallest fields of
definition of loci.  We construct them by linear algebra (reduced row echelon
bases), which is the most economical route for a formalization: it needs no
Gr\"obner bases and no Galois descent.

\subsection{Subspaces of a coordinate space}

Fix a finite set $\mathcal M$ (the index set of coordinates; later $\mathcal M$ is a set
of monomials) and consider $K$-subspaces $W\subseteq K^{\mathcal M}$.  Fix a total order
on $\mathcal M$.

\begin{definition}[defined over a subfield]\label{def:subspace-defined}
Let $F\subseteq K$ be a subfield.  $W$ is \emph{defined over $F$} if $W$ is spanned
by $W\cap F^{\mathcal M}$; equivalently (by comparing dimensions over $F$ and $K$) if
$W=(W\cap F^{\mathcal M})\otimes_F K$ under the natural map.
\end{definition}

\begin{definition}[reduced row echelon basis]\label{def:rref}
A finite sequence $w_1,\dots,w_r$ of vectors of $K^{\mathcal M}$ is in \emph{reduced
row echelon form} if, writing $\mu_i\in\mathcal M$ for the first index (in the fixed
order) at which $w_i$ is nonzero (the \emph{pivot}), we have $\mu_1<\mu_2<\dots<\mu_r$,
$w_i(\mu_i)=1$ and $w_j(\mu_i)=0$ for all $j\ne i$.
\end{definition}

\begin{lemma}[existence and uniqueness]\label{lem:rref}
Every subspace $W\subseteq K^{\mathcal M}$ has exactly one basis in reduced row echelon
form, the \emph{RREF basis} $\operatorname{rref}(W)$.  For an automorphism $\sigma$ of the
field $K$, $\operatorname{rref}(\sigma W)=\sigma(\operatorname{rref}(W))$.  The same holds
with $K$ replaced by any field.
\end{lemma}
\begin{proof}
Existence: Gaussian elimination on any basis.  Uniqueness: let $(w_i)$ and $(w'_j)$ be
two RREF bases.  The set of pivots of an RREF basis of $W$ is determined by $W$: it is
the set of $\mu\in\mathcal M$ such that some vector of $W$ has pivot $\mu$ (if $v\in W$
has pivot $\mu$, write $v=\sum c_i w_i$; the first $i$ with $c_i\neq0$ has
$\mu_i=\mu$; conversely each $w_i$ has pivot $\mu_i$).  So both bases have the same
pivots $\mu_1<\dots<\mu_r$.  Then $w_i-w'_i\in W$ vanishes at every pivot, so it is
zero (a nonzero vector of $W$ has a pivot).  Equivariance: $\sigma$ preserves the
property ``in reduced row echelon form''.
\end{proof}

\begin{definition}[field of definition of a subspace]\label{def:fod-subspace}
$\FOD(W)$ is the subfield of $K$ generated by the prime field and all coordinates
of all vectors of $\operatorname{rref}(W)$.
\end{definition}

\begin{lemma}\label{lem:fod-subspace}
Let $F\subseteq K$ be a subfield.  $W$ is defined over $F$ iff $\FOD(W)\subseteq F$.
Hence $\FOD(W)$ is the smallest field of definition of $W$, it is finitely generated
over the prime field, and $\sigma(\FOD(W))=\FOD(\sigma W)$ for every automorphism
$\sigma$ of $K$.  In particular, for $\sigma\in\Aut(K)$: $\sigma W=W$ iff $\sigma$ fixes
$\FOD(W)$ pointwise.
\end{lemma}
\begin{proof}
If $\FOD(W)\subseteq F$, the RREF basis lies in $F^{\mathcal M}$, so $W$ is defined over
$F$.  Conversely, if $W$ is spanned by $W\cap F^{\mathcal M}$, Gaussian elimination over
$F$ of an $F$-basis of $W\cap F^{\mathcal M}$ produces an RREF sequence in
$F^{\mathcal M}$ which spans $W$ over $K$; by uniqueness it is $\operatorname{rref}(W)$,
so $\FOD(W)\subseteq F$.  Equivariance is \cref{lem:rref}.  If $\sigma W=W$ then $\sigma$
maps $\operatorname{rref}(W)$ to $\operatorname{rref}(W)$ as a sequence, hence fixes each
coordinate; the converse is clear.
\end{proof}

\subsection{Ideals}

Let $X=(X_1,\dots,X_n)$ and, for $d\ge0$, let $K[X]_{\le d}$ be the subspace of
polynomials of total degree $\le d$, identified with $K^{\mathcal M_d}$ where
$\mathcal M_d$ is the set of monomials of degree $\le d$ (ordered, say,
lexicographically).  For an ideal $I\subseteq K[X]$ put $I_{\le d}=I\cap K[X]_{\le d}$.

\begin{lemma}\label{lem:ideal-filtration}
Let $F\subseteq K$ be a subfield and $I\subseteq K[X]$ an ideal.
\begin{enumerate}[nosep]
\item $I$ is defined over $F$ (\cref{def:locus}) iff every $I_{\le d}$ is defined over
  $F$ as a subspace of $K[X]_{\le d}$.
\item If $I_{\le d}$ is defined over $F$ and $d'\le d$ then $I_{\le d'}$ is defined over
  $F$; hence $\FOD(I_{\le d'})\subseteq\FOD(I_{\le d})$.
\item There is $d_0$ such that $I$ is generated by $I_{\le d_0}$; for $d\ge d_0$,
  $I$ is defined over $F$ iff $I_{\le d}$ is defined over $F$.
\end{enumerate}
\end{lemma}
\begin{proof}
(1) ``$\Rightarrow$'': $I=J\cdot K[X]$ with $J=I\cap F[X]$.  Choosing an $F$-basis of
$K$ containing $1$ (\cref{lem:extcontr}), $I=\bigoplus_\beta J\beta$ and the degree
filtration is compatible with this decomposition, so $I_{\le d}=\bigoplus_\beta J_{\le d}\beta$
is spanned by $J_{\le d}\subseteq I_{\le d}\cap F[X]$.  ``$\Leftarrow$'': $I=\bigcup_d I_{\le d}$
is generated by $\bigcup_d(I_{\le d}\cap F[X])\subseteq I\cap F[X]$.
(2) $I_{\le d'}=I_{\le d}\cap K[X]_{\le d'}$, and the intersection of two subspaces
defined over $F$ is defined over $F$: tensoring the exact sequence
$0\to U\cap U'\to U\oplus U'\to V$ of $F$-spaces with the flat $F$-module $K$ shows
$(U\otimes K)\cap(U'\otimes K)=(U\cap U')\otimes K$.  Then apply \cref{lem:fod-subspace}.
Here $\FOD(I_{\le d'})$ is computed in $K^{\mathcal M_{d'}}$ and $\FOD(I_{\le d})$ in
$K^{\mathcal M_d}$; since $\FOD$ of a subspace is its smallest field of definition
(\cref{lem:fod-subspace}), it does not depend on the ambient coordinate space in which
the subspace is placed (a subspace of $K^{\mathcal M_{d'}}$ is defined over $F$ iff its
image in $K^{\mathcal M_d}$ under the coordinate inclusion is), nor on the chosen
monomial order.
(3) $K[X]$ is noetherian; if $I$ is generated by polynomials of degree $\le d_0$ and
$I_{\le d}$ ($d\ge d_0$) is defined over $F$ then $I$ is generated by elements of
$I_{\le d}\cap F[X]$, so $I$ is defined over $F$; the converse is (1).
\end{proof}

\begin{definition}[smallest field of definition]\label{def:fod}
For an ideal $I\subseteq K[X]$, $\FOD(I):=\bigcup_{d}\FOD(I_{\le d})=\FOD(I_{\le d_0})$,
an increasing union which stabilises by \cref{lem:ideal-filtration}.
\end{definition}

\begin{theorem}[Weil]\label{thm:weil-fod}
Let $I\subseteq K[X]$ be an ideal.  A subfield $F\subseteq K$ is a field of definition
of $I$ iff $\FOD(I)\subseteq F$.  $\FOD(I)$ is finitely generated over the prime field.
For every automorphism $\sigma$ of $K$: $\FOD(\sigma I)=\sigma\FOD(I)$, and $\sigma I=I$
iff $\sigma$ fixes $\FOD(I)$ pointwise.
\end{theorem}
\begin{proof}
Combine \cref{lem:ideal-filtration} with \cref{lem:fod-subspace}, noting
$(\sigma I)_{\le d}=\sigma(I_{\le d})$.
\end{proof}

\begin{remark}
The same construction applies verbatim to ideals of $L[X]$ for any field $L$ with
subfields $F\subseteq L$, and for $J\subseteq L[X]$ one has $\FOD(J\cdot K[X])=\FOD(J)$
because $\operatorname{rref}((J\cdot K[X])_{\le d})=\operatorname{rref}(J_{\le d})$
(the RREF basis of $J_{\le d}$ over $L$ is in reduced row echelon form and spans
$(J K[X])_{\le d}$ by \cref{lem:ideal-filtration}(1)).  We shall only use ideals of
$K[X]$.
\end{remark}

\begin{leannote}
Mathlib has the predicate \lean{Matrix.IsReducedRowEchelon} but neither existence nor
uniqueness of the reduced row echelon basis; \cref{lem:rref} must be formalized (it is
elementary).  \lean{Submodule.restrictScalars}, \lean{Submodule.span} and flatness of
field extensions (\lean{Module.Flat}) supply the facts used in
\cref{lem:ideal-filtration}.  An alternative is to define $\FOD$ via Mathlib's
\lean{MonomialOrder} division algorithm and reduced Gr\"obner bases, but uniqueness of
reduced Gr\"obner bases is not in Mathlib either; the RREF route is shorter.
\end{leannote}

\subsection{Canonical bases}

\begin{definition}[canonical base]\label{def:cb}
Let $a$ be a tuple and $L$ a small algebraically closed subfield.  The
\emph{canonical base} of $a$ over $L$ is the subfield
\[
   \Cb(a/L):=k\cdot\FOD(P_L(a))\subseteq L,
\]
the field generated by $k$ and the smallest field of definition of the locus ideal
$P_L(a)$ (which lies in $L$ because $P_L(a)$ is defined over $L$).  It is a small
subfield.  For a finite tuple $w$ we write $\Cb(a/w)=\Cb(a/\acl(w))$.
\end{definition}

\begin{proposition}[properties of canonical bases]\label{prop:cb}
Let $a,a'$ be tuples, $L,L'$ small algebraically closed subfields and $C$ a small
algebraically closed subfield with $C\subseteq L$.
\begin{enumerate}[label=\textup{(Cb\arabic*)}]
\item\label{Cb:ind} $\rk(a/\Cb(a/L))=\rk(a/L)$, i.e.\ $a\ind{\Cb(a/L)}L$.
\item\label{Cb:char} $a\ind{C}L$ iff $\Cb(a/L)\subseteq C$.
\item\label{Cb:equiv} $\sigma(\Cb(a/L))=\Cb(\sigma a/\sigma L)$ for $\sigma\in\Gamma$; and for
  $\sigma\in\Gamma_L$, $\sigma$ fixes $\Cb(a/L)$ pointwise iff $P_L(\sigma a)=P_L(a)$ iff
  $\sigma a$ and $a$ are $\Gamma_L$-conjugate.
\item\label{Cb:two} If $a\ind{L}L'$ and $a\ind{L'}L$ then $\Cb(a/L)=\Cb(a/L')=\Cb(a/\acl(LL'))$.
\item\label{Cb:interalg} If $\acl(a)=\acl(a')$ then $\acl(\Cb(a/L))=\acl(\Cb(a'/L))$.
\item\label{Cb:basechange} If $e$ is a tuple with $e\ind{}aL$ then
  $\Cb(a/\acl(Le))=\Cb(a/L)$.
\end{enumerate}
\end{proposition}
\begin{proof}
\ref{Cb:ind}: Let $F=\Cb(a/L)$ and $F'=\acl(F)\subseteq L$, a small algebraically
closed field over which $P:=P_L(a)$ is defined.  Let $Q=P\cap F'[X]$, a prime of
$F'[X]$ with $Q\cdot K[X]=P$.  Since $a$ is a zero of $I_L(a)$ it is a zero of $P$ and
hence of $Q$, so $Q\subseteq I_{F'}(a)$.  By the ``consequently'' clause of \cref{lem:dim} ($\dim P$ may be computed with a
generic point over any small algebraically closed field over which $P$ is defined; a
generic point of $P$ over $F'$ has rank $\dim P=\dim P_L(a)=\rk(a/L)$, and its
vanishing ideal over $F'$ is $Q$ by \cref{lem:extcontr}) we get
$\trdeg_{F'}\operatorname{Frac}(F'[X]/Q)=\rk(a/L)$, while
$\trdeg_{F'}\operatorname{Frac}(F'[X]/I_{F'}(a))=\rk(a/F')$.  By \cref{lem:dimdrop}
(non-strict form, as $Q\subseteq I_{F'}(a)$), $\rk(a/F')\le\rk(a/L)$; the
reverse inequality is \ref{R:mono}.  Finally $\rk(a/F)=\rk(a/F')$ by \ref{R:acl}.

\ref{Cb:char}: If $a\ind{C}L$ then $P_C(a)=P_L(a)$ by \cref{lem:dim}(2), and $P_C(a)$
is defined over $C$, so $\Cb(a/L)=k\FOD(P_C(a))\subseteq C$.  Conversely, if
$\Cb(a/L)\subseteq C\subseteq L$ then by \ref{R:mono} and \ref{Cb:ind}
$\rk(a/L)\le\rk(a/C)\le\rk(a/\Cb(a/L))=\rk(a/L)$.

\ref{Cb:equiv}: $\sigma P_L(a)=P_{\sigma L}(\sigma a)$ and \cref{thm:weil-fod}.  For
$\sigma\in\Gamma_L$, $P_L(\sigma a)=\sigma P_L(a)$ equals $P_L(a)$ iff $\sigma$ fixes
$\FOD(P_L(a))$ pointwise; and $P_L(\sigma a)=P_L(a)$ iff $\sigma a$, $a$ are
$\Gamma_L$-conjugate by \cref{lem:dim}(3).

\ref{Cb:two}: Put $M=\acl(LL')$.  From $a\ind{L}L'$ we get $\rk(a/M)=\rk(a/LL')=\rk(a/L)$
by \ref{R:acl}, i.e.\ $a\ind{L}M$, so $P_L(a)=P_M(a)$ by \cref{lem:dim}(2) and
$\Cb(a/L)=\Cb(a/M)$.  Symmetrically $\Cb(a/L')=\Cb(a/M)$.

\ref{Cb:interalg}: Let $C=\acl(\Cb(a'/L))\subseteq L$.  By \ref{Cb:ind} and \ref{R:acl},
$a'\ind{C}L$; since $a\subseteq\acl(a')$, \ref{R:aclinv} gives $a\ind{C}L$, so
$\Cb(a/L)\subseteq C$ by \ref{Cb:char}, i.e.\ $\acl(\Cb(a/L))\subseteq\acl(\Cb(a'/L))$.
By symmetry the two closures coincide.

\ref{Cb:basechange}: $e\ind{}aL$ gives $e\ind{L}a$ and hence $a\ind{L}e$, i.e.
$a\ind{L}\acl(Le)$ by \ref{R:acl}; then $P_{\acl(Le)}(a)=P_L(a)$ by \cref{lem:dim}(2).
\end{proof}

\begin{remark}[what canonical bases are not]\label{rem:cb-not-acl}
In general $\Cb(a/L)\not\subseteq\acl(a)$: for independent $s,s',t$ and
$a=(t,st+s')$, $L=\acl(s,s')$, the locus of $a$ is the line $X_2=sX_1+s'$ and
$\Cb(a/L)=k(s,s')$, while $\rk(a)=2$ and $s\notin\acl(a)$.  The main theorem's
hypothesis ``$a\interalg\Cb(xy/a)$'' is therefore a genuine condition (see
\cref{sec:gc}).
\end{remark}

\begin{leannote}
$P_L(a)$ is \lean{(AclGeom.idealOf L a).map (MvPolynomial.map (algebraMap L K))}.  The
rank-two special case of \ref{Cb:char}--\ref{Cb:ind} (the coefficient closure of a plane
curve equation is the whole parameter plane) is \lean{AclGeom.QWitness.Psi.exists\_incidence\_germs\_with\_full\_coefficients}
in \lean{AclGeom/Config/ParameterField.lean}; there the field of definition of a principal
prime ideal is read off from the lexicographically monic generator
(\lean{AclGeom.FiniteCorrespondencePair.curveEquation}).  For principal ideals the two
descriptions agree: if $P=(f)$ with $f$ monic for a monomial order then
$\FOD(P)$ is the field generated by the coefficients of $f$ (the RREF basis of
$P_{\le d}$ consists of the monomial multiples of $f$ reduced against each other, whose
coefficients lie in the field generated by those of $f$, and conversely $f$ is the
unique element of $P$ with its leading monomial and coefficient $1$ reduced modulo
the others, so its coefficients lie in every field of definition).
\end{leannote}

\section{Definable closure, quasi-rational functions, and codes of finite sets}
\label{sec:codes}

\subsection{Definable closure}

\begin{lemma}[Galois characterisation of $\dcl$]\label{lem:dcl-galois}
Let $L\subseteq K$ be a subfield containing $k$ and $x\in K$.  Then $x\in\dcl(L)$ (the
perfect closure of $L$ in $K$) iff $\sigma x=x$ for every $\sigma\in\Gamma_L$.
\end{lemma}
\begin{proof}
If $x^{p^e}=c\in L$ then $x$ is the unique root in $K$ of $X^{p^e}-c$ (in
characteristic $0$, $e=0$ and $x=c$), and every $\sigma\in\Gamma_L$ permutes the roots of
this polynomial, so fixes $x$.  Conversely suppose $x\notin\dcl(L)$.  If $x$ is
transcendental over $L$, \cref{lem:orbit} (which only used \cref{lem:conj}, valid for
any subfield $L$) provides $\sigma\in\Gamma_L$ with $\sigma x\ne x$.  If $x$ is
algebraic over $L$, let $m\in L[X]$ be its minimal polynomial and $N=\deg m$.  If $x$
were the only root of $m$ in $K$ then $m=(X-x)^N$; write $N=p^eN'$ with $p\nmid N'$
(in characteristic $0$, $e=0$, $N'=N$), so $m=(X^{p^e}-x^{p^e})^{N'}$ and the
coefficient of $X^{p^e(N'-1)}$ is $-N'x^{p^e}\in L$, whence $x^{p^e}\in L$ and
$x\in\dcl(L)$, a contradiction.  So $m$ has a root $x'\ne x$; the vanishing ideals of $x$
and $x'$ over $L$ are both $(m)$, and \cref{lem:conj} gives $\sigma\in\Gamma_L$ with
$\sigma x=x'$.
\end{proof}

\begin{corollary}\label{cor:dcl-props}
For sets $A,B$: $k(A)\subseteq\dcl(A)\subseteq\acl(A)$; $\dcl(A)$ is a field;
$\dcl(\dcl(A))=\dcl(A)$; $A\subseteq B\Rightarrow\dcl(A)\subseteq\dcl(B)$; and for
$\sigma\in\Gamma$, $\sigma(\dcl A)=\dcl(\sigma A)$.  If $\dcl(CA)=\dcl(CA')$ then
$\acl(CA)=\acl(CA')$.
\end{corollary}
\begin{proof}
All clear from the definition (perfect closure); the last from
$\dcl\subseteq\acl$ and $\acl$ being a closure operator.
\end{proof}

\subsection{Quasi-rational functions}

\begin{definition}[quasi-rational functions]\label{def:qr}
Let $L$ be a subfield of $K$ containing $k$ and $U\subseteq K^n$.  A function
$f:U\to K^{n'}$ is \emph{quasi-rational over $L$} if there are $e\ge0$ (with $e=0$ in
characteristic $0$) and rational functions $R_1,\dots,R_{n'}\in L(X_1,\dots,X_n)$ whose
denominators do not vanish on $U$, such that $f(u)_j=R_j(u)^{1/p^e}$ for all $u\in U$
(the unique $p^e$-th root).  A \emph{quasi-rational expression} over $L$ is such a datum
$(e,R_1,\dots,R_{n'})$; it defines a function on the set where the denominators are
nonzero.
\end{definition}

\begin{lemma}\label{lem:qr-dcl}
Let $L\supseteq k$ be a subfield and $u\in K^n$, $w\in K^{n'}$.  Then $w\in\dcl(L,u)$
iff there is a quasi-rational expression $f$ over $L$ defined at $u$ with $f(u)=w$.
Moreover, if $f$ is a quasi-rational expression over $L$ defined at $u$, then for
every $\sigma\in\Gamma_L$ one has $f(\sigma u)=\sigma(f(u))$.  Compositions of
quasi-rational expressions are quasi-rational expressions.
\end{lemma}
\begin{proof}
$w_j\in\dcl(L,u)$ means $w_j^{p^e}\in L(u)$ for some $e$, i.e.\ $w_j^{p^e}=R_j(u)$
with $R_j\in L(X)$ whose denominator does not vanish at $u$; a common $e$ can be
chosen since $e$ may be increased by raising $R_j$ to a $p$-th power.  Equivariance:
$\sigma$ fixes the coefficients of $R_j$ and commutes with $p^e$-th roots.
Composition: $R_j\bigl(S_1(X)^{1/p^e},\dots\bigr)^{1/p^{e'}}
=\bigl(R_j^{(p^e)}(S_1(X),\dots)\bigr)^{1/p^{e+e'}}$, where $R_j^{(p^e)}$ is $R_j$ with all
coefficients raised to the power $p^e$ (Frobenius is a ring endomorphism of $K$).
\end{proof}

\begin{leannote}
In characteristic $0$ (the only case in which \lean{PerfectClosure} is trivial) all of
this collapses to $\dcl(A)=k(A)$ and rational functions.  In positive characteristic,
Mathlib provides \lean{perfectClosure} (as an intermediate field) and
\lean{IsPurelyInseparable}; \cref{lem:dcl-galois} is the statement that the fixed
field of \lean{Aut(K/L)} is the perfect closure, to be proved from
\lean{AclGeom.exists\_ambientAutomorphism\_of\_idealOf\_eq}.
\end{leannote}

\subsection{Codes of finite sets of tuples}

\begin{definition}[code of a finite set]\label{def:code}
Let $S=\{s^{1},\dots,s^{d}\}\subseteq K^m$ be a finite set of $d$ distinct tuples.
Put
\[
  P_S(T,U_1,\dots,U_m)=\prod_{i=1}^{d}\Bigl(T-\sum_{j=1}^{m}s^{i}_jU_j\Bigr)\in K[T,U],
\]
a polynomial of degree $d$, and let $\code{S}\in K^{\mathcal N_{d,m}}$ be its vector of
coefficients, indexed by the monomials $\mathcal N_{d,m}$ of degree $\le d$ in $T,U$.
\end{definition}

\begin{lemma}[properties of codes]\label{lem:code}
Let $S,S'\subseteq K^m$ be finite sets with $|S|=|S'|=d$ and let $\sigma\in\Gamma$.
\begin{enumerate}[nosep]
\item $\sigma\code{S}=\code{\sigma S}$, and $\code{S}\subseteq k(S)$.
\item $\code{S}=\code{S'}$ implies $S=S'$.  Hence $\sigma$ fixes $\code{S}$ pointwise
  iff $\sigma S=S$.
\item Every coordinate of every $s\in S$ lies in $\acl(\code{S})$; so
  $\acl(\code{S})=\acl(S)$.
\item If $L\supseteq k$ is a subfield and $\sigma S=S$ for all $\sigma\in\Gamma_L$, then
  $\code{S}\subseteq\dcl(L)$.
\end{enumerate}
\end{lemma}
\begin{proof}
(1) is clear.  (2): in the polynomial ring $K[U][T]$ over the domain $K[U]$, $P_S$ is
monic in $T$ and its roots in $K(U)$ are exactly the linear forms
$\ell_i=\sum_j s^i_jU_j$, each with multiplicity one since the $s^i$ are distinct; a
monic polynomial over a field determines its multiset of roots
(\lean{Polynomial.roots}), and $\ell_i$ determines $s^i$.  Hence $S$ is determined by
$P_S$.  The second sentence follows from (1).  (3): Let $F=k(\code{S})$, a small
subfield.  For $\sigma\in\Gamma_F$ we have $\sigma\code{S}=\code{S}$, hence $\sigma S=S$
by (2); so each coordinate $s^i_j$ has a finite $\Gamma_F$-orbit (contained in the set
of coordinates of elements of $S$) and lies in $\acl(F)$ by \cref{lem:orbit}.  The
inclusion $\code{S}\subseteq k(S)\subseteq\acl(S)$ gives the other direction.
(4): by (2), $\Gamma_L$ fixes $\code{S}$ pointwise; apply \cref{lem:dcl-galois}.
\end{proof}

\begin{definition}[conjugate sets]\label{def:conjset}
Let $L\supseteq k$ be a subfield and $z\in K^m$ a tuple with $z\subseteq\acl(L)$.  The
\emph{conjugate set} of $z$ over $L$ is the orbit $\Gamma_L\cdot z$, a finite set by
\cref{lem:orbit} (coordinatewise); its \emph{code} is $\code{\Gamma_L z}$.
\end{definition}

\begin{lemma}[codes of conjugate sets]\label{lem:conjcode}
In the situation of \cref{def:conjset}, with $c=\code{\Gamma_L z}$:
$c\subseteq\dcl(L)$, $\acl(c)=\acl(\Gamma_L z)$, and $z\subseteq\acl(c)$.
Moreover, if $L_0\subseteq L$ is a subfield containing $k$, $w$ is a tuple with
$w\in\dcl(L_0,z)$ and $w\subseteq L$, then $w\in\dcl(L_0,c)$.
\end{lemma}
\begin{proof}
The first three claims are \cref{lem:code}(4),(3) ($\Gamma_L$ permutes its own orbit).
For the last claim, write $w=f(z)$ with $f$ a quasi-rational expression over $L_0$
(\cref{lem:qr-dcl}), and let $L_1=k(L_0,c)$.  For $\sigma\in\Gamma_{L_1}$,
$\sigma$ fixes $c$ hence stabilises $\Gamma_Lz$ (\cref{lem:code}(2)), so
$\sigma z=\tau z$ for some $\tau\in\Gamma_L$.  Then, using equivariance of $f$ over
$L_0\subseteq L_1$ and over $L_0\subseteq L$,
\[
   \sigma w=\sigma f(z)=f(\sigma z)=f(\tau z)=\tau f(z)=\tau w=w,
\]
the last step because $w\subseteq L$ is fixed by $\tau\in\Gamma_L$.  So $w$ is fixed by
$\Gamma_{L_1}$ and $w\in\dcl(L_1)=\dcl(L_0,c)$ by \cref{lem:dcl-galois}.
\end{proof}

\begin{leannote}
Codes of finite sets of elements of $K$ (the case $m=1$) are the elementary symmetric
functions (\lean{Multiset.esymm}, \lean{Multiset.prod\_X\_sub\_C\_coeff}).  For tuples
the polynomial $P_S$ lives in \lean{MvPolynomial (Option (Fin m)) K} (or
\lean{Polynomial (MvPolynomial (Fin m) K)}); the uniqueness of roots argument uses
\lean{Polynomial.roots} over the fraction field of \lean{MvPolynomial (Fin m) K}.
\end{leannote}

\section{Germs of quasi-rational maps}
\label{sec:germs}

\subsection{Irreducible varieties and generic points}

\begin{definition}[irreducible $k$-variety]\label{def:variety}
An \emph{irreducible $k$-variety} $V$ in $n$ variables is a prime ideal
$I_V\subseteq K[Y_1,\dots,Y_n]$ defined over $k$ (\cref{def:locus}).  Its set of
$K$-points is $V(K)=\{y\in K^n: f(y)=0\ \forall f\in I_V\}$ and its dimension is
$\dim V=\dim I_V$ (\cref{lem:dim}).  For a small algebraically closed $M$, a tuple
$y$ is a \emph{generic point of $V$ over $M$} if $P_M(y)=I_V$.  Every tuple $y_0$
determines the irreducible $k$-variety $\Loc(y_0)$ with ideal $P_k(y_0)$, of which
$y_0$ is a generic point over $k$.
\end{definition}

\begin{lemma}\label{lem:generic-char}
Let $V$ be an irreducible $k$-variety and $M\subseteq M'$ small algebraically closed
subfields.
\begin{enumerate}[nosep]
\item Every $y\in V(K)$ has $\rk(y/M)\le\dim V$, and $y\in V(K)$ is generic over $M$
  iff $\rk(y/M)=\dim V$.  A generic point over $M'$ is generic over $M$, and a generic
  point $y$ over $M$ is generic over $M'$ iff $y\ind{M}M'$.
\item Generic points of $V$ over $M$ exist, and they form a single $\Gamma_M$-orbit.
\item If $y$ is generic over $M$ and $z\in\acl(M,y)$ then $\rk(yz/M)=\dim V$.
\end{enumerate}
\end{lemma}
\begin{proof}
(1) Let $y\in V(K)$.  Then $I_V\cap M[Y]\subseteq I_M(y)$, both prime ideals of $M[Y]$,
and $\trdeg_M\operatorname{Frac}(M[Y]/(I_V\cap M[Y]))=\dim V$ (compute $\dim I_V$ with a
generic point over $M$, \cref{lem:dim}; its vanishing ideal over $M$ is $I_V\cap M[Y]$ by
\cref{lem:extcontr}) while $\trdeg_M\operatorname{Frac}(M[Y]/I_M(y))=\rk(y/M)$.  The
non-strict form of \cref{lem:dimdrop} gives $\rk(y/M)\le\dim V$, and if equality holds
the strict form forces $I_V\cap M[Y]=I_M(y)$, i.e.\ $I_V=P_M(y)$ (\cref{lem:extcontr}).
Conversely if $P_M(y)=I_V$ then $\rk(y/M)=\dim P_M(y)=\dim V$ by \cref{lem:dim}.
``Generic over $M'$ implies generic over $M$'' is the last clause of \cref{lem:genexists}
($I_V$ is defined over $k\subseteq M$).  Finally, for $y$ generic over $M$: $y$ is
generic over $M'$ iff $\rk(y/M')=\dim V=\rk(y/M)$ iff $y\ind{M}M'$.
(2) is \cref{lem:genexists}.  (3) is \ref{R:add}.
\end{proof}

\begin{lemma}[agreement of quasi-rational expressions]\label{lem:qr-agree}
Let $N$ be a small algebraically closed subfield, $V$ an irreducible $k$-variety, and
$f,f'$ quasi-rational expressions over $N$ in the variables $Y$ with values in $K^{n'}$.
If $f(y')=f'(y')$ for one generic point $y'$ of $V$ over $N$, then $f(y)=f'(y)$ for every
$y\in V(K)$ at which both $f$ and $f'$ are defined.
\end{lemma}
\begin{proof}
Write the $j$-th coordinates as $R_j^{1/p^e}$ and $R'^{1/p^{e'}}_j$ with
$R_j=P_j/Q_j$, $R'_j=P'_j/Q'_j$, $P_j,Q_j,P'_j,Q'_j\in N[Y]$.  The polynomial
$A_j=P_j^{p^{e'}}Q'^{\,p^e}_j-P'^{\,p^e}_jQ_j^{p^{e'}}\in N[Y]$ vanishes at $y'$, hence
lies in $I_N(y')=I_V\cap N[Y]$, hence vanishes at every $y\in V(K)$.  If $Q_j(y)\ne0$
and $Q'_j(y)\ne0$ this gives $R_j(y)^{p^{e'}}=R'_j(y)^{p^e}$, and taking the unique
$p^{e+e'}$-th root in $K$, $R_j(y)^{1/p^e}=R'_j(y)^{1/p^{e'}}$.
\end{proof}

\subsection{Germs}

Throughout this section $V\subseteq K^n$ and $V'\subseteq K^{n'}$ are irreducible
$k$-varieties.

\begin{definition}[representatives and germs]\label{def:germ}
A \emph{representative} of a germ from $V$ to $V'$ is a pair $(M,f)$ where $M$ is a
small algebraically closed subfield and $f$ is a quasi-rational expression over $M$ in
the variables $Y$ with values in $K^{n'}$ such that for every generic point $y$ of $V$
over $M$, $f$ is defined at $y$ and $f(y)$ is a generic point of $V'$ over $M$.
Two representatives $(M,f)$ and $(M',f')$ are \emph{equivalent} if $f(y)=f'(y)$ for
every generic point $y$ of $V$ over $\acl(MM')$.  A \emph{germ} $g:V\to V'$ is an
equivalence class of representatives.  The set of germs is $\Germ(V,V')$ and
$\Germ(V)=\Germ(V,V)$.
\end{definition}

The relation is an equivalence relation.  Reflexivity and symmetry are clear.  For
transitivity let $(M,f)\sim(M',f')\sim(M'',f'')$ and let $y$ be a generic point of $V$
over $N=\acl(MM'')$.  Choose $y'$ generic over $\acl(MM'M'')$ (\cref{lem:genexists});
it is generic over $\acl(MM')$, $\acl(M'M'')$ and $N$ (\cref{lem:generic-char}(1)), so
$f(y')=f'(y')=f''(y')$, and then $f(y)=f''(y)$ by \cref{lem:qr-agree} (both are defined at
$y$ since $y$ is generic over $M$ and over $M''$).  Note also that if $(M,f)$ is a
representative and $M\subseteq M_1$ is small algebraically closed, then $(M_1,f)$ is a
representative equivalent to $(M,f)$: for $y$ generic over $M_1$, $y$ is generic over
$M$, so $f$ is defined at $y$ and $f(y)\in V'(K)$; moreover $y\ind{M}M_1$
(\cref{lem:generic-char}(1)) and $f(y)\subseteq\acl(M,y)$ give $f(y)\ind{M}M_1$
(\ref{R:aclinv}, \ref{R:sub}), so $\rk(f(y)/M_1)=\rk(f(y)/M)=\dim V'$ and $f(y)$ is
generic over $M_1$; the two representatives agree at every generic point over
$\acl(MM_1)=M_1$.

\begin{definition}[graph ideal, code, rank]\label{def:graph}
For a representative $(M,f)$ let $y$ be a generic point of $V$ over $M$ and put
$\Gamma(M,f)=P_M\bigl((y,f(y))\bigr)\subseteq K[Y,Z]$.
\end{definition}

\begin{lemma}[the graph ideal]\label{lem:graph}
\begin{enumerate}[nosep]
\item $\Gamma(M,f)$ does not depend on the choice of the generic point $y$, and
  equivalent representatives have the same graph ideal.  We write $\Gamma(g)$.
\item $\Gamma(g)\cap K[Y]=I_V$, $\Gamma(g)\cap K[Z]=I_{V'}$, and $\dim\Gamma(g)=\dim V$.
\item $\Gamma(g)$ determines $g$: if $\Gamma(g)=\Gamma(g')$ then $g=g'$.
\end{enumerate}
\end{lemma}
\begin{proof}
(1) Generic points of $V$ over $M$ are $\Gamma_M$-conjugate (\cref{lem:generic-char}),
and $\sigma(y,f(y))=(\sigma y,f(\sigma y))$ for $\sigma\in\Gamma_M$ by \cref{lem:qr-dcl},
so the locus ideals agree by \cref{lem:dim}(3).  If $(M,f)\sim(M',f')$, let $y$ be
generic over $N=\acl(MM')$; then $y$ is generic over $M$ and over $M'$,
$f(y)=f'(y)$, and $P_N((y,f(y)))=P_M((y,f(y)))$ because $(y,f(y))\ind{M}N$:
indeed $\rk(yf(y)/N)=\rk(y/N)=\dim V=\rk(y/M)=\rk(yf(y)/M)$ using
$f(y)\in\dcl(M,y)\subseteq\acl(M,y)$.  Symmetrically for $M'$.
(2) $I_M((y,f(y)))\cap M[Y]=I_M(y)=I_V\cap M[Y]$ and similarly for $Z$ since $f(y)$ is
generic in $V'$ over $M$; extend scalars and drop variables
(\cref{lem:extcontr}(2): $\Gamma(g)\cap K[Y]=(I_M((y,f(y)))\cap M[Y])K[Y]$).  The
dimension is $\rk(yf(y)/M)=\dim V$.
(3) Let $(M,f)$, $(M',f')$ represent $g,g'$ and let $y$ be generic over $N=\acl(MM')$.
As in (1), $(y,f(y))$ and $(y,f'(y))$ are generic points of $\Gamma(g)=\Gamma(g')$
over $N$, hence $\Gamma_N$-conjugate: there is $\sigma\in\Gamma_N$ with
$\sigma(y,f(y))=(y,f'(y))$.  Then $\sigma$ fixes $N(y)$ pointwise, hence
$\dcl(N,y)\ni f(y)$ pointwise (\cref{lem:dcl-galois}), so $f'(y)=\sigma f(y)=f(y)$.
\end{proof}

\begin{definition}[code, rank, closure of a germ]\label{def:germcode}
For a germ $g$ put $\Cb(g)=k\cdot\FOD(\Gamma(g))$, a small subfield;
$\rk(g)=\trdeg(\Cb(g)/k)$; $\acl(g)=\acl(\Cb(g))$ and, for a set $C$,
$\rk(g/C)=\rk(c/C)$ where $c$ is any finite tuple generating $\Cb(g)$ over $k$.
Independence $g\ind{C}B$ and interalgebraicity $g\interalg a$ with tuples are defined
through $\Cb(g)$ (e.g.\ $g\interalg a$ iff $\acl(\Cb(g))=\acl(a)$).  The
\emph{code tuple} $\code{g}$ is the finite tuple of all coordinates of the RREF bases
of $\Gamma(g)_{\le d}$ for $d\le d_0(\Gamma(g))$, where $d_0(I)$ is the least $d$ such
that $I$ is generated by $I_{\le d}$; so $k(\code g)=\Cb(g)$.
\end{definition}

\begin{lemma}[automorphisms act on germs]\label{lem:germ-aut}
For $\sigma\in\Gamma$ and a representative $(M,f)$ of $g$, $(\sigma M,\sigma f)$ is a
representative of a germ $\sigma g$ depending only on $g$; $\Gamma(\sigma g)=\sigma\Gamma(g)$,
$\Cb(\sigma g)=\sigma\Cb(g)$, $\code{\sigma g}=\sigma\code{g}$, and $\sigma g=g$ iff
$\sigma$ fixes $\Cb(g)$ pointwise.  If $(M,f)$ represents $g$ then $\Cb(g)\subseteq M$.
\end{lemma}
\begin{proof}
Equivariance of everything involved (\cref{lem:equiv}, \cref{lem:qr-dcl}, and
$\sigma(\Gamma(M,f))=\Gamma(\sigma M,\sigma f)$).  For the last two claims:
$\sigma g=g$ iff $\Gamma(\sigma g)=\Gamma(g)$ (\cref{lem:graph}(3)) iff
$\sigma\Gamma(g)=\Gamma(g)$ iff $\sigma$ fixes $\FOD(\Gamma(g))$ pointwise
(\cref{thm:weil-fod}), and $\sigma$ fixes $k$; and $\Gamma(g)=P_M(\cdot)$ is defined
over $M$, so $\FOD(\Gamma(g))\subseteq M$.
\end{proof}

\begin{convention}[germs in tuple positions]\label{conv:germ-tuple}
In ranks, closures, independence statements and tuples, a germ $g$ stands for its code
tuple $\code g$ (equivalently for the field $\Cb(g)=k(\code g)$): thus $\acl(g,w)$ means
$\acl(\code g,w)$, $h\ind{}gw$ means $\code h\ind{}\code g\,w$, and so on.
\end{convention}

\begin{lemma}[fibres and evaluation]\label{lem:eval}
Let $g\in\Germ(V,V')$ and let $F$ be a small algebraically closed subfield with
$\Cb(g)\subseteq F$.
\begin{enumerate}[nosep]
\item If $(y,z)$ and $(y,z')$ are generic points of $\Gamma(g)$ over $F$ then $z=z'$.
\item For every generic point $y$ of $V$ over $F$ there is a unique $z$, denoted
  $g(y)$, such that $(y,z)$ is a generic point of $\Gamma(g)$ over $F$; moreover
  $g(y)\in\dcl(F,y)$, $g(y)$ is a generic point of $V'$ over $F$, and $g(y)$ does not
  depend on $F$ (if $F\subseteq F'$ and $y$ is generic over $F'$, the two values agree).
\item $g$ has a representative over $F$; in particular over $\acl(\Cb(g))$.
\item (equivariance) $\sigma(g(y))=(\sigma g)(\sigma y)$ for $\sigma\in\Gamma$.
\item (agreement) If $g'\in\Germ(V,V')$ and $g(y)=g'(y)$ for one generic point $y$ of
  $V$ over $\acl(\Cb(g)\Cb(g'))$, then $g=g'$.
\end{enumerate}
\end{lemma}
\begin{proof}
(1) Choose a representative $(M,f)$ of $g$; replacing $M$ by $\acl(MF)$ (allowed by
the note after \cref{def:germ}) we may assume $M\supseteq F$, and write $M=\acl(F,m)$
for a finite tuple $m$.  By \cref{lem:fresh-conj} there is $\tau\in\Gamma_F$ with
$\tau m\ind{F}yzz'$; replacing $(M,f)$ by the equivalent representative
$(\tau M,\tau f)$ of $g$ (\cref{lem:germ-aut}, $\tau g=g$ as $\tau$ fixes
$F\supseteq\Cb(g)$) we may assume $M\ind{F}yzz'$.  Then $(y,z)$ and $(y,z')$ are
generic points of $\Gamma(g)$ over $M$ (\cref{lem:dim}(2)) and so is $(y,f(y))$:
indeed $y$ is generic in $V$ over $F$ (by \cref{lem:extcontr}(2), $\Gamma(g)\cap K[Y]=I_V$
and $I_F(y)=I_F((y,z))\cap F[Y]$, so $P_F(y)=I_V$) and $y\ind{F}M$ (\ref{R:sub}), so $y$ is
generic over $M$ and $(y,f(y))$ is generic for $\Gamma(g)=\Gamma(M,f)$ over $M$
(\cref{lem:graph}(1)).  All three are $\Gamma_M$-conjugate.  If $\sigma\in\Gamma_M$ carries
$(y,f(y))$ to $(y,z)$ then $\sigma$ fixes $M(y)$ pointwise, hence $\dcl(M,y)\ni f(y)$,
so $z=\sigma f(y)=f(y)$; likewise $z'=f(y)$.
(2) Existence: with $(M,f)$ as in (1) chosen with $M\ind{F}y$, $y$ is generic over
$M$ and $z=f(y)$ works: $(y,z)$ is a generic point of $\Gamma(g)$ over $M$
(\cref{lem:graph}(1)), and $\Gamma(g)$ is defined over $F$, so $(y,z)$ is generic over
$F$ by the last clause of \cref{lem:genexists}.  Uniqueness is (1).  Galois: for
$\sigma\in\Gamma_{F(y)}$, $(y,\sigma z)$ is again generic over $F$ (it is the image of a
generic point under an element of $\Gamma_F$), so $\sigma z=z$ and $z\in\dcl(F,y)$ by
\cref{lem:dcl-galois}.  Genericity of $z$ in $V'$: $I_F(z)=I_F((y,z))\cap F[Z]=
\Gamma(g)\cap F[Z]$ and $(\Gamma(g)\cap F[Z])K[Z]=\Gamma(g)\cap K[Z]=I_{V'}$ by
\cref{lem:extcontr}(2) and \cref{lem:graph}(2).  Independence of $F$: if $y$ is generic
over $F'\supseteq F$ then the value $z$ computed over $F$ satisfies $\rk(yz/F')=\rk(y/F')
=\dim V=\dim\Gamma(g)$ and $(y,z)$ is a zero of $\Gamma(g)$, so $(y,z)$ is generic over
$F'$ by \cref{lem:generic-char}(1) applied to the variety $\Gamma(g)$ (which is defined
over $F\subseteq F'$; \cref{lem:generic-char} holds verbatim for varieties defined over
$F$ in place of $k$).
(3) Pick $y$ generic over $F$ and $z=g(y)\in\dcl(F,y)$; by \cref{lem:qr-dcl} there is a
quasi-rational expression $f$ over $F$ with $f(y)=z$.  For any other generic point
$y'=\sigma y$ ($\sigma\in\Gamma_F$), $f$ is defined at $y'$ and $f(y')=\sigma z$ is generic
in $V'$ over $F$ with $(y',f(y'))$ generic in $\Gamma(g)$; so $(F,f)$ is a representative,
and it represents $g$ because its graph ideal is $\Gamma(g)$ (\cref{lem:graph}(3)).
(4) Apply $\sigma$ to ``$(y,g(y))$ is generic in $\Gamma(g)$ over $F$''.
(5) $(y,g(y))=(y,g'(y))$ is a generic point of both $\Gamma(g)$ and $\Gamma(g')$ over
$\acl(\Cb(g)\Cb(g'))$, so the two ideals coincide and \cref{lem:graph}(3) applies.
\end{proof}

\begin{definition}[generic points over $k$; domain of a germ]\label{def:dom}
$V_{\Gen}$ denotes the set of generic points of $V$ over $k$.
$\operatorname{Dom}(g)$ is the set of generic points of $V$ over $\acl(\Cb(g))$,
i.e.\ (by \cref{lem:generic-char}(1)) the points $y\in V_{\Gen}$ with
$y\ind{}\Cb(g)$.  By \cref{lem:eval}(2), $g(y)$ is defined for $y\in\operatorname{Dom}(g)$.
\end{definition}

\subsection{Composition and inverses}

\begin{lemma}[composition]\label{lem:comp}
Let $g\in\Germ(V,V')$ and $g'\in\Germ(V',V'')$.  There is a unique germ
$g'\circ g\in\Germ(V,V'')$ such that $(g'\circ g)(y)=g'(g(y))$ for every generic point
$y$ of $V$ over $\acl(\Cb(g)\Cb(g'))$.  Composition is associative, the germ
$\mathrm{id}_V$ of $(k,\mathrm{id})$ is a two-sided identity, and
$\Cb(g'\circ g)\subseteq\acl(\Cb(g)\Cb(g'))$.  For $\sigma\in\Gamma$,
$\sigma(g'\circ g)=\sigma g'\circ\sigma g$.
\end{lemma}
\begin{proof}
Let $F=\acl(\Cb(g)\Cb(g'))$ and take representatives $(F,f)$, $(F,f')$
(\cref{lem:eval}(3)).  For $y$ generic over $F$, $f(y)$ is generic in $V'$ over $F$,
so $f'$ is defined at $f(y)$ and $f'(f(y))$ is generic in $V''$ over $F$; thus
$(F,f'\circ f)$ is a representative (quasi-rational by \cref{lem:qr-dcl}) and
$\Cb$ of its germ lies in $F$ (\cref{lem:germ-aut}).  Uniqueness by \cref{lem:eval}(5);
associativity and identities by evaluation at a generic point over the closure of all
codes involved; equivariance by \cref{lem:eval}(4).
\end{proof}

\begin{lemma}[inverses]\label{lem:inverse}
For $g\in\Germ(V,V')$ the following are equivalent:
\begin{enumerate}[nosep]
\item there is $h\in\Germ(V',V)$ with $h\circ g=\mathrm{id}_V$ and $g\circ h=\mathrm{id}_{V'}$;
\item for some (equivalently every) small algebraically closed $F\supseteq\Cb(g)$ and
  some (equivalently every) generic point $y$ of $V$ over $F$: $y\in\dcl(F,g(y))$.
\end{enumerate}
In that case $h=:g^{-1}$ is unique, $\Gamma(g^{-1})$ is the image of $\Gamma(g)$ under
the swap of the variables $Y$ and $Z$, $\Cb(g^{-1})=\Cb(g)$, $\dim V=\dim V'$, and
$\sigma(g^{-1})=(\sigma g)^{-1}$.  Such germs are called \emph{invertible}; we write
$\Germ^\times(V,V')$ for the set of invertible germs $V\to V'$, and the invertible germs
in $\Germ(V)$ form a group $\Germ^\times(V)=\Germ^\times(V,V)$ under composition.
\end{lemma}
\begin{proof}
(1)$\Rightarrow$(2): a two-sided inverse $h$ is unique (if $h'g=\mathrm{id}$ then
$h'=h'gh=h$), so for $\sigma\in\Gamma$ with $\sigma g=g$ we get $\sigma h=(\sigma g)^{-1}=h$
(\cref{lem:comp}); hence every $\sigma$ fixing $\Cb(g)$ pointwise fixes $\Cb(h)$ pointwise
(\cref{lem:germ-aut}), and \cref{lem:dcl-galois} gives $\Cb(h)\subseteq\dcl(\Cb(g))\subseteq\acl(\Cb(g))$.
Now let $F\supseteq\Cb(g)$ be small algebraically closed and $y$ generic over $F$.  Then
$F\supseteq\Cb(h)$, $z=g(y)$ is generic in $V'$ over $F$ (\cref{lem:eval}(2)), so $h(z)$ is
defined and equals $y$ (evaluate $h\circ g=\mathrm{id}$ at $y$, \cref{lem:comp}), and
$y=h(z)\in\dcl(F,z)$ by \cref{lem:eval}(2) applied to $h$.
(2)$\Rightarrow$(1): let $\Gamma^{t}$ be the swap of $\Gamma(g)$; $(z,y)$ with $z=g(y)$ is a
generic point of $\Gamma^t$ over $F$, $z$ is generic in $V'$ over $F$, and
$y\in\dcl(F,z)$; by \cref{lem:qr-dcl} there is a quasi-rational $f'$ over $F$ with
$f'(z)=y$, and as in the proof of \cref{lem:eval}(3), $(F,f')$ is a representative of a
germ $h:V'\to V$ with $\Gamma(h)=\Gamma^t$.  Evaluating at generic points gives
$h\circ g=\mathrm{id}$ and $g\circ h=\mathrm{id}$.  The ``every'' clauses: for another
$F'$, the condition $y\in\dcl(F',g(y))$ for $y$ generic over $F'$ follows from (1).
Uniqueness of $h$ is \cref{lem:eval}(5).  $\FOD$ is invariant under renaming variables
(\cref{thm:weil-fod}: fields of definition are), so $\Cb(g^{-1})=\Cb(g)$;
$\dim V'=\rk(z/F)=\rk(yz/F)=\dim V$ since $y\in\acl(F,z)$.
\end{proof}

\begin{lemma}[composition and codes]\label{lem:comp-codes}
Let $g\in\Germ^\times(V,V')$ and $h\in\Germ^\times(V',V'')$.
\begin{enumerate}[nosep]
\item $\Cb(h\circ g)\subseteq\acl(\Cb g,\Cb h)$, $\Cb(h)\subseteq\acl(\Cb(hg),\Cb(g))$ and
  $\Cb(g)\subseteq\acl(\Cb(hg),\Cb(h))$.  In words: any two of $g,h,hg$ are
  interalgebraic with the third over each other, so $\rk(hg/Cg)=\rk(h/Cg)$ and
  $\rk(hg/Ch)=\rk(g/Ch)$ for every set $C$.
\item For $y\in\operatorname{Dom}(g)\cap\operatorname{Dom}(hg)$ with $g(y)\in\operatorname{Dom}(h)$:
  $(hg)(y)=h(g(y))$.  (If $y$ is generic over $\acl(\Cb g\,\Cb h)$ this is the defining
  property of $hg$ in \cref{lem:comp}.)
\end{enumerate}
\end{lemma}
\begin{proof}
(1) The first is \cref{lem:comp}; $h=(hg)g^{-1}$ and $g=h^{-1}(hg)$ with
\cref{lem:inverse} ($\Cb(g^{-1})=\Cb(g)$) and \cref{lem:comp} again.
(2) Let $N=\acl(\Cb g\,\Cb h)$ and take representatives $(F_1,f_1)$ of $g$ over
$F_1=\acl(\Cb g)$, $(F_3,f_3)$ of $h$ over $F_3=\acl(\Cb h)$ and $(F_2,f_2)$ of $hg$ over
$F_2=\acl(\Cb(hg))$ (\cref{lem:eval}(3)).  By \cref{lem:eval}(2) (with the respective
fields), $g(y)=f_1(y)$, $h(g(y))=f_3(f_1(y))$ and $(hg)(y)=f_2(y)$, all defined.  The
quasi-rational expressions $f_2$ and $f_3\circ f_1$ over $N$ agree at a generic point
of $V$ over $N$ (\cref{lem:comp}), hence at $y$ by \cref{lem:qr-agree}.
\end{proof}

\begin{lemma}[values are definable over the codes]\label{lem:eval-dcl}
Let $g\in\Germ(V,V')$ and $y\in\operatorname{Dom}(g)$.  Then $g(y)\in\dcl(\code g,y)$,
where $\code g$ is the code tuple of \cref{def:germcode}.  More generally, if
$g_1,\dots,g_s$ are invertible germs between irreducible $k$-varieties, $w$ is a
composite of the $g_i$ and their inverses which is a germ $V\to V'$
(\cref{lem:comp,lem:inverse}), and $y\in V_{\Gen}$ is independent from
$\code{g_1}\cdots\code{g_s}$, then $y\in\operatorname{Dom}(w)$ and
$w(y)\in\dcl(\code{g_1},\dots,\code{g_s},y)$.
\end{lemma}
\begin{proof}
Let $L=k(\code g,y)$ and $\sigma\in\Gamma_L$.  Since $\sigma$ fixes $k(\code g)=\Cb(g)$
pointwise, $\sigma g=g$ (\cref{lem:germ-aut}); as $\sigma y=y$, \cref{lem:eval}(4) gives
$\sigma(g(y))=g(y)$.  By \cref{lem:dcl-galois}, $g(y)\in\dcl(L)$.  For a composite $w$:
iterating $\Cb(g'\circ g)\subseteq\acl(\Cb g\,\Cb g')$ (\cref{lem:comp}) and
$\Cb(g^{-1})=\Cb(g)$ (\cref{lem:inverse}) gives
$\Cb(w)\subseteq\acl(\code{g_1}\cdots\code{g_s})$, so $y\in\operatorname{Dom}(w)$; and every
$\sigma$ fixing $k(\code{g_1},\dots,\code{g_s},y)$ pointwise fixes each $g_i$
(\cref{lem:germ-aut}), hence $w$ (equivariance clauses of \cref{lem:comp,lem:inverse}:
$\sigma(g'\circ g)=\sigma g'\circ\sigma g$, $\sigma(g^{-1})=(\sigma g)^{-1}$), hence
$w(y)$ (\cref{lem:eval}(4)); conclude by \cref{lem:dcl-galois}.
\end{proof}

\subsection{The dictionary between triples and germs}

\begin{lemma}[germs from triples]\label{lem:triple}
Let $b$ be a tuple, $V=\Loc(y_0)$, $V'=\Loc(z_0)$ irreducible $k$-varieties, and let
$y\in V(K)$, $z\in V'(K)$ with
\[
   y\ind{}b,\qquad z\ind{}b,\qquad z\in\dcl(b,y),
\]
$y$ generic in $V$ and $z$ generic in $V'$ over $k$.  Then there is a unique germ
$g_b\in\Germ(V,V')$ with $\Cb(g_b)\subseteq\acl(b)$ and $g_b(y)=z$.  Its graph ideal is
$P_{\acl(b)}((y,z))$, so $\Cb(g_b)=\Cb(yz/b)$; moreover $\Cb(g_b)\subseteq\dcl(b)$.  It is
invertible iff $y\in\dcl(b,z)$.  Moreover:
\begin{enumerate}[nosep]
\item if $z=f(b,y)$ for a quasi-rational expression $f$ over $k$ in the variables $(B,Y)$,
  then for every $\sigma\in\Gamma$, $g_{\sigma b}:=\sigma g_b$ is the germ determined in
  the same way by $(\sigma b,\sigma y,\sigma z)$, and it depends only on $\sigma b$: it is
  the germ of the representative $(\acl(\sigma b),f(\sigma b,\cdot))$;
\item (generic transitivity) if moreover $y\ind{}z$ then for every $b'$ with
  $b'\ind{}y$ in the $\Gamma$-orbit of $b$, $g_{b'}(y)\ind{}y$;
\item $b\interalg\Cb(g_b)$ iff $b\subseteq\acl(\Cb(yz/b))$ (``canonicity'').
\end{enumerate}
\end{lemma}
\begin{proof}
$y$ is generic in $V$ over $M=\acl(b)$ (\cref{lem:generic-char}(1)) and $z$ in $V'$.
By \cref{lem:qr-dcl} there is a quasi-rational $f_b$ over $M$ with $f_b(y)=z$; for
$y'=\sigma y$, $\sigma\in\Gamma_M$, $f_b(y')=\sigma z$ is generic in $V'$ over $M$, so
$(M,f_b)$ is a representative; let $g_b$ be its germ.  Then $\Gamma(g_b)=P_M((y,z))$ and
$\Cb(g_b)=k\FOD(P_{\acl(b)}((y,z)))=\Cb(yz/b)\subseteq M$; uniqueness by
\cref{lem:eval}(5).

$\Cb(g_b)\subseteq\dcl(b)$: write $z=f(b,y)$ for a quasi-rational $f$ over $k$ in the
variables $(B,Y)$ (\cref{lem:qr-dcl}, as $z\in\dcl(b,y)$).  Then $(M,f(b,\cdot))$ is a
representative of $g_b$: $f(b,\cdot)$ is an expression over $k(b)\subseteq M$ which agrees
with $f_b$ at $y$, hence (\cref{lem:qr-dcl}, equivariance under $\Gamma_M$) at every
$\Gamma_M$-conjugate of $y$, i.e.\ at every generic point of $V$ over $M$
(\cref{lem:generic-char}(2)).  For $\sigma\in\Gamma_{k(b)}$ this representative is carried
to $(\sigma M,f(b,\cdot))=(M,f(b,\cdot))$ itself, so $\sigma g_b=g_b$; hence $\sigma$ fixes
$\Cb(g_b)$ pointwise (\cref{lem:germ-aut}) and $\Cb(g_b)\subseteq\dcl(b)$ by
\cref{lem:dcl-galois}.

Invertibility: if $y\in\dcl(b,z)\subseteq\dcl(M,z)$ then $g_b$ is invertible by
\cref{lem:inverse}(2) with $F=M$.  Conversely, if $g_b$ is invertible then
$\Cb(g_b^{-1})=\Cb(g_b)\subseteq\dcl(b)$ (\cref{lem:inverse}), $z$ is generic in $V'$ over
$M\supseteq\acl(\Cb g_b^{-1})$, so $z\in\operatorname{Dom}(g_b^{-1})$ and
$y=g_b^{-1}(z)$ (evaluate $g_b^{-1}\circ g_b=\mathrm{id}$ at $y$, \cref{lem:comp}); by
\cref{lem:eval-dcl}, $y\in\dcl(\code{g_b^{-1}},z)\subseteq\dcl(\dcl(b),z)=\dcl(b,z)$
(\cref{cor:dcl-props}).
(1) The triple $(\sigma b,\sigma y,\sigma z)$ satisfies the same hypotheses and
$\sigma g_b$ has $\Gamma(\sigma g_b)=P_{\acl(\sigma b)}(\sigma(y,z))$.  If $z=f(b,y)$ with
$f$ over $k$, then $(\acl(b),f(b,\cdot))$ is a representative of $g_b$ (shown above), and
$\sigma$ carries it to $(\acl(\sigma b),f(\sigma b,\cdot))$, which depends only on
$\sigma b$.
(2) Let $b'=\tau b$ with $\tau\in\Gamma$.  Both $\tau y$ and $y$ are generic points of
$V$ over $\acl(b')$ ($\tau y$ because $y$ is generic over $\acl(b)$ and $\tau$ is an
automorphism; $y$ because $y\ind{}b'$ and \cref{lem:generic-char}(1)), so there is
$\rho\in\Gamma_{\acl(b')}$ with $\rho\tau y=y$ (\cref{lem:generic-char}(2)).  Put
$\sigma=\rho\tau$; then $\sigma b=b'$, $\sigma y=y$, and by (1) $g_{b'}=\sigma g_b$, so
$g_{b'}(y)=\sigma(g_b(y))=\sigma z$ (\cref{lem:eval}(4)) and
$y=\sigma y\ind{}\sigma z$ by \cref{lem:equiv}.
(3) $\Cb(g_b)\subseteq\acl(b)$ always; so $\acl(\Cb g_b)=\acl(b)$ iff $b\subseteq\acl(\Cb g_b)$.
\end{proof}

\begin{remark}[canonical parametrisation by codes]\label{rem:code-param}
Let $g$ be a germ and $c=\code{g}$ (\cref{def:germcode}).  For $\sigma\in\Gamma$ the germ
$\sigma g$ has code $\sigma c$ and is the germ with that code; thus the map
$\sigma c\mapsto\sigma g$ is a well-defined bijection between the $\Gamma$-orbit of $c$
(the generic points over $k$ of the irreducible $k$-variety $\Loc(c)$) and the
$\Gamma$-orbit of $g$.  Moreover, for $y$ generic in $V$ over $\acl(c)$ one has
$g(y)\in\dcl(c,y)$ (\cref{lem:eval-dcl}), so $g(y)=f(c,y)$ for a
quasi-rational $f$ over $k$, and then $(\sigma g)(y')=f(\sigma c,y')$ for all generic
$y'$ over $\acl(\sigma c)$.  Every $\Gamma$-orbit of germs is therefore a ``definable
family parametrised by its codes'' in the sense of \cref{lem:triple}(1), with
$b=c$ and canonicity automatic.
\end{remark}

\begin{leannote}
The repository's \lean{AclGeom.FiniteCorrespondencePair} is a pair $(y,z)$ of
interalgebraic generic elements; germs in the present sense are its single-valued
analogue (no branch choice), and \lean{AclGeom.FiniteCorrespondenceGerm.Composes}
corresponds to \cref{lem:comp} in the special case of a literally shared middle
point.  A Lean implementation should define \lean{Germ V V'} as a quotient of the type
of representatives by the relation of \cref{def:germ}, or equivalently as the subtype of
prime ideals of \lean{MvPolynomial (Fin n ⊕ Fin n') K} characterised by
\cref{lem:graph}(2) together with the single-valuedness of \cref{lem:eval}(1); the two
presentations are identified by \cref{lem:graph}(3) and \cref{lem:eval}(3).
\end{leannote}

\section{The Hrushovski--Weil theorem}
\label{sec:hw}

\subsection{Generic sets of germs}

\begin{lemma}[independent generic tuples are conjugate]\label{lem:indep-conj}
Let $V_1,\dots,V_r$ be irreducible $k$-varieties and let $(y_1,\dots,y_r)$,
$(y'_1,\dots,y'_r)$ be two independent tuples ($\rk(y_1\cdots y_r)=\sum_i\dim V_i$,
likewise for the primed tuple) with $y_i,y'_i$ generic points of $V_i$ over $k$.  Then
there is $\sigma\in\Gamma$ with $\sigma y_i=y'_i$ for all $i$.  The same holds with
$\Gamma$ replaced by $\Gamma_M$ for a small algebraically closed $M$ when all the tuples
are generic over $M$ and independent over $M$ (same proof with $k$ replaced by $M$).
\end{lemma}
\begin{proof}
Induction on $r$.  For $r=1$ this is \cref{lem:generic-char}(2).  Given $\sigma_1\in\Gamma$
with $\sigma_1(y_1,\dots,y_{r-1})=(y'_1,\dots,y'_{r-1})$, the tuples $\sigma_1y_r$ and
$y'_r$ are both generic points of $V_r$ over $\acl(y'_1\cdots y'_{r-1})$ (independence
and \cref{lem:generic-char}(1)), hence conjugate under
$\Gamma_{\acl(y'_1\cdots y'_{r-1})}$ by $\tau$; take $\sigma=\tau\sigma_1$.
\end{proof}

Throughout this section $V$ is an irreducible $k$-variety of dimension $m$ and
$V_{\Gen}$ is the set of generic points of $V$ over $k$ (\cref{def:dom}); germs in
tuple positions stand for their codes (\cref{conv:germ-tuple}).

\begin{definition}[generic set of germs]\label{def:genset}
A \emph{generic set of germs} on $V$ is a nonempty subset
$\mathcal G\subseteq\Germ^\times(V)$ which is a single $\Gamma$-orbit.  All its members
have the same rank, denoted $n=\rk(\mathcal G)$, and their code tuples are the generic
points of one irreducible $k$-variety $W=\Loc(\code g)$ (\cref{rem:code-param}), and
$g\mapsto\code g$ is a bijection from $\mathcal G$ onto $W_{\Gen}$.
\end{definition}

\begin{lemma}\label{lem:genset-basic}
Let $\mathcal G$ be a generic set of germs on $V$ of rank $n$.
\begin{enumerate}[nosep]
\item For every finite tuple $w$ there is $h\in\mathcal G$ with $h\ind{}w$; any two such
  $h,h'$ are $\Gamma_{\acl(w)}$-conjugate.
\item If $g\in\mathcal G$ and $y\in V_{\Gen}$ with $y\ind{}g$, then $y\in\operatorname{Dom}(g)$
  and $g(y)\in V_{\Gen}$ with $g(y)\ind{}g$.
\item $\mathcal G^{-1}=\{g^{-1}:g\in\mathcal G\}$ is a generic set of germs of rank $n$.
\end{enumerate}
\end{lemma}
\begin{proof}
(1) The codes of elements of $\mathcal G$ are the generic points of $W$ over $k$ and
determine the germ (\cref{rem:code-param}); take a generic point of $W$ over $\acl(w)$
(\cref{lem:genexists}) and the germ with that code.  Two such codes are generic over
$\acl(w)$, hence conjugate under $\Gamma_{\acl(w)}$ (\cref{lem:generic-char}(2)), and
so are the germs (\cref{lem:germ-aut}).
(2) \cref{lem:eval}(2) with $F=\acl(\Cb g)$.  (3) $\Cb(g^{-1})=\Cb(g)$ and
$\sigma(g^{-1})=(\sigma g)^{-1}$ (\cref{lem:inverse}).
\end{proof}

\subsection{Statement}

\begin{definition}[Hrushovski--Weil hypotheses]\label{def:hw-hyp}
A generic set of germs $\mathcal G$ on $V$ is \emph{generically transitive} if
\begin{enumerate}[label=\textup{(H\arabic*)}]
\item\label{H:trans} for all $g\in\mathcal G$ and $y\in V_{\Gen}$ with $y\ind{}g$:
  $g(y)\ind{}y$;
\end{enumerate}
it is \emph{closed under inverse and generic composition} if
\begin{enumerate}[label=\textup{(H\arabic*)},start=2]
\item\label{H:inv} $g\in\mathcal G\Rightarrow g^{-1}\in\mathcal G$;
\item\label{H:comp} for all $g_1,g_2\in\mathcal G$ with $g_1\ind{}g_2$: $g_1g_2\in\mathcal G$.
\end{enumerate}
\end{definition}

\begin{lemma}[products of independent generics are generic over each factor]
\label{lem:prod-indep}
Let $\mathcal G$ be a generic set of germs of rank $n$, and $g_1,g_2\in\mathcal G$ with
$g_1\ind{}g_2$ and $g_1g_2\in\mathcal G$.  Then $g_1g_2\ind{}g_1$ and $g_1g_2\ind{}g_2$.
\end{lemma}
\begin{proof}
By \cref{lem:comp-codes}(1), $g_1g_2$ and $g_2$ are interalgebraic over $g_1$, so
$\rk(g_1g_2/g_1)=\rk(g_2/g_1)=\rk(g_2)=n=\rk(g_1g_2)$; symmetrically
$\rk(g_1g_2/g_2)=\rk(g_1/g_2)=n$.
\end{proof}

\begin{theorem}[Hrushovski--Weil for germs]\label{thm:hw}
Let $\mathcal G$ be a generically transitive generic set of germs on $V$ of rank $n$,
closed under inverse and generic composition, and let $G\le\Germ^\times(V)$ be the
subgroup generated by $\mathcal G$.  Then:
\begin{enumerate}[label=\textup{(\alph*)}]
\item\label{hw:two} $G=\mathcal G\cdot\mathcal G=\{g_1g_2:g_1,g_2\in\mathcal G\}$, and $G$ is
  $\Gamma$-invariant.
\item\label{hw:stab} (generic stabilisation) for $\gamma\in G$ and $h\in\mathcal G$ with
  $h\ind{}\gamma$: $\gamma h\in\mathcal G$, $h\gamma\in\mathcal G$, $\gamma h\ind{}\gamma$ and
  $h\gamma\ind{}\gamma$.
\item\label{hw:rank} every $\gamma\in G$ has $\rk(\gamma)\le n$, with equality iff
  $\gamma\in\mathcal G$.  In particular $\mathcal G$ is the set of elements of $G$ of maximal
  rank and it is a single $\Gamma$-orbit (``$G$ is connected with generic type
  $\mathcal G$'').
\item\label{hw:action} (generic action) for $\gamma\in G$ and $y\in V_{\Gen}$ with
  $y\ind{}\gamma$: $\gamma(y)\in V_{\Gen}$ and $\gamma(y)\ind{}\gamma$; and
  $(\gamma\delta)(y)=\gamma(\delta(y))$ whenever $y\ind{}\gamma\delta$.
\item\label{hw:space} (homogeneous space) there is a set $S$ with a transitive and
  faithful action of $G$, and an injection $\iota:V_{\Gen}\to S$ such that
  $\gamma\cdot\iota(y)=\iota(\gamma(y))$ whenever $y\ind{}\gamma$; every element of $S$ is
  of the form $\gamma\cdot\iota(y)$.  Explicitly $S=(G\times V_{\Gen})/E$ where
  \[
    (\gamma,y)\,E\,(\gamma',y')\iff (h\gamma)(y)=(h\gamma')(y')
    \text{ for some (equivalently every) } h\in\mathcal G\text{ with } h\ind{}\gamma\gamma'yy',
  \]
  $\delta\cdot[\gamma,y]=[\delta\gamma,y]$ and $\iota(y)=[\mathrm{id},y]$.
\end{enumerate}
\end{theorem}

\subsection{Proof}

We use freely that for $\delta,\gamma\in\Germ^\times(V)$: $\Cb(\delta\gamma)\subseteq
\acl(\delta,\gamma)$ and $\Cb(\delta)\subseteq\acl(\delta\gamma,\gamma)$ (\cref{lem:comp-codes}(1)),
so that \emph{$\delta\gamma$ and $\delta$ are interalgebraic over $\gamma$}; consequently
$\rk(\delta\gamma/C\gamma)=\rk(\delta/C\gamma)$ for every $C$.  Whenever \ref{H:comp} is
applied, \cref{lem:prod-indep} gives in addition that the product is independent from
each factor.

\begin{lemma}[three generators reduce to two]\label{lem:threetwo}
For $g_1,g_2,g_3\in\mathcal G$ there are $g_5,g_6\in\mathcal G$ with $g_1g_2g_3=g_5g_6$.
\end{lemma}
\begin{proof}
Choose $g_0\in\mathcal G$ with $g_0\ind{}g_1g_2g_3$ (\cref{lem:genset-basic}(1)).  Then
$g_1g_2g_3=(g_1g_0)(g_0^{-1}g_2g_3)$.  As $g_0^{-1}\in\mathcal G$ (\ref{H:inv}) and
$g_0^{-1}\ind{}g_2$, \ref{H:comp} and \cref{lem:prod-indep} give $g'':=g_0^{-1}g_2\in\mathcal G$
with $g''\ind{}g_2$.
Since $g''$ and $g_0$ are interalgebraic over $g_2$,
$\rk(g''/g_2g_3)=\rk(g_0/g_2g_3)=n$, so $g''\ind{}g_3$ and $g''g_3\in\mathcal G$ by
\ref{H:comp}.  Also $g_1\ind{}g_0$ gives $g_1g_0\in\mathcal G$.  Take $g_5=g_1g_0$,
$g_6=g''g_3$.
\end{proof}

\begin{proof}[Proof of \cref{thm:hw}\ref{hw:two}]
Let $P=\mathcal G\mathcal G$.  $P$ contains $\mathrm{id}=gg^{-1}$ and $\mathcal G$: for
$g\in\mathcal G$ pick $h\in\mathcal G$ with $h\ind{}g$; then $gh\in\mathcal G$ and
$g=(gh)h^{-1}\in P$.  $P$ is closed under inverses by \ref{H:inv}, and under products:
$(g_1g_2)(g_3g_4)=(g_5g_6)g_4=g_5(g_6g_4)$ is a product of two elements of $\mathcal G$ by
two applications of \cref{lem:threetwo}.  Hence $P$ is a subgroup containing
$\mathcal G$, so $P=G$.  $\Gamma$-invariance: $\sigma(g_1g_2)=\sigma g_1\,\sigma g_2$ and
$\mathcal G$ is $\Gamma$-invariant.
\end{proof}

\begin{proof}[Proof of \ref{hw:stab}]
Let $\gamma\in G$ and $h\in\mathcal G$ with $h\ind{}\gamma$.  By \ref{hw:two} write
$\gamma=g_1g_2$ with $g_i\in\mathcal G$.  The pair $(g_1,g_2)$ may be replaced by any
$\Gamma_{\acl(\Cb\gamma)}$-conjugate pair (its product is still $\gamma$, as such
automorphisms fix $\gamma$ by \cref{lem:germ-aut}), so by \cref{lem:fresh-conj} applied
to the code pair over $\acl(\Cb\gamma)$ we may assume
$g_1g_2\ind{\gamma}h$; together with $h\ind{}\gamma$ this gives $h\ind{}g_1g_2\gamma$
(\ref{R:trans}), in particular $h\ind{}g_1g_2$.  Then $g_2h\in\mathcal G$ by \ref{H:comp}, and
$\rk(g_2h/g_1)\ge\rk(g_2h/g_1g_2)=\rk(h/g_1g_2)=n$ (interalgebraicity of $g_2h$ and $h$
over $g_2$), so $g_2h\ind{}g_1$ and $\gamma h=g_1(g_2h)\in\mathcal G$ by \ref{H:comp}.
Moreover $\rk(\gamma h/\gamma)\ge\rk(\gamma h/g_1g_2)=\rk(h/g_1g_2)=n$ since $\gamma h$ and
$h$ are interalgebraic over $g_1g_2$ (which contains the code of $\gamma$); hence
$\gamma h\ind{}\gamma$.  For $h\gamma$ apply this to $\gamma^{-1}$ and $h^{-1}$
(\ref{H:inv}, $\Cb(h^{-1})=\Cb(h)$, $\Cb(\gamma^{-1})=\Cb(\gamma)$): $\gamma^{-1}h^{-1}\in\mathcal G$
and $\gamma^{-1}h^{-1}\ind{}\gamma^{-1}$, i.e.\ $h\gamma\in\mathcal G$ and $h\gamma\ind{}\gamma$.
\end{proof}

\begin{proof}[Proof of \ref{hw:rank}]
Let $\gamma\in G$ and pick $h\in\mathcal G$ with $h\ind{}\gamma$.  By \ref{hw:stab}
$h\gamma\in\mathcal G$, so $\rk(h\gamma)=n$.  Since $h\gamma$ and $\gamma$ are interalgebraic
over $h$, $\rk(\gamma)=\rk(\gamma/h)=\rk(h\gamma/h)\le\rk(h\gamma)=n$.  If $\rk(\gamma)=n$ then
$\rk(h\gamma/h)=n=\rk(h\gamma)$, i.e.\ $h\gamma\ind{}h$, so $h^{-1}\ind{}h\gamma$ and
$\gamma=h^{-1}(h\gamma)\in\mathcal G$ by \ref{H:inv}, \ref{H:comp}.  Conversely elements
of $\mathcal G$ have rank $n$.
\end{proof}

\begin{proof}[Proof of \ref{hw:action}]
The first claim is \cref{lem:eval}(2) with $F=\acl(\Cb\gamma)$.  The second is the
defining property of the composite $\gamma\circ\delta$ in \cref{lem:comp}, since
$y\ind{}\gamma\delta$ says that $y$ is generic over $\acl(\Cb\gamma\,\Cb\delta)$.
\end{proof}

\begin{lemma}[the key rank computation]\label{lem:keyrank}
Let $\gamma\in G$, $h\in\mathcal G$ and $w$ a tuple with $h\ind{}\gamma w$.  Then
$h\gamma\in\mathcal G$ and $h\gamma\ind{}w$; in particular, if $y\in V_{\Gen}$ and
$h\ind{}\gamma y$ then $y\in\operatorname{Dom}(h\gamma)$.
\end{lemma}
\begin{proof}
$h\gamma\in\mathcal G$ by \ref{hw:stab}.  As $h\gamma$ and $h$ are interalgebraic over
$\gamma$, $\rk(h\gamma/w)\ge\rk(h\gamma/w\gamma)=\rk(h/w\gamma)=n$.
\end{proof}

\begin{proof}[Proof of \ref{hw:space}]
\emph{$E$ is well defined.}  Let $\gamma,\gamma'\in G$, $y,y'\in V_{\Gen}$ and
$w:=\gamma\gamma'yy'$.  By \cref{lem:genset-basic}(1) there is at least one $h\in\mathcal G$
with $h\ind{}w$ (so the ``for some'' and ``for every'' forms are equivalent once the
condition is shown to be independent of $h$; a formalization should define $E$ by the
universal form).  Let $h,\tilde h\in\mathcal G$ both be independent from $w$.  By
\cref{lem:keyrank}, $y\in\operatorname{Dom}(h\gamma)$ and $y'\in\operatorname{Dom}(h\gamma')$, so
both sides of the defining equation are defined; likewise for $\tilde h$.  By
\cref{lem:genset-basic}(1) there is $\sigma\in\Gamma_{\acl(w)}$ with $\sigma h=\tilde h$;
$\sigma$ fixes the codes of $\gamma,\gamma'$ and the points $y,y'$, hence (\cref{lem:germ-aut},
\cref{lem:eval}(4)) transforms the equation for $h$ into the equation for $\tilde h$.
So the condition does not depend on $h$.

\emph{$E$ is an equivalence relation.}  Reflexivity and symmetry are clear.  For
transitivity, given $(\gamma,y)E(\gamma',y')$ and $(\gamma',y')E(\gamma'',y'')$ choose one
$h\in\mathcal G$ independent from all six objects and compare.

\emph{The action is well defined.}  Let $(\gamma,y)E(\gamma',y')$ and $\delta\in G$.
Choose $\tilde h\in\mathcal G$ with $\tilde h\ind{}\gamma\gamma'yy'\delta$.  By
\cref{lem:keyrank} (with $w=\gamma\gamma'yy'$), $\tilde h\delta\in\mathcal G$ and
$\tilde h\delta\ind{}\gamma\gamma'yy'$, so $\tilde h\delta$ is an admissible witness for
$(\gamma,y)E(\gamma',y')$: $(\tilde h\delta\gamma)(y)=(\tilde h\delta\gamma')(y')$.  Since
$\tilde h\ind{}\delta\gamma\,\delta\gamma'\,yy'$ as well (its rank over these is still $n$
because they lie in $\acl$ of the previous tuple), $\tilde h$ witnesses
$(\delta\gamma,y)E(\delta\gamma',y')$.  Associativity and unit law are inherited from $G$.

\emph{$\iota$ is injective.}  If $(\mathrm{id},y)E(\mathrm{id},y')$ then $h(y)=h(y')$ for
$h\in\mathcal G$ independent from $yy'$, and $y=h^{-1}(h(y))=h^{-1}(h(y'))=y'$ by
\cref{lem:comp} ($y,y'$ are generic over $\acl(\Cb h)=\acl(\Cb h^{-1})$).

\emph{Compatibility.}  If $y\ind{}\gamma$, choose $h\ind{}\gamma y\gamma(y)$; then
$y\ind{}\gamma h$ (\ref{R:trans}: $y\ind{}\gamma$ and $y\ind{\gamma}h$, the latter from
$h\ind{}\gamma y$), so $(h\gamma)(y)=h(\gamma(y))$ by \ref{hw:action}, i.e.\
$(\gamma,y)E(\mathrm{id},\gamma(y))$, which is $\gamma\cdot\iota(y)=\iota(\gamma(y))$.

\emph{Transitivity.}  Since $[\gamma,y]=\gamma\cdot\iota(y)$ it suffices to carry $\iota(y)$
to $\iota(y')$ for $y,y'\in V_{\Gen}$.  Choose $c\in V_{\Gen}$ with $c\ind{}yy'$.
The pairs $(y,c)$ and $(y',c)$ are independent generic pairs of $V$.  By \ref{H:trans}
and \cref{lem:genset-basic}(2), for any $g_0\in\mathcal G$ with $g_0\ind{}y$ the pair
$(y,g_0(y))$ is an independent generic pair; by \cref{lem:indep-conj} there is
$\sigma\in\Gamma$ with $\sigma(y,g_0(y))=(y,c)$, and then $g:=\sigma g_0\in\mathcal G$
satisfies $g\ind{}y$ and $g(y)=c$ (\cref{lem:eval}(4)).  Likewise there is $g'\in\mathcal G$
with $g'\ind{}c$ and $g'(c)=y'$.  By compatibility, $g\cdot\iota(y)=\iota(c)$ and
$g'\cdot\iota(c)=\iota(y')$, so $(g'g)\cdot\iota(y)=\iota(y')$.

\emph{Faithfulness.}  Suppose $\delta\in G$ fixes every element of $S$.  Choose
$y\in V_{\Gen}$ with $y\ind{}\delta$ and $h\in\mathcal G$ with $h\ind{}\delta y$.  From
$(\delta,y)E(\mathrm{id},y)$ we get $(h\delta)(y)=h(y)$.  Now $h\delta,h\in\mathcal G$ and
$y\ind{}h\delta$ (as $y\ind{}\delta$ and $h\ind{}\delta y$ give $y\ind{}\delta h$), so by
\cref{lem:eval}(5) $h\delta=h$ and $\delta=\mathrm{id}$.
\end{proof}

\begin{remark}[what ``connected'' and ``homogeneous space of rank $m$'' mean here]
\label{rem:connected}
Statement \ref{hw:rank} is the germ-level content of connectedness: the elements of
maximal rank form one $\Gamma$-orbit, and \ref{hw:stab} says that this orbit is
stabilised by left and right multiplication by independent elements.  The set $S$
contains (a copy of) $V_{\Gen}$, on which $G$ acts generically by \ref{hw:action}:
elements of $G$ move generic points of $V$ independent from them to generic points of
$V$ independent from them.  We define no rank on $S$ and make no claim about ``generic
elements of $S$''.  Every element of $S$ is $\gamma\cdot\iota(y)$ for some $\gamma\in G$ and
$y\in V_{\Gen}$ (by definition of $S$), but $y$ can \emph{not} in general be taken
independent from $\gamma$: for the group of translations of the affine line
($V=\mathbb A^1$, $\mathcal G$ the germs $y\mapsto y+a$ with $a$ transcendental, $G$ all
translations, $(\gamma,y)\,E\,(\gamma',y')$ iff $\gamma+y=\gamma'+y'$, so $S\cong K$ and
$\iota(V_{\Gen})$ is the set of transcendental elements) the class $[-t,t]=0$ has no
presentation $\gamma\cdot\iota(y)$ with $y\ind{}\gamma$, since by compatibility such a
presentation would equal $\iota(\gamma(y))\in\iota(V_{\Gen})$.  No later statement needs
such presentations.  In \cref{sec:constructibility} both $G$ and $S$ are shown to be
constructible quotients.
\end{remark}

\begin{remark}[comparison with the literature]
\Cref{thm:hw} is Theorem 5.1 of Bays' Hrushovski--Weil notes (= Lemma 5.2 of Bays'
geometric stability theory notes, which follow Pillay's \emph{Geometric Stability
Theory}, Chapter 5, \S1; Hrushovski's thesis, Theorem 3.1.1 together with the
Proposition on p.~37) written for germs of quasi-rational maps.  The only
``stability'' used is the uniqueness of generic points up to conjugacy
(\cref{lem:generic-char}, \cref{lem:indep-conj}) and the rank calculus.  The hypothesis
\ref{H:inv} replaces the chain-condition argument of the thesis, which is unnecessary in
the group-configuration application since inverses are available there by construction.
\end{remark}

\section{Hrushovski--Weil for families of bijections}
\label{sec:bij}

Let $V,V'$ be irreducible $k$-varieties.  A \emph{generic set of invertible germs
from $V$ to $V'$} is a nonempty $\Gamma$-orbit $\mathcal F\subseteq\Germ^\times(V,V')$
(invertible germs $V\to V'$, \cref{lem:inverse}); its members have a common rank
$n=\rk(\mathcal F)$ and $\dim V=\dim V'=:m$.  It is \emph{generically transitive} if
$f(y)\ind{}y$ for all $f\in\mathcal F$ and $y\in V_{\Gen}$ with $y\ind{}f$
(exactly as \ref{H:trans}).  \Cref{lem:genset-basic} holds verbatim for $\mathcal F$.

\begin{definition}[the difference set]\label{def:diff}
For $f_1,f_2\in\mathcal F$ with $f_1\ind{}f_2$ put $\delta(f_1,f_2)=f_1^{-1}f_2\in\Germ^\times(V)$,
and let $\mathcal G(\mathcal F)$ be the set of all such germs.
\end{definition}

\begin{lemma}\label{lem:diff-basic}
\begin{enumerate}[nosep]
\item $\mathcal G(\mathcal F)$ is a single $\Gamma$-orbit, hence a generic set of germs on $V$.
\item For $f_1\ind{}f_2$ in $\mathcal F$ and $\delta=\delta(f_1,f_2)$: $\delta$ and $f_2$ are
  interalgebraic over $f_1$, $\delta$ and $f_1$ are interalgebraic over $f_2$, and
  $\rk(\delta)\ge n$.  The following are equivalent: $\rk(\delta)=n$; $\delta\ind{}f_1$;
  $\delta\ind{}f_2$.
\end{enumerate}
\end{lemma}
\begin{proof}
(1) Independent pairs $(f_1,f_2)$ correspond to independent pairs of generic points of
the code variety $W$, which are all $\Gamma$-conjugate (\cref{lem:indep-conj}); and
$\sigma(f_1^{-1}f_2)=(\sigma f_1)^{-1}\sigma f_2$.
(2) $\Cb(\delta)\subseteq\acl(f_1,f_2)$, $\Cb(f_2)\subseteq\acl(f_1,\delta)$ (as $f_2=f_1\delta$),
$\Cb(f_1)\subseteq\acl(f_2,\delta)$ (as $f_1=f_2\delta^{-1}$), by \cref{lem:comp-codes}(1)
applied to $\delta:V\to V$ and $f_1:V\to V'$.  Hence
$\rk(\delta/f_1)=\rk(f_2/f_1)=n$ and $\rk(\delta)\ge n$; and $\rk(\delta)=n$ iff
$\rk(\delta)=\rk(\delta/f_1)$ iff $\delta\ind{}f_1$.  Symmetrically for $f_2$.
\end{proof}

\begin{theorem}[Hrushovski--Weil for bijections]\label{thm:hw-bij}
Let $\mathcal F$ be a generically transitive generic set of invertible germs from $V$
to $V'$ of rank $n$, and suppose
\begin{enumerate}[label=\textup{(F)}]
\item\label{F:rank} $\rk(\delta)=n$ for one (equivalently every) $\delta\in\mathcal G(\mathcal F)$.
\end{enumerate}
Then $\mathcal G:=\mathcal G(\mathcal F)$ is a generically transitive generic set of germs on
$V$ of rank $n$ closed under inverse and generic composition, so \cref{thm:hw} applies to
it.  Moreover, for $f\in\mathcal F$ and $\gamma$ in the group $G$ generated by $\mathcal G$
with $\gamma\ind{}f$: $f\gamma\in\mathcal F$ and $f\gamma\ind{}\gamma$.
\end{theorem}

\begin{proof}
\emph{Inverse \ref{H:inv}.}  $\delta(f_1,f_2)^{-1}=f_2^{-1}f_1=\delta(f_2,f_1)$.

\emph{Generic composition \ref{H:comp}.}  Let $\delta,\delta'\in\mathcal G$ with
$\delta\ind{}\delta'$, and choose $f\in\mathcal F$ with $f\ind{}\delta\delta'$
(\cref{lem:genset-basic}(1)).  Write $\delta=f_1^{-1}f_2$ with $f_1\ind{}f_2$.
By \cref{lem:diff-basic}(2) and \ref{F:rank}, $(f_2,\delta)$ is an independent pair
consisting of a member of $\mathcal F$ and a member of $\mathcal G$; so is $(f,\delta)$.
By \cref{lem:indep-conj} (applied to the code varieties of $\mathcal F$ and $\mathcal G$)
there is $\sigma\in\Gamma$ with $\sigma(f_2,\delta)=(f,\delta)$.  Put $f'_1=\sigma f_1\in\mathcal F$;
then $f'^{-1}_1f=\sigma(f_1^{-1}f_2)=\sigma\delta=\delta$.  Similarly $(f_1,\delta)$ and
$(f,\delta')$ are independent $\mathcal F\times\mathcal G$ pairs, so there is $\sigma'$ with
$\sigma'(f_1,\delta)=(f,\delta')$; put $f'_2=\sigma'f_2$, so $f^{-1}f'_2=\delta'$.  Hence
\[
   \delta\delta'=f'^{-1}_1\,f\,f^{-1}\,f'_2=f'^{-1}_1f'_2 .
\]
It remains to check that $f'_1\ind{}f'_2$.  First, $f'_1\ind{}f$ and $f'_2\ind{}f$ (images of
independent pairs).  Next, $f'_1$ and $\delta$ are interalgebraic over $f$
($\delta=f'^{-1}_1f$, $f'_1=f\delta^{-1}$), and $f'_2$ and $\delta'$ are interalgebraic
over $f$; and $\delta\ind{f}\delta'$ because $\delta\ind{}\delta'$ and $f\ind{}\delta\delta'$
(\ref{R:trans}: $\rk(\delta/f\delta')=\rk(\delta/\delta')=\rk(\delta)=\rk(\delta/f)$).
Hence $f'_1\ind{f}f'_2$ by \ref{R:aclinv}, and with $f'_1\ind{}f$ we get $f'_1\ind{}ff'_2$
(\ref{R:trans}), so $f'_1\ind{}f'_2$.  Therefore $\delta\delta'=\delta(f'_1,f'_2)\in\mathcal G$.
(That $\delta\delta'$ is then independent from $\delta$ and from $\delta'$ is
\cref{lem:prod-indep}.)

\emph{Generic transitivity \ref{H:trans}.}  Let $\delta\in\mathcal G$ and $y\in V_{\Gen}$
with $y\ind{}\delta$.  The set of pairs $(f_1,f_2)$ with $f_1\ind{}f_2$ and
$f_1^{-1}f_2=\delta$ is $\Gamma_{\acl(\Cb\delta)}$-invariant (\cref{lem:germ-aut}) and
nonempty, so by \cref{lem:fresh-conj} we may choose one with $f_1f_2\ind{\delta}y$; then
$y\ind{}\delta f_1f_2$, in particular $y\ind{}f_2$ and $y\ind{}f_1f_2$.  Put $c=f_2(y)$;
by generic transitivity of $\mathcal F$, $c\ind{}y$, and $c\in\acl(y,f_2)$.  From
$y\ind{}f_1f_2$ and $f_1\ind{}f_2$ we get $\rk(f_1/f_2y)=\rk(f_1f_2y)-\rk(f_2y)=
(m+2n)-(m+n)=n$, so $f_1\ind{}f_2y$ and hence $f_1\ind{}yc$.  Then $c\ind{}y$ and
$c\ind{y}f_1$ (from $f_1\ind{}yc$) give $c\ind{}yf_1$ (\ref{R:trans}), so $c$ is a generic
point of $V'$ over $F:=\acl(\Cb f_1,y)$.  By \cref{lem:eval}(2) applied to $f_1^{-1}$ and
$F$, $z:=f_1^{-1}(c)$ is a generic point of $V$ over $F$, so $z\ind{}y$.  Finally
$z=f_1^{-1}(f_2(y))=\delta(y)$ by \cref{lem:comp}, as $y$ is generic over
$\acl(\Cb f_1\,\Cb f_2)$.

\emph{The last claim.}  First let $f\in\mathcal F$ and $\delta\in\mathcal G$ with
$\delta\ind{}f$.  Fix any independent $f_1\ind{}f_2$ in $\mathcal F$ and $\delta_0=f_1^{-1}f_2$;
by \ref{F:rank} and \cref{lem:diff-basic}(2), $(f_1,\delta_0)$ is an independent
$\mathcal F\times\mathcal G$ pair, as is $(f,\delta)$, so \cref{lem:indep-conj} gives
$\sigma\in\Gamma$ with $\sigma(f_1,\delta_0)=(f,\delta)$.  Then $f\delta=\sigma(f_1\delta_0)=\sigma f_2\in\mathcal F$,
$f\delta\ind{}f$ (image of $f_2\ind{}f_1$) and $f\delta\ind{}\delta$ (image of
$f_2\ind{}\delta_0$, which is \ref{F:rank}).  Now let $\gamma\in G$ with $\gamma\ind{}f$.  By
\cref{thm:hw}\ref{hw:two} write $\gamma=\delta_1\delta_2$ with $\delta_i\in\mathcal G$,
conjugated over $\acl(\Cb\gamma)$ (\cref{lem:fresh-conj}) so that
$\delta_1\delta_2\ind{\gamma}f$; then
$f\ind{}\gamma\delta_1\delta_2$.  By the first part, $f_2:=f\delta_1\in\mathcal F$ with
$f_2$ and $f$ interalgebraic over $\delta_1$, so
$\rk(f_2/\delta_2)\ge\rk(f_2/\delta_1\delta_2)=\rk(f/\delta_1\delta_2)=n$, i.e.\ $f_2\ind{}\delta_2$,
and again by the first part $f\gamma=f_2\delta_2\in\mathcal F$.  Finally $f\gamma$ and $f$
are interalgebraic over $\gamma$ (\cref{lem:comp-codes}(1)), so
$\rk(f\gamma/\gamma)\ge\rk(f\gamma/\delta_1\delta_2)=\rk(f/\delta_1\delta_2)=n$ (as
$\Cb(\gamma)\subseteq\acl(\delta_1\delta_2)$), i.e.\ $f\gamma\ind{}\gamma$.
\end{proof}

\begin{remark}
In the proof of generic composition the only use of \ref{F:rank} is to know that
$(f_2,\delta)$ and $(f_1,\delta)$ are independent pairs, which is what allows
``re-basing'' the representation $\delta=f_1^{-1}f_2$ on the fresh element $f$
(it is used once more, in the last claim, for $f\delta\ind{}\delta$).  In the
finite-rank setting \ref{F:rank} is exactly the hypothesis ``$c\ind{}b_i$'' of Bays'
Lemma 5.4, by \cref{lem:diff-basic}(2).  In \cref{sec:proof} it is verified for the
configuration using the non-degeneracy clause.
\end{remark}

\section{Group configurations}
\label{sec:gc}

\subsection{Definition and first consequences}

\begin{definition}[group configuration]\label{def:gc}
A \emph{group configuration over $k$} is a $6$-tuple $(a,b,c,x,y,z)$ of tuples of
$K$ such that, with the \emph{lines}
\[
   L_1=\{a,b,c\},\quad L_2=\{a,x,y\},\quad L_3=\{b,y,z\},\quad L_4=\{c,z,x\},
\]
the following hold:
\begin{enumerate}[label=\textup{(GC\arabic*)}]
\item\label{GC:group} $\acl(ab)=\acl(ac)=\acl(bc)$;
\item\label{GC:space} $\acl(ax)=\acl(ay)$, $\acl(by)=\acl(bz)$, $\acl(cx)=\acl(cz)$;
\item\label{GC:indep} every triple of points not contained in a line is independent:
  $\rk(pqr)=\rk(p)+\rk(q)+\rk(r)$ for all non-collinear $\{p,q,r\}$;
\item\label{GC:cb} $a\subseteq\acl(\Cb(xy/a))$, $b\subseteq\acl(\Cb(yz/b))$, $c\subseteq\acl(\Cb(xz/c))$;
\item\label{GC:rank} $\rk(a)\ge1$.
\end{enumerate}
The points $a,b,c$ are the \emph{group points}, $x,y,z$ the \emph{space points}.
\end{definition}

Two distinct lines meet in exactly one point ($L_1\cap L_2=\{a\}$, $L_1\cap L_3=\{b\}$,
$L_1\cap L_4=\{c\}$, $L_2\cap L_3=\{y\}$, $L_2\cap L_4=\{x\}$, $L_3\cap L_4=\{z\}$), each
point lies on exactly two lines, and any two points lie on at most one line (the three
pairs $\{a,z\}$, $\{b,x\}$, $\{c,y\}$ lie on no line).  The group line $L_1$ is
distinguished by the definition (\ref{GC:group} and \ref{GC:cb} concern group points
only), and the definition is invariant under the six relabelings induced by
permutations of the three space lines $L_2,L_3,L_4$: these permute $a,b,c$ and
correspondingly $x,y,z$ (for instance $L_3\leftrightarrow L_4$ is $b\leftrightarrow c$,
$y\leftrightarrow x$, $z$ fixed), and each of \ref{GC:group}--\ref{GC:cb} is symmetric in
the two space points of a line and in the three group points (\ref{GC:cb} because $\FOD$
is invariant under renaming variables, so $\Cb(xy/a)=\Cb(yx/a)$); \ref{GC:rank} is
invariant because $\rk(a)=\rk(b)=\rk(c)$ by \cref{lem:gc-ranks}.

\begin{lemma}[ranks]\label{lem:gc-ranks}
Let $(a,b,c,x,y,z)$ be a group configuration.  Then $\rk(a)=\rk(b)=\rk(c)=:n$,
$\rk(x)=\rk(y)=\rk(z)=:m$, $1\le m\le n$; all pairs of points are independent;
$\rk(abc)=2n$; $\rk(axy)=\rk(byz)=\rk(cxz)=n+m$; $\rk(abcxyz)=2n+m$; and each point
is independent from each line not through it: for instance $a\ind{}byz$, $a\ind{}cxz$,
$b\ind{}axy$, $b\ind{}cxz$, $c\ind{}axy$, $c\ind{}byz$, $x\ind{}abc$, $x\ind{}byz$, and
so on.  Moreover $m\ge 1$ follows from $n\ge1$.
\end{lemma}
\begin{proof}
Every pair lies in a non-collinear triple (a triple containing a given pair and lying
on a line lies on the unique line through that pair, if any; so any third point off that
line works), so pairs are independent by \ref{GC:indep}.  By \ref{GC:group}, $\rk(ab)=\rk(ac)$, i.e.\ $n_b=n_c$,
and similarly all three agree.  By \ref{GC:space}, $\rk(ax)=\rk(ay)$ so $\rk(x)=\rk(y)$,
etc.  $\rk(axy)=\rk(ay)=n+m\ge\rk(xy)=2m$ gives $m\le n$.  For $\rk(abcxyz)$: $(a,b,y)$ is
non-collinear, so $\rk(aby)=2n+m$, and $c\in\acl(ab)$, $x\in\acl(ay)$, $z\in\acl(by)$.
Independence of a point $p$ from a line $\ell$ not through $p$: pick two points $q,r$
of $\ell$, taking $q$ to be the group point of $\ell$ if $\ell$ is a space line; then
$(p,q,r)$ is non-collinear (a line containing $q$ and $r$ is $\ell$, as two distinct
lines share only one point, and $p\notin\ell$), so $p\ind{}qr$, and the third point of $\ell$ is algebraic over $qr$ by \ref{GC:group} (if
$\ell=L_1$) or \ref{GC:space} (a space point is algebraic over the group point and the
other space point of its line); hence $p\ind{}\ell$.  E.g.\ $a\ind{}by$ and
$z\in\acl(by)$ give $a\ind{}byz$.  If $m=0$
then $x,y\subseteq k$, $P_{\acl(a)}(xy)$ is a $k$-rational point and $\Cb(xy/a)=k$,
contradicting \ref{GC:cb} and $n\ge1$.
\end{proof}

\begin{lemma}[invariance]\label{lem:gc-invariance}
Let $\mathfrak c=(a,b,c,x,y,z)$ be a group configuration over $k$.
\begin{enumerate}[nosep]
\item For $\sigma\in\Gamma$, $\sigma\mathfrak c$ is a group configuration over $k$.
\item (base change) If $e$ is a tuple with $e\ind{}abcxyz$, then $\mathfrak c$ is a group
  configuration over $k'=\acl(k,e)$ (all notions of
  \cref{def:closures,def:cb} now taken relative to the base $k'$; $K$ still has
  infinite transcendence degree over $k'$).
\item (interalgebraic replacement) If $a'\interalg a,\dots,z'\interalg z$ then
  $(a',b',c',x',y',z')$ is a group configuration over $k$.
\end{enumerate}
\end{lemma}
\begin{proof}
(1) \cref{lem:equiv} and \ref{Cb:equiv}.
(2) \ref{GC:group}, \ref{GC:space}: e.g.\ $c\in\acl(ab)\subseteq\acl(k'ab)$ and
symmetrically.  \ref{GC:indep}: for a non-collinear triple, $\rk(pqr/e)=\rk(pqr)$ and
$\rk(p/e)=\rk(p)$ etc.\ by \ref{R:basechange}.  \ref{GC:cb}: the canonical base over the
new base is $\Cb'(xy/a)=k'\cdot\FOD(P_{\acl(k'a)}(xy))$ and
$\FOD(P_{\acl(ea)}(xy))=\FOD(P_{\acl(a)}(xy))$ by \ref{Cb:basechange} (as
$e\ind{}xya$); hence $\acl(\Cb'(xy/a))=\acl(k',\Cb(xy/a))\supseteq\acl(\Cb(xy/a))\ni a$.
\ref{GC:rank}: $\rk(a/e)=\rk(a)$.
(3) \ref{GC:group}--\ref{GC:indep} and \ref{GC:rank} are invariant under replacing
points by interalgebraic ones (\ref{R:aclinv}).  For \ref{GC:cb}:
$\Cb(x'y'/a')=\Cb(x'y'/\acl(a))$ and $\acl(\Cb(x'y'/\acl a))=\acl(\Cb(xy/\acl a))$ by
\ref{Cb:interalg}, which contains $a$, hence $a'$.
\end{proof}

\begin{lemma}[conjugate sets of a space point]\label{lem:conjset-gc}
Let $(a,b,c,x,y,z)$ be a group configuration and let $r$ be a space point with the two
lines $\{p_1,q_1,r\}$ and $\{p_2,q_2,r\}$ through it, $p_1,p_2$ being their group points
and $q_1,q_2$ their other space points.  Put $L=k(p_1q_1p_2q_2)$.  Then
$r\subseteq\acl(L)$, every element of the conjugate set $\Gamma_L\cdot r$ lies in
$\acl(r)$, and the code $\tilde r=\code{\Gamma_L r}$ satisfies $\tilde r\subseteq\dcl(L)$
and $\tilde r\interalg r$.
\end{lemma}
\begin{proof}
$r\in\acl(p_1q_1)\cap\acl(p_2q_2)$ by \ref{GC:space} (a space point is algebraic over
the group point and the other space point of each line through it).  The group points
$p_1,p_2$ lie on different lines through $r$, and the line through $p_1,p_2$ is $L_1$,
which does not contain $r$; so $(r,p_1,p_2)$ is non-collinear, $p_1\ind{}rp_2$ and hence
$p_1\ind{r}p_2$.  Since $q_1\in\acl(rp_1)$ and $q_2\in\acl(rp_2)$ (\ref{GC:space} again),
\ref{R:aclinv} gives $p_1q_1\ind{r}p_2q_2$, and \ref{R:meet} gives
$\acl(rp_1q_1)\cap\acl(rp_2q_2)=\acl(r)$.  For $\sigma\in\Gamma_L$ we have
$\sigma r\in\sigma\acl(p_1q_1)\cap\sigma\acl(p_2q_2)=\acl(p_1q_1)\cap\acl(p_2q_2)\subseteq\acl(r)$.
The claims about $\tilde r$ are \cref{lem:conjcode}.  (Note that a group point need
\emph{not} be algebraic over the two space points of its line; only the space-point
clauses of \ref{GC:space} were used.)
\end{proof}

\begin{remark}
For a group point the same statement is false in general (cf.\ \cref{rem:cb-not-acl});
this is why the definability step of \cref{sec:proof} only ever codes space points.
\end{remark}

\subsection{Configurations of a homogeneous space}

\begin{proposition}[the converse]\label{prop:converse}
Let $\mathcal G$ be a generically transitive generic set of germs on $V$ closed under
inverse and generic composition, with group $G$ (\cref{thm:hw}), $n=\rk(\mathcal G)$,
$m=\dim V\ge1$.  Let $g,h\in\mathcal G$ and $y\in V_{\Gen}$ with $(g,h,y)$ independent,
and put $x=g(y)$, $z=h(y)$, $c=hg^{-1}\in G$.  Then
\[
  \bigl(\code g,\ \code h,\ \code c,\ x,\ y,\ z\bigr)
\]
is a group configuration over $k$ with $\rk(\code g)=n$ and $\rk(y)=m$; moreover
$z=c(x)$ and $h=cg$.
\end{proposition}
\begin{proof}
Write $a=\code g$, $b=\code h$ (so $\acl(a)=\acl(g)$ etc.).  Since $\Cb(hg^{-1})\subseteq\acl(g,h)$,
$\Cb(h)\subseteq\acl(c,g)$ ($h=cg$) and $\Cb(g)\subseteq\acl(c,h)$ ($g=c^{-1}h$), \ref{GC:group}
holds.  $x=g(y)\in\dcl(g,y)$ and $y=g^{-1}(x)$ ($x\in\operatorname{Dom}(g^{-1})$ because
$x$ is generic over $\acl(\Cb g)$, \cref{lem:eval}(2)), so $\acl(ax)=\acl(ay)$; similarly
$\acl(by)=\acl(bz)$.  Next, $x\ind{}gh$: $y\ind{}gh$ and $x\interalg y$ over $g$ give
$\rk(x/gh)=\rk(y/gh)=m$.  Hence $x\in\operatorname{Dom}(c)$ and by \cref{lem:comp-codes}(2)
$c(x)=h(g^{-1}(x))=h(y)=z$; so $z\in\dcl(c,x)$, $x=c^{-1}(z)$ and $\acl(cx)=\acl(cz)$.
This proves \ref{GC:space}.

\ref{GC:cb}: by \cref{lem:triple} the germ determined by the triple $(a;y,x)$ (with
$y\ind{}a$ by hypothesis, $x\ind{}a$ by \cref{lem:genset-basic}(2), $x\in\dcl(a,y)$ by
\cref{lem:eval-dcl}) is the unique germ with code in $\acl(a)$ sending $y$ to $x$, hence
is $g$ itself; so $\Cb(yx/a)=\Cb(g)=k(a)$ and, as $\FOD$ is invariant under renaming the
variables (\cref{thm:weil-fod}), $\Cb(xy/a)=\Cb(yx/a)$, whence
$a\subseteq\acl(\Cb(xy/a))$.  Likewise for $b$ with $(y,z)$ and for $c$ with $(x,z)$
(using $x\ind{}c$, $z\ind{}c$, which follow from $\rk(x/gh)=m$ and
$\rk(z/gh)=\rk(y/gh)=m$).

\ref{GC:indep}: all six points lie in $\acl(ghy)$, which has rank $2n+m$.  We check the
$16$ non-collinear triples.  Since $\acl(ab)=\acl(ac)=\acl(bc)$ it suffices, for triples
with two group points, to treat $(g,h,p)$ with $p\in\{x,y,z\}$: $\rk(ghy)=2n+m$ by
hypothesis, and $\rk(x/gh)=\rk(z/gh)=m$ was shown above.  Triples with one group point
and two space points not on its line: $(g,y,z)$, $(g,x,z)$, $(h,x,y)$, $(h,x,z)$,
$(c,x,y)$, $(c,y,z)$.  From $g\ind{}hy$ and $z\in\acl(hy)$ we get $g\ind{}yz$, so
$\rk(z/gy)=\rk(yz/g)-\rk(y/g)=\rk(yz)-m=m$, using $\rk(yz)=2m$ (which is \ref{H:trans}
for $h$); hence $\rk(gyz)=n+2m$, and $\rk(gxz)=\rk(gyz)$ as $\acl(gx)=\acl(gy)$.
Symmetrically $h\ind{}gy$ gives $\rk(hxy)=\rk(hxz)=n+2m$.  For $c$: $c$ and $h$ are
interalgebraic over $g$, so $\rk(c/xyg)=\rk(h/xyg)=\rk(h/yg)=n$ (as $x\in\acl(gy)$);
therefore $\rk(c/xy)=n$, $c\ind{}xy$ and $\rk(cxy)=n+2m$; similarly $c$ and $g$ are
interalgebraic over $h$, $\rk(c/yzh)=\rk(g/yzh)=\rk(g/yh)=n$, so $\rk(cyz)=n+2m$.
Finally $(x,y,z)$: $\rk(xyz)=\rk(y)+\rk(x/y)+\rk(z/xy)$ with $\rk(x/y)=m$ (\ref{H:trans})
and $\rk(z/xy)=\rk(z/xyg)+\bigl(\rk(g/xy)-\rk(g/xyz)\bigr)$; here $\rk(z/xyg)=\rk(z/yg)=m$
(shown above: $g\ind{}hy$ and $z\in\acl(hy)$ give $g\ind{}yz$, so $\rk(z/gy)=m$) and $\rk(g/xy)\ge\rk(g/xyz)$, so
$\rk(z/xy)\ge m$, hence $=m$ and $\rk(xyz)=3m$.  \ref{GC:rank} is $n\ge1$, which holds
since $\mathcal G$ consists of germs of rank $n$ moving generic points
(\ref{H:trans} with $m\ge1$ forces $n\ge1$: a germ of rank $0$ is $\Gamma$-invariant and
$g(y)\in\dcl(k,y)$ would then be interalgebraic with $y$, contradicting $g(y)\ind{}y$).
\end{proof}

\section{Proof of the group configuration theorem}
\label{sec:proof}

Throughout, $\mathfrak c=(a,b,c,x,y,z)$ is a group configuration over $k$ with ranks
$n=\rk(a)$, $m=\rk(y)$ (\cref{lem:gc-ranks}).  ``Over $k_1$'' means that the base field
$k$ is replaced by $k_1$ in all of \cref{sec:ud,sec:fod,sec:codes,sec:germs,sec:hw,sec:bij,sec:gc};
when the base changes we say so explicitly, and $\Gamma_{(k_1)}=\Aut(K/k_1)$.

\subsection{Step I: making the lines definable}

\begin{lemma}[definability step]\label{lem:defstep}
Let $r$ be a space point of $\mathfrak c$, let $\ell_1=\{p_1,q_1,r\}$ be one of the two
lines through $r$ ($p_1$ its group point, $q_1$ its other space point), let
$\ell_2=\{p_2,q_2,r\}$ be the other line through $r$, and let $\ell_3=\{p_3,q_1,q_2\}$
be the line through $q_1$ and $q_2$; so $\{p_1,p_2,p_3\}=\{a,b,c\}$.  Then there is a
tuple $e$ with $e\ind{}abcxyz$, and tuples $p'_2,q'_2$, such that with $k_1=\acl(k,e)$:
\begin{enumerate}[nosep]
\item $p'_2\subseteq\acl(e\,p_1)$ and $q'_2\subseteq\acl(e\,q_1)$;
\item with $L=k(p_1q_1p'_2q'_2)$ and $\tilde r=\code{\Gamma_L\,r}$: $\tilde r\subseteq\dcl(L)$
  and $\tilde r\interalg r$;
\item the tuple $\mathfrak c_1$ obtained from $\mathfrak c$ by replacing
  $p_1\mapsto(p_1,p'_2)$, $q_1\mapsto(q_1,q'_2)$, $r\mapsto\tilde r$ is a group configuration
  over $k_1$, pointwise interalgebraic over $k_1$ with $\mathfrak c$, and
  $\tilde r\in\dcl\bigl(k_1,(p_1,p'_2),(q_1,q'_2)\bigr)$;
\item if moreover $q_1\in\dcl(k,p_1,r)$ and $q_2\in\dcl(k,p_2,r)$, then
  $(q_1,q'_2)\in\dcl\bigl(k_1,(p_1,p'_2),\tilde r\bigr)$.
\end{enumerate}
\end{lemma}
\begin{proof}
Let $e$ be a generic point of $\Loc(p_3)$ over $\acl(abcxyz)$ (\cref{lem:genexists}),
so $e\ind{}abcxyz$.  By \cref{lem:gc-ranks}, $p_3\ind{}p_1q_1r$ (the line $\ell_1$ does
not contain $p_3$), so $p_3$ and $e$ are both generic points of $\Loc(p_3)$ over
$\acl(p_1q_1r)$ and there is $\sigma\in\Gamma_{\acl(p_1q_1r)}$ with $\sigma p_3=e$
(\cref{lem:generic-char}(2)).  Put $p'_2=\sigma p_2$, $q'_2=\sigma q_2$.  Then
$\sigma\mathfrak c$ is a group configuration (\cref{lem:gc-invariance}(1)) with points
$p_1,q_1,r$ (fixed), $p'_2,q'_2,e$ and lines $\ell_1$, $\{p'_2,q'_2,r\}$, $\{e,q_1,q'_2\}$,
$\{p_1,p'_2,e\}$.

(1) \ref{GC:group} for the group line $\{p_1,p'_2,e\}$ of $\sigma\mathfrak c$ gives
$p'_2\subseteq\acl(ep_1)$; \ref{GC:space} for the line $\{e,q_1,q'_2\}$ gives
$q'_2\subseteq\acl(eq_1)$.
(2) \Cref{lem:conjset-gc} applied to $\sigma\mathfrak c$, the space point $r$ and its lines
$\ell_1$, $\{p'_2,q'_2,r\}$.
(3) $\mathfrak c$ is a configuration over $k_1$ by \cref{lem:gc-invariance}(2); by (1),(2)
the three replaced points are interalgebraic over $k_1$ with the originals, so
$\mathfrak c_1$ is a configuration over $k_1$ by \cref{lem:gc-invariance}(3); and
$\dcl(L)\subseteq\dcl(k_1,(p_1,p'_2),(q_1,q'_2))$.
(4) $q'_2=\sigma q_2\in\sigma\dcl(k,p_2,r)=\dcl(k,p'_2,r)$ (\cref{cor:dcl-props}, $\sigma$
fixes $r$).  So $w:=(q_1,q'_2)\in\dcl(L_0,r)$ with $L_0=k(p_1p'_2)\subseteq L$, and
$w\subseteq L$; \cref{lem:conjcode} gives $w\in\dcl(L_0,\tilde r)$.
\end{proof}

\begin{proposition}[Step I]\label{prop:stepI}
There is a tuple $e\ind{}abcxyz$ such that, over $k'=\acl(k,e)$, there is a group
configuration $\mathfrak c'=(a',b',c',x',y',z')$ pointwise interalgebraic over $k'$ with
$\mathfrak c$ and satisfying
\[
   z'\in\dcl(k',b',y'),\qquad y'\in\dcl(k',b',z').
\]
Consequently (over $k'$): with $V=\Loc_{k'}(y')$, $V'=\Loc_{k'}(z')$ (irreducible
$k'$-varieties of dimension $m$), the triple $(b';y',z')$ determines by \cref{lem:triple}
an invertible germ $g_{b'}:V\to V'$ with $g_{b'}(y')=z'$ and $b'\interalg\code{g_{b'}}$;
replacing $b'$ by $\code{g_{b'}}$ we may assume $b'=\code{g_{b'}}$.  Then
$\mathcal F:=\Gamma_{(k')}\cdot g_{b'}$ is a generically transitive generic set of
invertible germs from $V$ to $V'$ of rank $n$ over $k'$.
\end{proposition}
\begin{proof}
Apply \cref{lem:defstep} three times.
\begin{enumerate}[nosep]
\item To $\mathfrak c$ with $r=z$, $\ell_1=\{b,y,z\}$ ($p_1=b$, $q_1=y$),
  $\ell_2=\{c,z,x\}$ ($p_2=c$, $q_2=x$), $\ell_3=\{a,x,y\}$, $p_3=a$.  Get $e_1$,
  $k_1=\acl(k,e_1)$ and $\mathfrak c_1=(a,(b,c'),c,x,(y,x'),\tilde z)$ with
  $\tilde z\in\dcl(k_1,b_1,y_1)$, where $b_1=(b,c')$, $y_1=(y,x')$.
\item Over $k_1$, to $\mathfrak c_1$ with $r=x_1=x$, $\ell_1=\{a_1,x_1,y_1\}$,
  $\ell_2=\{c_1,z_1,x_1\}$, $\ell_3=\{b_1,y_1,z_1\}$, $p_3=b_1$.  Get $e_2\ind{k_1}\mathfrak c_1$,
  $k_2=\acl(k_1,e_2)$ and $\mathfrak c_2=((a,c''),b_1,c,\tilde x,(y_1,z''),z_1)$ with
  $\tilde x\in\dcl(k_2,a_2,y_2)$.  Since $b_2=b_1$, $z_2=z_1$ and $y_2\supseteq y_1$, still
  $z_2\in\dcl(k_2,b_2,y_2)$.
\item Over $k_2$, to $\mathfrak c_2$ with $r=y_2$, $\ell_1=\{b_2,y_2,z_2\}$ ($p_1=b_2$,
  $q_1=z_2$), $\ell_2=\{a_2,x_2,y_2\}$ ($p_2=a_2$, $q_2=x_2$), $\ell_3=\{c_2,z_2,x_2\}$,
  $p_3=c_2$.  The hypotheses of \cref{lem:defstep}(4) hold: $z_2\in\dcl(k_2,b_2,y_2)$
  and $x_2\in\dcl(k_2,a_2,y_2)$.  Get $e_3$, $k_3$ and
  $\mathfrak c_3=(a_2,(b_2,a'),c_2,x_2,\tilde y,(z_2,x'''))$ with
  $\tilde y\in\dcl(k_3,b_3,z_3)$ and, by (4), $z_3=(z_2,x''')\in\dcl(k_3,b_3,\tilde y)$.
\end{enumerate}
Let $e=(e_1,e_2,e_3)$ and $k'=k_3=\acl(k,e)$.  Since $e_1\ind{}\mathfrak c$,
$e_2\ind{e_1}\mathfrak c_1$ with $\acl(e_1\mathfrak c_1)=\acl(e_1\mathfrak c)$, and
$e_3\ind{e_1e_2}\mathfrak c_2$ similarly, \ref{R:trans} gives $e\ind{}\mathfrak c$.
Interalgebraicity of $\mathfrak c_3$ with $\mathfrak c$ over $k'$ follows from the three
pointwise interalgebraicities (each over a subfield of $k'$).  Put $\mathfrak c'=\mathfrak c_3$.

For the consequences: over $k'$, $y'\ind{}b'$ and $z'\ind{}b'$ (pairs are independent),
$y'$ is a generic point of $V$ and $z'$ of $V'$ (by definition of $\Loc$),
$z'\in\dcl(b',y')$ and $y'\in\dcl(b',z')$; \cref{lem:triple} gives the invertible germ
$g_{b'}$ with $\Cb(g_{b'})=\Cb(y'z'/b')$, and \ref{GC:cb} together with
\cref{lem:triple}(3) gives $b'\interalg\code{g_{b'}}$; the replacement is allowed by
\cref{lem:gc-invariance}(3), and the two $\dcl$ clauses survive it since
$z'=g_{b'}(y')\in\dcl(\code{g_{b'}},y')$ and $y'=g_{b'}^{-1}(z')\in\dcl(\code{g_{b'}},z')$
(\cref{lem:eval-dcl}).  Generic transitivity: let $f\in\mathcal F$ and $y''\in V_{\Gen}$
with $y''\ind{}f$.  Both $(\code{g_{b'}},y')$ and $(\code f,y'')$ are independent pairs
of generic points of $W\times V$ (the former because $b'\ind{}y'$), so
\cref{lem:indep-conj} gives $\sigma\in\Gamma_{(k')}$ with
$\sigma(\code{g_{b'}},y')=(\code f,y'')$; then $f=\sigma g_{b'}$ (\cref{rem:code-param})
and $f(y'')=\sigma z'$ (\cref{lem:eval}(4)), which is independent from $\sigma y'=y''$
because $z'\ind{}y'$ (pairs independent; \cref{lem:equiv}).  Its rank is $\rk(b')=n$.
\end{proof}

\subsection{Step II: the independence hypothesis}

\begin{proposition}[Step II]\label{prop:stepII}
In the situation at the end of \cref{prop:stepI} (base $k'$, written $k$ again;
$\mathfrak c=(a,b,c,x,y,z)$ with $b=\code{g_b}$, $g_b\in\mathcal F$, $z=g_b(y)$),
hypothesis \ref{F:rank} holds for $\mathcal F$: $\rk(\delta)=n$ for every
$\delta\in\mathcal G(\mathcal F)$.
\end{proposition}
\begin{proof}
Let $b'$ be a generic point of $\Loc(b)$ over $\acl(abcxyz)$, so $b'\ind{}abcxyz$, and
let $g_{b'}\in\mathcal F$ be the germ with code $b'$.  By \cref{lem:gc-ranks},
$b\ind{}cxz$, so $b,b'$ are generic points of $\Loc(b)$ over $\acl(xcz)$ and there is
$\sigma\in\Gamma_{\acl(xcz)}$ with $\sigma b=b'$.  Put $y'=\sigma y$, $a'=\sigma a$.  Then
$\sigma\mathfrak c=(a',b',c,x,y',z)$ is a group configuration, $g_{b'}=\sigma g_b$ and
$z=\sigma z=g_{b'}(y')$ (\cref{lem:eval}(4)).

\emph{Claim 1: $y'\ind{}bb'$.}  From $b'\ind{}bz$: $zb'\ind{b'}b$, hence (as
$y'\in\acl(zb')$ by the line $\{b',y',z\}$ of $\sigma\mathfrak c$) $y'\ind{b'}b$; with
$y'\ind{}b'$ (pairs in $\sigma\mathfrak c$) this gives $y'\ind{}bb'$.

Now $y\ind{}bb'$ ($y\ind{}b$ and $b'\ind{}$ everything), so $y\in\operatorname{Dom}(g_b)$,
$z=g_b(y)$, $z\ind{}b'$ puts $z\in\operatorname{Dom}(g_{b'}^{-1})$, and
$g_{b'}^{-1}(z)=y'$ because $z=g_{b'}(y')$ with $y'\in\operatorname{Dom}(g_{b'})$.  Hence
$\delta:=g_{b'}^{-1}g_b\in\mathcal G(\mathcal F)$ satisfies $\delta(y)=y'$
(\cref{lem:comp-codes}(2)).  By \cref{lem:triple} applied to the triple $(bb';y,y')$
(hypotheses: $y\ind{}bb'$, $y'\ind{}bb'$ by Claim 1, and $y'=\delta(y)\in\dcl(bb',y)$ by
\cref{lem:eval-dcl} since $b,b'$ are the codes of $g_b,g_{b'}$), the germ with code in
$\acl(bb')$ carrying $y$ to $y'$ is unique, hence equals $\delta$, and
\[
  \Cb(\delta)=\Cb(yy'/bb').
\]
\emph{Claim 2: $\Cb(yy'/bb')=\Cb(yy'/aa')$.}  By \cref{lem:gc-ranks}, $y\ind{}ab$ and
$c\in\acl(ab)$, so $y\ind{}abc$; $b'\ind{}abcy$ gives $y\ind{abc}b'$, hence $y\ind{}abcb'$;
and $a'\in\acl(cb')$ (group line of $\sigma\mathfrak c$) gives $y\ind{}aa'bb'$.  Also
$y'\in\acl(a'x)\subseteq\acl(a'ay)$ (lines $\{a',x,y'\}$ of $\sigma\mathfrak c$ and
$\{a,x,y\}$), so $\rk(yy'/aa'bb')=\rk(y/aa'bb')=m=\rk(y/aa')=\rk(yy'/aa')$, i.e.\
$yy'\ind{aa'}bb'$; and $y'\in\acl(b'z)\subseteq\acl(b'by)$ gives likewise
$yy'\ind{bb'}aa'$.  Now \ref{Cb:two} applies.

\emph{Claim 3: $b\ind{}aa'$.}  $b'\ind{}bc$ gives $b'\ind{c}b$; since $a\in\acl(cb)$ and
$a'\in\acl(cb')$, \ref{R:aclinv} gives $ab\ind{c}a'b'$, so $a'\ind{c}ab$; as $a'\ind{}c$
(pairs in $\sigma\mathfrak c$), \ref{R:trans} gives $a'\ind{}abc$, so $a'\ind{}ab$, i.e.\
$b\ind{a}a'$; with $b\ind{}a$ this is $b\ind{}aa'$.

By Claims 2 and 3, $\Cb(\delta)\subseteq\acl(aa')$ and $b\ind{}aa'$, so $\delta\ind{}b$,
i.e.\ $\delta\ind{}g_b$.  By \cref{lem:diff-basic}(2) (with $f_1=g_{b'}$, $f_2=g_b$,
which are independent), $\rk(\delta)=n$.  Since $\mathcal G(\mathcal F)$ is a single
$\Gamma$-orbit this holds for every element.
\end{proof}

\subsection{Base change of the group}

\begin{lemma}[base change of a Hrushovski--Weil datum]\label{lem:hw-basechange}
Let $\mathcal G$ be a generically transitive generic set of germs on $V$ over $k$ closed
under inverse and generic composition, with group $G$, space $S=(G\times V_{\Gen})/E$
and embedding $\iota$ (\cref{thm:hw}).  Let $e$ be any tuple, $k''=\acl(k,e)$.  Over
$k''$, the set $\mathcal G''=\{\gamma\in\mathcal G:\gamma\ind{}e\}$ is a nonempty single
$\Gamma_{(k'')}$-orbit of invertible germs of the $k''$-variety $V$ (same ideal $I_V$), it
satisfies \ref{H:trans}--\ref{H:comp} over $k''$, it generates $G$, and the equivalence
relation $E''$ of \cref{thm:hw}\ref{hw:space} computed over $k''$ coincides with $E$ on
$G\times V''_{\Gen}$ where $V''_{\Gen}=\{y\in V_{\Gen}:y\ind{}e\}$.  Consequently $G$, its
generic set $\mathcal G''$, and $S''=(G\times V''_{\Gen})/E''\subseteq S$ (a $G$-invariant
subset on which $G$ acts transitively and faithfully, equal to $S$ by transitivity) are the
structures produced by \cref{thm:hw} over $k''$.  Likewise, for a generically transitive
generic set $\mathcal F$ of invertible germs $V\to V'$ satisfying \ref{F:rank},
$\mathcal F''=\{f\in\mathcal F:f\ind{}e\}$ has the same properties over $k''$ and
$\mathcal G(\mathcal F'')=\mathcal G''$.
\end{lemma}
\begin{proof}
First, germs over $k$ and germs over $k''$ are the same objects: a representative
$(M,f)$ over $k$ is equivalent to $(\acl(Me),f)$ (note after \cref{def:germ}), and two
representatives over fields containing $k''$ are equivalent over $k''$ iff over $k$
(both conditions quantify over the generic points of $V$ over $\acl(MM')$).  Hence
graph ideals, $\FOD$, code tuples and values $g(y)$ (\cref{lem:eval}(2), independence
of the field) are unchanged, while $\Cb''(g)=k''\cdot\Cb(g)$, $\rk''(g)=\rk(\code g/e)$,
$V''_{\Gen}$ is as stated (\cref{lem:generic-char}(1)) and
$\operatorname{Dom}''(g)=\{y\in V''_{\Gen}:y\ind{e}g\}$.
Codes of elements of $\mathcal G''$ are the generic points of $W=\Loc(\code g)$ over
$\acl(e)$, which exist and form one $\Gamma_{(k'')}$-orbit (\cref{lem:generic-char}), and
$W$ is the same variety over $k''$ (\cref{lem:dim}(2)).
Ranks over $k''$ of tuples independent from $e$ equal ranks over $k$
(\ref{R:basechange}).  \ref{H:trans}: let $\gamma\in\mathcal G''$ and $y\in V_{\Gen}$ with
$y\ind{e}\gamma$ (independence over $k''$, i.e.\ $y\ind{}e$ and $y\ind{e}\gamma$).  Then
$y\ind{}e\gamma$, and with $\gamma\ind{}e$, \ref{R:swap} gives $e\ind{}\gamma y$, hence
$e\ind{}\gamma\,y\,\gamma(y)$ ($\gamma(y)\in\acl(\gamma y)$); so $\rk(\gamma(y)/ey)=\rk(\gamma(y)/y)=m$.
\ref{H:inv} is clear.  \ref{H:comp}: if $\gamma_1,\gamma_2\in\mathcal G''$ and
$\gamma_1\ind{e}\gamma_2$ then $\gamma_1\ind{}e\gamma_2$, so $\gamma_1\gamma_2\in\mathcal G$,
and $\rk(\gamma_1\gamma_2/e)\ge\rk(\gamma_1\gamma_2/e\gamma_2)=\rk(\gamma_1/e\gamma_2)=n$
(\cref{lem:comp-codes}(1)), giving $\gamma_1\gamma_2\in\mathcal G''$.  Generation: for $\gamma\in G$ pick $h\in\mathcal G$ with
$h\ind{}\gamma e$; then $h\gamma\in\mathcal G$ and $\rk(h\gamma/e)\ge\rk(h\gamma/e\gamma)=\rk(h/e\gamma)=n$,
so $h,h\gamma\in\mathcal G''$ and $\gamma=h^{-1}(h\gamma)$.  $E''=E$: a witness $h$ for $E''$ is a
witness for $E$, and $E$ may be tested with any admissible witness.  The statements for
$\mathcal F$ are proved in the same way, using \cref{lem:diff-basic}(2) for
$\mathcal G(\mathcal F'')=\mathcal G''$: if $f_1,f_2\in\mathcal F''$, $f_1\ind{e}f_2$, then
$f_1\ind{}f_2$ and $\rk(f_1^{-1}f_2/e)\ge\rk(f_1^{-1}f_2/ef_1)=\rk(f_2/ef_1)=n$; conversely for
$\delta\in\mathcal G''$ choose a representation $\delta=f_1^{-1}f_2$ with $f_1f_2\ind{\delta}e$
(conjugating over $\acl(\Cb\delta)$, \cref{lem:fresh-conj}; the set of such pairs is
$\Gamma_{\acl(\Cb\delta)}$-invariant by \cref{lem:germ-aut}).  Then $\delta\ind{}e$
($\delta\in\mathcal G''$) and $e\ind{\delta}f_1f_2$ give $e\ind{}\delta f_1f_2$ by
\ref{R:trans}, hence $f_1f_2\ind{}e$, i.e.\ $\rk(f_1f_2/e)=2n$; so $f_1\ind{}e$,
$f_2\ind{}e$ and $\rk(f_1/ef_2)=n=\rk(f_1/e)$, i.e.\ $f_1\ind{e}f_2$, and
$\delta=f_1^{-1}f_2\in\mathcal G(\mathcal F'')$.
\end{proof}

\subsection{Step III and the main theorem}

\begin{theorem}[group configuration theorem for algebraically closed fields]\label{thm:main}
Let $\mathfrak c=(a,b,c,x,y,z)$ be a group configuration over $k$ with ranks $n,m$.  Then
there is a tuple $e$ with $e\ind{}abcxyz$ such that, over $k'=\acl(k,e)$ (all notions
relative to $k'$), there exist
\begin{enumerate}[nosep]
\item irreducible $k'$-varieties $V$, $V'$ of dimension $m$;
\item a generically transitive generic set $\mathcal F$ of invertible germs $V\to V'$ of
  rank $n$ satisfying \ref{F:rank}, with code variety $W=\Loc_{k'}(\code f)$
  ($f\in\mathcal F$), and its difference set $\mathcal G=\mathcal G(\mathcal F)$, a
  generically transitive generic set of germs on $V$ of rank $n$ closed under inverse
  and generic composition (\cref{thm:hw-bij}), with the group $G=\langle\mathcal G\rangle$,
  the transitive faithful $G$-set $S$ and the embedding $\iota:V_{\Gen}\to S$ of
  \cref{thm:hw};
\item elements $g_a,g_b\in\mathcal G$ and $y'\in V_{\Gen}$ with $(g_a,g_b,y')$ independent,
\end{enumerate}
such that, with $x'=g_a(y')$, $z'=g_b(y')$ and $g_c=g_bg_a^{-1}\in G$,
\[
   a\interalg\code{g_a},\quad b\interalg\code{g_b},\quad c\interalg\code{g_c},\quad
   x\interalg x',\quad y\interalg y',\quad z\interalg z'
   \qquad(\text{interalgebraic over }k').
\]
The tuple $(\code{g_a},\code{g_b},\code{g_c},x',y',z')$ is a group configuration over $k'$
(\cref{prop:converse}), with $x'=g_a\ast y'$, $z'=g_b\ast y'=g_c\ast x'$ and
$g_b=g_cg_a$ in $G$, and $\iota(x')=g_a\cdot\iota(y')$, $\iota(z')=g_b\cdot\iota(y')=g_c\cdot\iota(x')$.
\end{theorem}

\begin{proof}
By \cref{prop:stepI,prop:stepII} there is $e_0\ind{}abcxyz$ such that over
$k_0=\acl(k,e_0)$ there is a configuration, which we rename $(a,b,c,x,y,z)$, pointwise
interalgebraic over $k_0$ with the original one, with $b=\code{g_b}$, $g_b\in\mathcal F$,
$z=g_b(y)$, where $\mathcal F$ is a generically transitive generic set of invertible
germs $V\to V'$ ($V=\Loc_{k_0}(y)$, $V'=\Loc_{k_0}(z)$) of rank $n$ satisfying
\ref{F:rank}.  \Cref{thm:hw-bij,thm:hw} give $\mathcal G=\mathcal G(\mathcal F)$, $G$, $S$,
$\iota$ over $k_0$.  We now rename $k_0$ as $k$.

\emph{Step III.1.}  Let $b'$ be a generic point of $\Loc(b)$ over $\acl(abcxyz)$ and
$g_{b'}\in\mathcal F$ the germ with code $b'$.  Since $b\ind{}z$, there is
$\sigma\in\Gamma_{\acl(z)}$ with $\sigma b=b'$; put $y'=\sigma y$, so $z=g_{b'}(y')$.
Let $g=g_{b'}^{-1}g_b\in\mathcal G$ (independent pair).  As in the proof of
\cref{prop:stepII}, $g(y)=y'$ (using $y'\ind{}b'$, a pair of the configuration
$\sigma\mathfrak c$).  Moreover $\Cb(g)\subseteq\acl(bb')$ and
$\Cb(g_b)\subseteq\acl(g_{b'},g)=\acl(b',g)$ (since $g_b=g_{b'}g$), so $b$ and $\code g$ are
interalgebraic over $\acl(k,b')$; and $y'=g_{b'}^{-1}(z)\in\dcl(b',z)$, $z=g_{b'}(y')\in\dcl(b',y')$,
so $z$ and $y'$ are interalgebraic over $\acl(k,b')$.  \textbf{Base change} to
$k_1=\acl(k,b')$ ($b'\ind{}abcxyz$; \cref{lem:gc-invariance}(2) for the configuration and
\cref{lem:hw-basechange} for $\mathcal F,\mathcal G,G,S$, noting $g\ind{}b'$ by
\ref{F:rank}), and \textbf{replace} $b\mapsto\code g$, $z\mapsto y'$
(\cref{lem:gc-invariance}(3)).  The new configuration is $(a,\code g,c,x,y,y')$ with
$g\in\mathcal G$ (generic over $k_1$), $y,y'\in V_{\Gen}$ (both independent from $b'$) and
$y'=g(y)$.

\emph{Step III.2.}  Rename $k_1$ as $k$.  Let $c''$ be a generic point of $\Loc(c)$ over
$\acl(abcxyy')$ (with $b=\code g$ now).  Since $c\ind{}axy$ (\cref{lem:gc-ranks}), there is
$\sigma\in\Gamma_{\acl(axy)}$ with $\sigma c=c''$.  Put $g''=\sigma g\in\mathcal G$ and
$z''=\sigma y'=\sigma(g(y))=g''(y)$ (\cref{lem:eval}(4), $\sigma y=y$).  The image
$\sigma\mathfrak c=(a,\code{g''},c'',x,y,z'')$ is a configuration; its group line gives
$\code{g''}\subseteq\acl(ac'')$ and $a\subseteq\acl(\code{g''}\,c'')$, and its line
$\{c'',z'',x\}$ gives $x\interalg z''$ over $\acl(k,c'')$.  \textbf{Base change} to
$k_2=\acl(k,c'')$ ($c''\ind{}abcxyy'$; \cref{lem:hw-basechange} again; $g$, $y$, $y'$
are independent from $c''$ by the choice of $c''$, and $g''$, $z''$ are independent from
$c''$ as pairs of $\sigma\mathfrak c$, so all of them survive the base change) and
\textbf{replace} $a\mapsto\code{g''}$, $x\mapsto z''$.  The configuration is now
$(\code{g''},\code g,c,z'',y,y')$ over $k_2$ with $g'',g\in\mathcal G$, $z''=g''(y)$,
$y'=g(y)$.

\emph{Step III.3.}  Rename $k_2$ as $k$ and write $g_a=g''$, $g_b=g$, $x'=z''$, $z'=y'$,
$y'=y$ (so $x'=g_a(y')$, $z'=g_b(y')$).  Since $(\code{g_a},\code{g_b},y')$ is a
non-collinear triple of the current configuration, it is independent, and $g_a\ind{}g_b$.
Put $g_c=g_bg_a^{-1}\in G$.  By \cref{lem:gc-ranks}, $x'\ind{}\code{g_a}\code{g_b}$, so
$x'\in\operatorname{Dom}(g_c)$ and $g_c(x')=g_b(g_a^{-1}(x'))=g_b(y')=z'$
(\cref{lem:comp-codes}(2)).  By \cref{lem:triple} applied to the triple
$(\code{g_a}\code{g_b};x',z')$ ($x'$ and $z'$ are independent from $\code{g_a}\code{g_b}$,
and $z'=g_c(x')\in\dcl(\code{g_a}\code{g_b},x')$ by \cref{lem:eval-dcl}), $g_c$ is the unique germ with code in
$\acl(g_ag_b)$ carrying $x'$ to $z'$, so $\Cb(g_c)=\Cb(x'z'/g_ag_b)$.  By \ref{GC:space},
$\acl(cx')=\acl(cz')$, and by \ref{GC:group}, $c\in\acl(g_ag_b)$; hence
$x'z'\ind{c}g_ag_b$ (both sides have rank $m$ over the respective bases, as
$\rk(x'/cg_ag_b)=\rk(x'/g_ag_b)=m$) and trivially $x'z'\ind{g_ag_b}c$; so by \ref{Cb:two}
$\Cb(g_c)=\Cb(x'z'/c)$.  Now \ref{GC:cb} for $c$ says $c\subseteq\acl(\Cb(x'z'/c))=\acl(g_c)$,
and $\Cb(x'z'/c)\subseteq\acl(c)$; so $c\interalg\code{g_c}$.  \textbf{Replace}
$c\mapsto\code{g_c}$.

The resulting configuration is $(\code{g_a},\code{g_b},\code{g_c},x',y',z')$ over the
final base $k'=\acl(k,e)$ with $e=(e_0,b',c'')$, and it is pointwise interalgebraic over
$k'$ with the original configuration (each step replaced points by interalgebraic ones
over a subfield of $k'$).  Here $e\ind{}abcxyz$: $e_0\ind{}abcxyz$; $b'$ was chosen
generic over the closure of $k_0$ and of a configuration interalgebraic over $k_0$ with
$abcxyz$, so $b'\ind{e_0}abcxyz$; likewise $c''\ind{e_0b'}abcxyz$; conclude by
\ref{R:trans}.  Also $g_c\in G$ is in fact in $\mathcal G$ (\cref{thm:hw}\ref{hw:stab},
as $g_b\ind{}g_a^{-1}$).  The family $\mathcal F$ of the statement is the family of
\cref{prop:stepI} restricted along the two base changes (\cref{lem:hw-basechange}:
$\mathcal F''=\{f\in\mathcal F:f\ind{}b'c''\}$, by \ref{R:trans}), which still satisfies
\ref{F:rank} and has $\mathcal G(\mathcal F'')=\mathcal G$; its code variety is
$W=\Loc_{k'}(b_0)$, where $b_0$ denotes the code $b=\code{g_b}$ of \cref{prop:stepI} (not
the $g_b$ of the statement, which is the $g$ of Step III.1); $b_0$ is independent from
$b'c''$ ($b_0\ind{}b'$ by choice of $b'$, and $c''$ is generic over a field containing
$k_1$ and $b_0$), so $W$ is unchanged by the base changes (\cref{lem:dim}(2)).  The identities $x'=g_a(y')$, $z'=g_b(y')=g_c(x')$, $g_b=g_cg_a$
hold by construction, and \cref{thm:hw}\ref{hw:space} translates them into the stated
identities in $S$.  \Cref{prop:converse} confirms that the final tuple is a group
configuration.
\end{proof}

\begin{remark}[parameters used]\label{rem:params}
The tuple $e$ consists of five independent generic copies of group points: three in
\cref{prop:stepI} (copies of $a$, of the enlarged $b$, of $c$) and two in Step III
(copies of $b$ and $c$).  Each is chosen generic over the algebraic closure of
everything constructed before, so $\rk(e)=5n$.  (This counts only the explicitly
adjoined parameters; \cref{rem:finite-trdeg} explains why it does not yield a version
of the theorem for ambient fields of finite transcendence degree.)
\end{remark}

\begin{remark}[faithfulness, connectedness, ranks]\label{rem:main-props}
By \cref{thm:hw}, $G$ acts faithfully and transitively on $S$, the elements of $G$ of
maximal rank $n$ form the single orbit $\mathcal G$ (connectedness), and $V_{\Gen}$ embeds
in $S$ with $G$ acting generically on it in the sense of \ref{hw:action}
(\cref{rem:connected}).  The hypothesis
\ref{GC:cb} was used three times: in \cref{lem:gc-ranks} to see $m\ge1$ (any group
point serves there), for $b$ in \cref{prop:stepI} (canonicity of
$\mathcal F$, which is what makes the germs $g_{b'}$ determined by $b'$ and the rank of
$\mathcal G$ equal to $n$) and for $c$ in Step III.3; the clause for $a$ is used only
through the symmetry of the configuration (the proof could be run with the roles of
$a$ and $c$ exchanged) and is in any case a consequence of the conclusion: the final
tuple satisfies \ref{GC:cb} by \cref{prop:converse}, \cref{lem:gc-invariance}(3)
transfers it to $a$ over $k'$, and \ref{Cb:basechange} with \ref{R:basechange} bring it
back to $k$.  Without
\ref{GC:cb} the theorem is false for the faithful conclusion (\cref{rem:cb-not-acl}; see
also the example of a product group acting through one factor in
\cref{rem:psi-hyp}(b)).
\end{remark}

\section{Constructibility of the group}
\label{sec:constructibility}

The group $G$ of \cref{thm:main} is an abstract group of germs.  This section shows
that it is a \emph{constructible group}: a constructible set modulo a constructible
equivalence relation, with constructible operations, all defined over the base
extended by finitely many fresh generic points.  This is the exact analogue of
``definable group in $K^{\mathrm{eq}}$'' and is the input for any later
algebraization or classification.  Only this section uses Chevalley's theorem.

\subsection{Constructible sets and functions}

\begin{definition}\label{def:constructible}
Let $L\subseteq K$ be a subfield.  A subset $C\subseteq K^N$ is \emph{constructible
over $L$} if it is a finite Boolean combination of zero sets
$Z(f)=\{u: f(u)=0\}$ of polynomials $f\in L[X_1,\dots,X_N]$; equivalently a finite
union of sets $Z(g_1,\dots,g_s)\setminus Z(f)$ with $g_i,f\in L[X]$.  A partial
function $\phi:K^N\rightharpoonup K^M$ is constructible over $L$ if its graph is.
A relation on $K^N$ is constructible if it is a constructible subset of $K^{2N}$.
\end{definition}

\begin{theorem}[Chevalley, $K$-points]\label{thm:chevalley}
If $C\subseteq K^{N+M}$ is constructible over $L$ then its projection
$\pi(C)=\{u\in K^N:\exists v\,(u,v)\in C\}$ is constructible over $L$.
\end{theorem}
\begin{proof}
Let $C_{\mathrm{sp}}\subseteq\operatorname{Spec}L[X,Y]$ be the constructible set given by
the same Boolean combination of the closed sets $V(f)$ (\lean{PrimeSpectrum.zeroLocus});
it is constructible in the topological sense (\lean{Topology.IsConstructible}) because
$\operatorname{Spec}L[X,Y]$ is noetherian, so every open set is retrocompact.
For a $K$-point $(u,v)$ let $\mathfrak p_{(u,v)}=I_L((u,v))\in\operatorname{Spec}L[X,Y]$
(\lean{MvPolynomial.pointToPoint}); then $(u,v)\in C$ iff $\mathfrak p_{(u,v)}\in C_{\mathrm{sp}}$,
since $(u,v)\in Z(f)$ iff $f\in\mathfrak p_{(u,v)}$.  Let
$\pi^*:\operatorname{Spec}L[X,Y]\to\operatorname{Spec}L[X]$ be the contraction map
(\lean{PrimeSpectrum.comap} of the inclusion $L[X]\to L[X,Y]$, a map of finite
presentation).  By Mathlib's Chevalley theorem
\lean{PrimeSpectrum.isConstructible\_comap\_image}, $\pi^*(C_{\mathrm{sp}})$ is
constructible, hence (\lean{PrimeSpectrum.exists\_constructibleSetData\_iff}) a finite
union of sets $V(g_1,\dots,g_s)\setminus V(f)$ with $g_i,f\in L[X]$; let $D\subseteq K^N$
be the $L$-constructible set defined by the same data.  As for $C$, a point $u\in K^N$
lies in $D$ iff $\mathfrak p_u=I_L(u)$ lies in $\pi^*(C_{\mathrm{sp}})$ (again because
$u\in Z(h)$ iff $h\in\mathfrak p_u$ for $h\in L[X]$).  We claim $\pi(C)=D$.

If $u=\pi(u,v)$ with $(u,v)\in C$ then $\pi^*(\mathfrak p_{(u,v)})=\mathfrak p_{(u,v)}\cap L[X]
=\mathfrak p_u$, so $\mathfrak p_u\in\pi^*(C_{\mathrm{sp}})$ and $u\in D$.  Conversely let
$u\in D$, so $\mathfrak p_u=\mathfrak q\cap L[X]$ for some prime $\mathfrak q\in C_{\mathrm{sp}}$,
say $\mathfrak q\in V(g_1,\dots,g_s)\setminus V(f)$ with $g_i,f\in L[X,Y]$.  The domain
$A=L[X,Y]/\mathfrak q$ contains $L[X]/\mathfrak p_u\cong L[u]\subseteq K$ and $\bar f\ne0$.
Localising at $\bar f$ and at the nonzero elements of $L[u]$ (nonzerodivisors in the
domain $A$) gives a nonzero finitely generated $L(u)$-algebra $B$; a maximal ideal of
$B$ has residue field finite over $L(u)$ (Zariski's lemma,
\lean{MvPolynomial.comp\_C\_integral\_of\_surjective\_of\_isJacobsonRing}), which embeds into
the algebraically closed field $K$ over $L(u)$ (\lean{IsAlgClosed.lift}).  The images
$v\in K^M$ of the classes of $Y$ satisfy $g_i(u,v)=0$ and $f(u,v)\ne0$, so $(u,v)\in C$
and $u\in\pi(C)$.
\end{proof}

\begin{corollary}\label{cor:constructible-closure}
Constructible sets and relations over $L$ are closed under finite Boolean operations,
projections, and substitution of constructible functions; the domain, range, image and
preimage of a constructible partial function are constructible; compositions and
inverses (of injective functions) of constructible partial functions are constructible;
and every quasi-rational expression over $L$ defines a constructible partial function
over $L$ (graph $\{(u,w): w_j^{p^e}q_j(u)=p_j(u),\ q_j(u)\ne0\}$).  Universal quantifiers
are handled by complementation.
\end{corollary}

\begin{lemma}[generic points of invariant and of constructible sets]\label{lem:gen-constructible}
Let $V$ be an irreducible $k$-variety and $L$ a small subfield.
\begin{enumerate}[nosep]
\item A $\Gamma_{\acl(L)}$-invariant subset $C\subseteq V(K)$ contains either every generic
  point of $V$ over $\acl(L)$ or none of them.
\item If $C\subseteq V(K)$ is constructible over $L$ then $C$ is $\Gamma_{\acl(L)}$-invariant,
  so (1) applies; and if $C$ contains no generic point of $V$ over $\acl(L)$ then
  $C\subseteq Z(f)$ for some $f\in L[X]\setminus I_V$.
\end{enumerate}
\end{lemma}
\begin{proof}
(1) The generic points of $V$ over $\acl(L)$ form one $\Gamma_{\acl(L)}$-orbit
(\cref{lem:generic-char}(2)).  (2) Automorphisms over $\acl(L)\supseteq L$ fix the
coefficients of the defining polynomials.  For the last claim write
$C=\bigcup_i Z(g_i)\setminus Z(f_i)$ with $g_i$ a tuple of polynomials and $f_i$ a
polynomial, all over $L$ (\cref{def:constructible}).  A generic point $e$ of $V$ over
$\acl(L)$ lies in a closed set $Z(h)$, $h\in L[X]$, iff $h\in I_{\acl(L)}(e)$ iff
$h\in I_V$ (as $P_{\acl(L)}(e)=I_V$).  Hence if $C$ contains no generic point then for
each $i$ either some $g\in g_i$ is not in $I_V$, or all entries of $g_i$ and also $f_i$
lie in $I_V$ (and then $(Z(g_i)\setminus Z(f_i))\cap V(K)=\emptyset$).  Take $f$ to be the
product of one $g\notin I_V$ for each $i$ of the first kind (and $f=1$ if there is none).
\end{proof}

\subsection{The generic quantifier}

\begin{lemma}[generic quantifier]\label{lem:genquant}
Let $V$ be an irreducible $k$-variety of dimension $m\ge1$, $L$ a small subfield, and
$\Phi\subseteq K^N\times V(K)$ constructible over $L$.  Let $t_1,\dots,t_r$ ($r\ge2N+1$) be
generic points of $V$ independent over $L$ (so $\rk(t_1\cdots t_r/L)=rm$).  Then
\[
  \{\,u\in K^N:\ \Phi(u,y)\text{ for every generic point $y$ of $V$ over $\acl(L,u)$}\,\}
  =\{\,u\in K^N:\ \#\{i\le r:\Phi(u,t_i)\}\ge N+1\,\},
\]
and the right-hand side is constructible over $L(t_1,\dots,t_r)$.  The left-hand side
also equals $\{u:\Phi(u,y)\text{ for some generic $y$ over }\acl(L,u)\}$.
\end{lemma}
\begin{proof}
Fix $u$ and let $J=\{i: t_i\not\ind{L}u\}$.  For $i\in J$, $\rk(t_i/Lu)\le m-1$, so
$\rk(t_1\cdots t_r/Lu)\le rm-|J|$; also $\rk(t_1\cdots t_r/Lu)\ge rm-\rk(u/L)\ge rm-N$.
Hence $|J|\le N$.  If $\Phi(u,y)$ holds for all generic $y$ over $\acl(L,u)$, it holds
for all $t_i$ with $i\notin J$ (these are generic over $\acl(L,u)$), i.e.\ for at least
$r-N\ge N+1$ indices.  Conversely, if it holds for $N+1$ indices, one of them is outside
$J$, so $\Phi(u,t_i)$ holds for a generic $t_i$ over $\acl(L,u)$, and then for all
generic points by \cref{lem:gen-constructible}(2) applied to the $L(u)$-constructible set
$\{y:\Phi(u,y)\}$.  The same lemma gives the ``some = every'' equivalence.
\end{proof}

\subsection{Germ equality and admissible parameters}

Fix a quasi-rational expression $F(B,Y)$ over $k$ in variables $B$ (of length $N$) and
$Y$, with values in $K^{n'}$, and irreducible $k$-varieties $V$ (in the variables $Y$) and
$V'$ (in $n'$ variables) of the same dimension $m\ge1$.  For $u\in K^N$ write
$F_u=F(u,\cdot)$.

\begin{definition}[admissible parameters]\label{def:admissible}
$u\in K^N$ is \emph{admissible} if $(\acl(u),F_u)$ is a representative of a germ $V\to V'$
(\cref{def:germ}), and \emph{invertibly admissible} if moreover that germ is invertible.
Write $\mathrm{Adm}$ and $\mathrm{Adm}^\times$ for these sets and $g_u$ for the germ.
\end{definition}

\begin{lemma}\label{lem:adm-constructible}
Let $M$ be any small subfield.  There is a tuple $t_0$ consisting of generic points of
$V$ and of $V'$, independent over $M$, such that $\mathrm{Adm}$, $\mathrm{Adm}^\times$ and the
relation $\{(u,u'):u,u'\in\mathrm{Adm},\ g_u=g_{u'}\}$ are constructible over $k(t_0)$.  More
generally, for finitely many quasi-rational expressions $F^{(1)},\dots,F^{(s)}$ over $k$
(with their own varieties and admissible sets $\mathrm{Adm}_j$, $\mathrm{Adm}^\times_j$) and
any type-correct word $w$ in the symbols $F^{(j)}_{u_j},(F^{(j)}_{u_j})^{-1}$ with source $V$
and target $V''$, the relation
\[
  \bigl\{(u_1,\dots,u_s,u'):\ u_j\in\mathrm{Adm}^\times_j\text{ for all }j,\ u'\in\mathrm{Adm}_{j'},\
  \text{and the composite germ } w(u_1,\dots,u_s)\text{ equals } g_{u'}\bigr\}
\]
is constructible over $k(t_0)$, for a tuple $t_0$ as above (generic points of the finitely
many varieties involved, independent over $M$); the restriction to invertibly admissible
factors is part of the relation.
\end{lemma}
\begin{proof}
$u$ is admissible iff (A) $F_u$ is defined at every generic point of $V$ over $\acl(u)$,
and (B) $F_u(y)$ is a generic point of $V'$ over $\acl(u)$ for every such $y$.
(A) is ``for generic $y$: the denominators of $F_u$ do not vanish at $y$'', constructible
by \cref{lem:genquant}.  For (B) we use: \emph{(B) holds iff for every generic point $w$
of $V'$ over $\acl(u)$ there is $y\in V(K)$ with $F_u(y)=w$.}  Indeed, if (B) holds then
$F_u(y_0)$ is generic in $V'$ over $\acl(u)$ for a generic $y_0$, and every generic point
of $V'$ over $\acl(u)$ is a $\Gamma_{\acl(u)}$-conjugate of it, hence of the form
$F_u(\sigma y_0)$.  Conversely, if $w$ is generic over $\acl(u)$ and $F_u(y)=w$ with
$y\in V(K)$ then $w\in\dcl(u,y)$ gives $\rk(y/u)\ge\rk(w/u)=m$, so $y$ is generic over
$\acl(u)$ (\cref{lem:generic-char}(1)) and $F_u(y)$ is generic; the set of generic $y$
over $\acl(u)$ at which $F_u$ takes a generic value is $\Gamma_{\acl(u)}$-invariant
(\cref{lem:qr-dcl}), so by \cref{lem:gen-constructible}(1) the same then holds for all
generic $y$.  The displayed condition is ``for generic $w\in V'$: $\exists y\in V(K)$,
$F_u(y)=w$'', constructible by \cref{thm:chevalley} and \cref{lem:genquant}.

Invertibility.  Let $u\in\mathrm{Adm}$ (so $F_u$ is defined at every generic point of $V$
over $\acl(u)$ and represents $g_u$).  We claim: $g_u$ is invertible iff for every
generic $w$ of $V'$ over $\acl(u)$ the fibre $\{y\in V(K):F_u(y)=w\}$ is a singleton.
If the fibre over $w=F_u(y)$ is $\{y\}$ then $y$ is fixed by $\Gamma_{k(u,w)}$ (which
maps the fibre to itself, \cref{lem:qr-dcl}), so $y\in\dcl(u,w)\subseteq\dcl(\acl(u),g_u(y))$
by \cref{lem:dcl-galois}, and $g_u$ is invertible by \cref{lem:inverse}(2) with
$F=\acl(u)$.  Conversely let $g_u$ be invertible, $w$ generic over $\acl(u)$,
$y=g_u^{-1}(w)$ (so $F_u(y)=w$), and $y'$ any point of the fibre.  As above $y'$ is
generic over $\acl(u)$, so $(y',w)=(y',F_u(y'))$ is a generic point of $\Gamma(g_u)$ over
$\acl(u)$ (\cref{lem:graph}(1)) and $(w,y')$ a generic point of $\Gamma(g_u^{-1})$
(\cref{lem:inverse}); likewise $(w,y)$.  \Cref{lem:eval}(1) applied to $g_u^{-1}$ gives
$y'=y$.  The fibre condition is constructible by \cref{thm:chevalley} and \cref{lem:genquant}.

Germ equality: for $u,u'\in\mathrm{Adm}$, $g_u=g_{u'}$ iff $F_u(y)=F_{u'}(y)$ for generic
$y$ over $\acl(uu')$ (\cref{lem:eval}(5)), again constructible by \cref{lem:genquant}.

Words.  Let $w$ be a word of length $l$ in the symbols $F^{(j)}_{u_j}$ and
$(F^{(j)}_{u_j})^{-1}$, type-correct with respect to the source and target varieties.
Type-correctness prescribes for each position $i$ of the word an intermediate
variety $V_i$ ($V_0=V$, $V_l=V''$).  By \cref{cor:constructible-closure} the \emph{chain
relation} $R_w(u_1,\dots,u_s;y,y')$ := ``there are $y=y_0,y_1,\dots,y_l=y'$ with
$y_i\in V_i(K)$ for all $i$, such that for each $i$, $(y_{i-1},y_i)$ lies in the graph of
the $i$-th symbol (the graph of $F^{(j)}_{u_j}$, or its transpose for an inverse symbol)''
is constructible over $k$ (each $V_i(K)$ is the zero set of finitely many polynomials
over $k$).  The conjuncts $y_i\in V_i(K)$ are essential: without them an inverse step
would admit points of the ambient affine space outside $V_i$ in the fibre, and the chain
would not be unique.  Suppose every $u_j\in\mathrm{Adm}^\times$, put $u=(u_1,\dots,u_s)$,
and let $y$ be generic over $\acl(u)$.
Then \emph{the chain starting at $y$ is unique and ends at $w(y)$}, where $w$ is also read
as the composite germ (\cref{lem:comp,lem:inverse}): inductively $y_{i-1}$ is generic
over $\acl(u)$; for a forward symbol $y_i=F^{(j)}_{u_j}(y_{i-1})=g_{u_j}(y_{i-1})$
(\cref{lem:eval}(2)), which is generic over $\acl(u)$; for an inverse symbol,
$y_{i-1}=F^{(j)}_{u_j}(y_i)\in\dcl(u,y_i)$ gives $\rk(y_i/u)\ge\rk(y_{i-1}/u)=\dim V_{i-1}
=\dim V_i$ (\cref{lem:inverse}), so $y_i\in V_i(K)$ is generic over $\acl(u)$
(\cref{lem:generic-char}(1)), $(y_i,y_{i-1})$ is a generic point of
$\Gamma(g_{u_j})$ over $\acl(u)$ and $(y_{i-1},y_i)$ one of $\Gamma(g_{u_j}^{-1})$, so
$y_i=g_{u_j}^{-1}(y_{i-1})$ by \cref{lem:eval}(1) applied to $g_{u_j}^{-1}$ over $\acl(u)$.
Consequently ``$w(u_1,\dots,u_s)$ is a germ $V\to V''$ equal to $g_{u'}$'' is equivalent
to: all $u_j\in\mathrm{Adm}^\times$, $u'\in\mathrm{Adm}$, and for generic $y$ of $V$ over
$\acl(uu')$ there is $y'$ with $R_w(u;y,y')$ and $F^{(j')}_{u'}(y)=y'$ (\cref{lem:eval}(5):
agreement at one generic point over $\acl(\Cb w\,\Cb g_{u'})\subseteq\acl(uu')$ suffices,
and the point is generic over the larger field).  This is constructible by
\cref{thm:chevalley} and \cref{lem:genquant}.  In every application of \cref{lem:genquant}
above the relation $\Phi$ is constructible over $k$ itself (its parameters are the
$u$'s, which are quantified, not adjoined), so the generic points $t_1,\dots,t_r$ it needs
are generic points of $V$ or $V'$ independent over $k$; one tuple $t_0$ containing
$r\ge2N'+1$ generic points of each variety involved, with $N'$ the largest parameter
length among all the relations to which \cref{lem:genquant} is applied (including the
longer words used by the consumers of this lemma, e.g.\ the four- and six-symbol words
of \cref{thm:constructible}, whose parameter tuples have length $4N$ and $6N$), serves all
finitely many applications, and by \cref{lem:fresh} it may be chosen independent over
any prescribed small field $M$.
\end{proof}

\subsection{Constructible groups}

\begin{definition}[constructible group, constructible $G$-set]\label{def:constructible-group}
A \emph{constructible group over $L$} is a group $G$ together with a surjection
$\rho:D_G\to G$ from a constructible set $D_G\subseteq K^N$ over $L$ such that the kernel
relation $\{(d,d'):\rho d=\rho d'\}$ and the relations
$\{(d_1,d_2,d_3):\rho d_1\cdot\rho d_2=\rho d_3\}$ and $\{(d,d'):(\rho d)^{-1}=\rho d'\}$ are
constructible over $L$.  A \emph{constructible $G$-set} is a $G$-set $X$ with a
surjection $\rho_X:D_X\to X$ from a constructible set such that the kernel relation and
$\{(d,d_X,d'_X):\rho d\cdot\rho_Xd_X=\rho_Xd'_X\}$ are constructible.  A subgroup $H\le G$
is constructible if $\rho^{-1}(H)$ is.
\end{definition}

\begin{theorem}[the group of \cref{thm:main} is constructible]\label{thm:constructible}
In the situation of \cref{thm:main} (base $k'=\acl(k,e)$; varieties $V,V'$; family
$\mathcal F$ of invertible germs $V\to V'$ with code variety $W$ and
$\mathcal G=\mathcal G(\mathcal F)$), let $F(B,Y)$ be the quasi-rational expression over
$k'$ with $F(\code f,y)=f(y)$ for $f\in\mathcal F$ (\cref{rem:code-param}), and let
$y_0\in V_{\Gen}$ be generic over $\acl(k',abcxyz)$.  There is a tuple $t=(t_0,t_1)$ of generic points of the varieties $V$, $V'$ (in $t_0$) and
$W$ (in $t_1$), independent over $\acl(k',abcxyz,y_0)$ (\cref{lem:fresh}; in particular
$(e,t)\ind{}abcxyz$ over $k$), chosen in two stages as described in the proof, such that,
over $k''=k'(t)$:
\begin{enumerate}[nosep]
\item $G=\{f_1^{-1}f_2: f_1,f_2\in\mathcal F\}$, and $G$ is a constructible group over $k''$
  with $D_G=\{(u_1,u_2)\in\mathrm{Adm}^\times\times\mathrm{Adm}^\times:
  (u_1,u_2)\text{ satisfies }(\ast)\}$, $\rho(u_1,u_2)=g_{u_1}^{-1}g_{u_2}$, where
  $(\ast)$ is the condition: \emph{for every $u$ generic in $W$ over $\acl(k'',u_1,u_2)$
  there is $u'\in W(K)\cap\mathrm{Adm}^\times$ with $g_u\,g_{u_1}^{-1}g_{u_2}=g_{u'}$};
\item the stabiliser $H=\{\gamma\in G:\gamma\cdot\iota(y_0)=\iota(y_0)\}$ is a constructible
  subgroup over $k''(y_0)$, and the orbit map $G/H\to S$, $\gamma H\mapsto\gamma\cdot\iota(y_0)$,
  is a $G$-equivariant bijection; hence $S$ is a constructible $G$-set over $k''(y_0)$.
\end{enumerate}
\end{theorem}
\begin{proof}
(1) \emph{Two-step description.}  Every $f_1^{-1}f_2$ lies in $G$: for $f_1,f_2\in\mathcal F$
choose $f\in\mathcal F$ with $f\ind{}f_1f_2$; then $f_1^{-1}f_2=(f_1^{-1}f)(f^{-1}f_2)$ with
both factors in $\mathcal G(\mathcal F)$.  Conversely, for $\gamma\in G$ choose $f\in\mathcal F$
with $f\ind{}\gamma$; by \cref{thm:hw-bij}, $f\gamma\in\mathcal F$ and $\gamma=f^{-1}(f\gamma)$.

\emph{The condition $(\ast)$ characterises $G$ inside the germs $g_{u_1}^{-1}g_{u_2}$
($u_i\in\mathrm{Adm}^\times$).}  If $\gamma=g_{u_1}^{-1}g_{u_2}\in G$ then for $u$ generic in
$W$ over $\acl(k'',u_1,u_2)\supseteq\acl(\gamma)$, $g_u\in\mathcal F$ and $g_u\gamma\in\mathcal F$
(\cref{thm:hw-bij}, as $g_u\ind{}\gamma$), i.e.\ $g_u\gamma=g_{u'}$ with
$u'=\code{g_u\gamma}$, a generic point of $W$, so $(\ast)$ holds.
Conversely suppose $(\ast)$ holds for $\gamma=g_{u_1}^{-1}g_{u_2}$ (a germ in
$\Germ^\times(V)$ by \cref{lem:inverse,lem:comp}).  Pick $u$ generic in $W$ over
$\acl(k'',u_1,u_2)$ and $u'\in W(K)\cap\mathrm{Adm}^\times$ with $g_u\gamma=g_{u'}$.  Then
$\rk(g_{u'})\ge\rk(g_u\gamma/\gamma)=\rk(g_u/\gamma)=n$, so $\Cb(g_{u'})\subseteq\acl(u')$ forces
$\rk(u')\ge n$, i.e.\ $u'$ is a generic point of $W$ and $g_{u'}\in\mathcal F$; hence
$\gamma=g_u^{-1}g_{u'}\in G$ by the two-step description.

\emph{Constructibility; the staged choice of $t$.}  First choose $t_0$ as in
\cref{lem:adm-constructible} (generic points of $V$ and $V'$) independent over
$M_0=\acl(k',abcxyz,y_0)$, and put $L_0=k'(t_0)$.  Then $\mathrm{Adm}^\times$, the relation
$g_{u_1}^{-1}g_{u_2}=g_{u'_1}^{-1}g_{u'_2}$ (kernel of $\rho$), the relations
$\rho(d_1)\rho(d_2)=\rho(d_3)$ and $\rho(d)^{-1}=\rho(d')$, and the relation
$\Psi(u_1,u_2;u):=$ ``$\exists u'\in W(K)\cap\mathrm{Adm}^\times$ with
$g_ug_{u_1}^{-1}g_{u_2}=g_{u'}$'' (all word equalities, then a projection; note $W(K)$ is
constructible over $k'$) are constructible over $L_0$.  Next choose $t_1$, consisting of
$r\ge2N'+1$ generic points of $W$ independent over $\acl(M_0,t_0)$ ($N'$ the largest
parameter length of the relations quantified below, i.e.\ $6N$ for the action relation of
$G/H$), and put
$k''=k'(t_0,t_1)$.  By \cref{lem:genquant} (with $L=L_0$, $\Phi=\Psi$) the set of
$(u_1,u_2)$ such that $\Psi(u_1,u_2;u)$ holds for every $u$ generic in $W$ over
$\acl(L_0,u_1,u_2)$ is constructible over $L_0(t_1)=k''$; and by
\cref{lem:gen-constructible}(2) applied to the $L_0(u_1,u_2)$-constructible set
$\{u:\Psi(u_1,u_2;u)\}$, quantifying over the points generic over $\acl(L_0,u_1,u_2)$ is the
same as quantifying over those generic over the larger field $\acl(k'',u_1,u_2)$ (all or
none of either kind satisfy $\Psi$), which is the form $(\ast)$ in the statement.  Hence
$D_G$ and the structure relations of $G$ are constructible over $k''$.

(2) Let $\gamma=\rho(d)\in G$, $d=(u_1,u_2)\in D_G$.  We claim: \emph{$\gamma\in H$ iff for
every $v$ generic in $W$ over $\acl(k'',d,y_0)$ there is $u'\in W(K)\cap\mathrm{Adm}^\times$
with $g_v\gamma=g_{u'}$ and $F(u',y_0)=F(v,y_0)$ (both sides defined).}  Granting the claim,
$H$ is constructible over $k''(y_0)$: the relation in $(d,v,u')$ is constructible over
$L_0(y_0)$ (``$g_vg_{u_1}^{-1}g_{u_2}=g_{u'}$'' is a word relation, \cref{lem:adm-constructible};
the two values of $F$ are quasi-rational in $(u',y_0)$ and $(v,y_0)$,
\cref{cor:constructible-closure}), its projection $\Phi_H(d,v)$ is constructible over
$L_0(y_0)$ (\cref{thm:chevalley}), and the generic quantifier over $v\in W$ is
\cref{lem:genquant} with $L=L_0(y_0)$ and the same tuple $t_1$ (its points are generic over
$\acl(M_0,t_0)\supseteq L_0(y_0)$), giving a set constructible over $L_0(y_0,t_1)=k''(y_0)$;
as in (1), \cref{lem:gen-constructible}(2) identifies ``for every $v$ generic over
$\acl(L_0,y_0,d)$'' with ``for every $v$ generic over $\acl(k'',d,y_0)$''.  Note that the claim never evaluates a parametrised expression at $y_0$
with parameters that may depend on $y_0$: it only evaluates single members $F(u',\cdot)$,
$F(v,\cdot)$ of the family, and the rank argument below shows that $y_0$ is automatically
generic over $\acl(u')$ for every witness $u'$.

\emph{Proof of the claim.}  Let $v$ be generic in $W$ over $\acl(k'',d,y_0)$, so $v\ind{}dy_0$
and $f:=g_v\in\mathcal F$ satisfies $f\ind{}\gamma y_0$ (as $\Cb(\gamma)\subseteq\acl(d)$).  By
\cref{thm:hw-bij} (last claim), $f\gamma\in\mathcal F$; let $w=\code{f\gamma}$, a generic point
of $W$.  Since $f\gamma$ and $f$ are interalgebraic over $\gamma$ (\cref{lem:comp-codes}(1)),
$\rk(w/k''dy_0)=\rk(v/k''dy_0)=n$ (ranks here and below are over $k'$; $v$ is generic over
$\acl(k'',d,y_0)$), so $w\ind{}k''dy_0$, in particular $w\ind{}y_0$.  Choose $f_1\in\mathcal F$ generic
over $\acl(k'',v,d,y_0)$ and put $h=f_1^{-1}f\in\mathcal G$; then $h\ind{}\gamma y_0$
($\Cb(h)\subseteq\acl(f_1,f)$, $f\ind{}\gamma y_0$ and $f_1\ind{}f\gamma y_0$, \ref{R:swap}), so
$h$ is an admissible witness in the definition of $E$: $\gamma\in H$ iff
$(h\gamma)(y_0)=h(y_0)$.  Now $h\gamma=f_1^{-1}(f\gamma)=f_1^{-1}g_w$ and $y_0$ is generic over
$\acl(f_1,w)$: $\rk(f_1w/y_0)=\rk(w/y_0)+\rk(f_1/wy_0)\ge n+\rk(f_1/vdy_0)=2n$ (using
$w\subseteq\acl(v,d)$).  Hence $(h\gamma)(y_0)=f_1^{-1}(g_w(y_0))$ and $h(y_0)=f_1^{-1}(f(y_0))$
(\cref{lem:comp}; $y_0$ is generic over $\acl(f_1,v)$ likewise).  Both $g_w(y_0)$ and
$f(y_0)$ are generic points of $V'$ over $\acl(f_1)$ (each is interalgebraic with $y_0$
over its code, and $y_0\ind{}f_1w$, $y_0\ind{}f_1v$), and $f_1^{-1}$ is injective on such
points (if $f_1^{-1}(z)=f_1^{-1}(z')$ then $(y,z)$ and $(y,z')$ are generic points of
$\Gamma(f_1)$ over $\acl(f_1)$ for $y=f_1^{-1}(z)$, so $z=z'$ by \cref{lem:eval}(1)).  Therefore
$\gamma\in H$ iff $g_w(y_0)=g_v(y_0)$, i.e.\ iff $F(w,y_0)=F(v,y_0)$ (\cref{lem:eval}(2),
as $y_0$ is generic over $\acl(w)$ and over $\acl(v)$).

If $\gamma\in H$, then $u':=w$ is a witness.  Conversely let $u'\in W(K)\cap\mathrm{Adm}^\times$
with $g_{u'}=g_v\gamma=g_w$ and $F(u',y_0)=F(v,y_0)$.  Then $k'(w)=\Cb(g_{u'})\subseteq\acl(u')$,
so $\rk(u'/k''dy_0)\ge\rk(w/k''dy_0)=n\ge\rk(u')$ (as $u'\in W(K)$ and $\dim W=n$,
\cref{lem:generic-char}(1)); hence $u'$ is generic in $W$ over $\acl(k'',d,y_0)$ and $y_0$ is
generic over $\acl(u')$.  Since $(\acl(u'),F_{u'})$ represents $g_{u'}=g_w$,
$F(u',y_0)=g_w(y_0)$ (\cref{lem:eval}(2)), and $F(v,y_0)=g_v(y_0)$; so $g_w(y_0)=g_v(y_0)$ and
$\gamma\in H$ by the equivalence just proved.  The orbit map is $G$-equivariant and injective by
definition of $H$, and surjective because $G$ acts transitively on $S$
(\cref{thm:hw}\ref{hw:space}).  The quotient $G/H$ is presented by $D_G$ with the kernel
relation $\{(d,d'):\rho(d)^{-1}\rho(d')\in H\}$, constructible by (1) and the claim (applied
to the word $g_{u_2}^{-1}g_{u_1}g_{u'_1}^{-1}g_{u'_2}$ in place of $g_{u_1}^{-1}g_{u_2}$), and the
action is multiplication.
\end{proof}

\begin{remark}[connectedness and rank as properties of the constructible group]
\label{rem:constr-connected}
The generic elements of $G$ are the elements of rank $n$ over $\acl(k'')$ (these are
the elements of $\mathcal G(\mathcal F)$ independent from $t$, \cref{lem:hw-basechange}), they
form one $\Gamma_{(k'')}$-orbit, and every element of $G$ is a product of two of them
(\cref{thm:hw}).  In the language of definable groups: $G$ has Morley rank $n$ and degree
one.  We do not need, and do not prove, the ``definability of rank and degree in
families'' that would express these as constructible conditions on parameters; see
\cref{rem:descent-limits}.
\end{remark}

\begin{leannote}
Mathlib's \lean{PrimeSpectrum.isConstructible\_comap\_image} is stated for
\lean{RingHom.FinitePresentation}; the inclusion \lean{MvPolynomial (Fin N) L → MvPolynomial (Fin N ⊕ Fin M) L}
is of finite presentation: identify \lean{MvPolynomial (Fin N ⊕ Fin M) L} with
\lean{MvPolynomial (Fin M) (MvPolynomial (Fin N) L)} via \lean{MvPolynomial.sumAlgEquiv},
apply \lean{Algebra.FinitePresentation.mvPolynomial\_of\_finitePresentation} (with
\lean{Algebra.FinitePresentation.self}), and pass to \lean{RingHom.FinitePresentation}
through \lean{RingHom.finitePresentation\_algebraMap}.  The
translation between $L$-constructible subsets of $K^N$ and
\lean{PrimeSpectrum.ConstructibleSetData (MvPolynomial (Fin N) L)} is the content of
\cref{thm:chevalley}; it is new (not in Mathlib) but elementary given
\lean{MvPolynomial.pointToPoint}, \lean{IsAlgClosed.lift} and Zariski's lemma.
\end{leannote}

\section{Descent of independent parameters}
\label{sec:descent}

The main theorem adjoins to the base a tuple $e$ independent from the configuration.
First-order consequences of the theorem can be brought back to the original base $k$
by the following lemma, which is the universal-domain form of the ``place''
argument of Evans--Hrushovski (1991, Lemma 2.1.1) and of the coheir argument used by
Bays--Zou.

\begin{lemma}[density of $k$-points]\label{lem:kpoints-dense}
Let $U$ be an irreducible $k$-variety and $f\in K[X]\setminus I_U$.  Then there is
$u_0\in U(k)=U(K)\cap k^N$ with $f(u_0)\ne0$.
\end{lemma}
\begin{proof}
Choose a $k$-basis $(\beta_j)$ of $K$ and write $f=\sum_j f_j\beta_j$ with $f_j\in k[X]$;
since $I_U$ is defined over $k$, $f\notin I_U$ means some $f_j\notin I_U\cap k[X]=:J$,
a prime ideal of $k[X]$.  If $f_j$ vanished on $U(k)=Z_k(J)$ then, by the Nullstellensatz
over the algebraically closed field $k$ (\lean{MvPolynomial.vanishingIdeal\_zeroLocus\_eq\_radical}
with both fields equal to $k$), $f_j\in\sqrt J=J$, a contradiction.  So some
$u_0\in U(k)$ has $f_j(u_0)\neq0$; as the $\beta_j$ are linearly independent over $k$,
$f(u_0)\ne0$.
\end{proof}

\begin{theorem}[descent]\label{thm:descent}
Let $\Phi\subseteq K^{N}\times K^{M}$ be constructible over $k$, let $w\in K^M$ and let
$e\in K^N$ with $e\ind{}w$ and $\Phi(e,w)$.  Then there is $e_0\in k^N$ with $\Phi(e_0,w)$.
\end{theorem}
\begin{proof}
Let $U=\Loc(e)$ and $C=\{u\in U(K):\Phi(u,w)\}$, a set constructible over $k(w)$
containing $e$.  Since $e\ind{}w$, $e$ is a generic point of $U$ over $\acl(w)$
(\cref{lem:dim}(2)).  By \cref{lem:gen-constructible} (with $L=k(w)$), $C$ contains
every generic point of $U$ over $\acl(w)$, so its complement $U(K)\setminus C$, which is
constructible over $k(w)$ and contains no generic point, is contained in $Z(f)$ for some
$f\in k(w)[X]\setminus I_U$.  By \cref{lem:kpoints-dense} there is $u_0\in U(k)$ with
$f(u_0)\ne0$; then $u_0\in C$, i.e.\ $\Phi(u_0,w)$.
\end{proof}

\begin{remark}[how to use descent; its limits]\label{rem:descent-limits}
Any consequence of \cref{thm:main} written as ``$\Phi(e,abcxyz)$'' for a
constructible $\Phi$ over $k$ (with $e$ the adjoined parameters, and all further
ambient-field witnesses existentially quantified) therefore holds with parameters
in $k$. To use this for interalgebraicity or interdefinability, fix the actual
finite polynomial or quasi-rational witness formulas supplied by the proof,
including their coefficients in $k$ and their nonvanishing conditions. A fixed
polynomial equation and a fixed polynomial inequation define constructible sets;
finite Boolean combinations preserve constructibility, as does existential
quantification over finite tuples of $K$ when $K$ is algebraically closed. For
example, a chosen relation
$P(e,v,u)=0$, together with the nonvanishing of a chosen coefficient of $P$ as a
polynomial in $u$, gives a constructible algebraicity witness. A chosen relation
$Q(e,v)u^{p^d}=R(e,v)$ with $Q(e,v)\ne0$ gives a constructible definability witness
in positive characteristic; use exponent $1$ in characteristic zero.

A degree bound alone does not make the full relations ``$u\in\acl(k,e,v)$'' or
``$u\in\dcl(k,e,v)$'' constructible in the ambient tuples. Already with empty
$(e,v)$, the roots of nonzero polynomials over $k$ of degree at most one form
exactly $k$. Constant rational-function witnesses also give exactly $k$. This
proper infinite subfield is not a constructible subset of the line: constructible
subsets are finite or cofinite, while the complement of $k$ contains $t+k$ for
any $t\notin k$. Quantifying coefficients over $K$ does not restrict them to $k$.
The earlier wording ``once a degree bound is fixed'' is refuted in issue \#31;
its original provenance is \texttt{6b76e01}, this remark. The Lean statements
\lean{GroupConfig.not\_isConstructibleOver\_range\_algebraMap} and
\lean{GroupConfig.not\_isConstructibleOver\_setOf\_exists\_natDegree\_le\_one}
record the base-field obstruction and its degree-one algebraicity case. The
corrected fixed-formula statement is given by
\lean{GroupConfig.isConstructibleOver\_iff\_exists\_polynomial\_data},
\lean{GroupConfig.IsConstructibleOver.polynomial\_zero},
\lean{GroupConfig.IsConstructibleOver.polynomial\_ne\_zero},
\lean{GroupConfig.QuasiRational.isConstructibleOver\_graph}, the finite Boolean
operations and \lean{GroupConfig.IsConstructibleOver.exists}.

The assertions ``$\rk(u/e)=r$'', ``$G$ is connected'' and ``$\rk(G)=n$'' are not
first-order in the parameter as stated; definability of Morley rank and degree
in $\mathrm{ACF}$ treats families, but we do not prove that here. The application
in \cref{sec:application} still requires fixing formulas for the
interalgebraicity and interdefinability conclusions to be descended, and proving
that the needed independence of the descended witnesses follows from the rank
equalities of the original configuration. These obligations are not discharged
by a degree bound or by the constructible-relation descent theorem alone.
\end{remark}

\section{Application: the $\Psi$-witness of the reconstruction blueprint}
\label{sec:application}

\subsection{The $\Psi$-witness is a group configuration}

Recall the setting of the reconstruction blueprint (\texttt{sources/blueprint.tex},
\S``The geometric definition of $Q$''; Lean: \lean{AclGeom.QWitness.Psi}).  There
$k\subseteq K$ are fields with $K$ an arbitrary extension, $\mathcal G(K/k)$ is the lattice
of relatively algebraically closed intermediate fields (\lean{AclGeom.ClosedIF}), points
are its atoms (\lean{AclGeom.Point}), and a witness $w$ consists of points
$A_1,A_2,B_1,B_2,C_1,C_2,D,\dots,Z$ with $A=A_1\vee A_2$ etc.  The clauses of $\Psi$ used
here are:
\begin{enumerate}[label=(\roman*),nosep]
\item $\rk(A\vee B\vee C)=\rk(A\vee B)=\rk(B\vee C)=\rk(A\vee C)=4$;
\item $X\le A\vee Y$ and $Z\le(B\vee Y)\wedge(C\vee X)$;
\item $X,Y,Z\not\le A\vee B\vee C$;
\item for every atom $A'\le A$: $X\not\le A'\vee Y$; for every atom $B'\le B$:
  $Z\not\le B'\vee Y$; for every atom $C'\le C$: $Z\not\le C'\vee X$.
\end{enumerate}
We assume in addition that $k$ and $K$ are algebraically closed and that
$\trdeg(K/k)$ is infinite (\cref{conv:ud}); see \cref{rem:psi-hyp} for the general case.

\begin{proposition}[$\Psi$ gives a group configuration]\label{prop:psi-gc}
Let $w$ satisfy (i)--(iv) and choose representatives: $a=(A_1.\mathrm{rep},A_2.\mathrm{rep})$,
$b=(B_1.\mathrm{rep},B_2.\mathrm{rep})$, $c=(C_1.\mathrm{rep},C_2.\mathrm{rep})$,
$x=X.\mathrm{rep}$, $y=Y.\mathrm{rep}$, $z=Z.\mathrm{rep}$ (so that $\acl(a)=A$, etc.).  Then
$(a,b,c,x,y,z)$ is a group configuration over $k$ with $n=2$, $m=1$, in the orientation
of \cref{conv:orientation}: $X=A\ast Y$, $Z=B\ast Y$, $Z=C\ast X$.
\end{proposition}
\begin{proof}
Throughout, joins of closed subfields are relative algebraic closures of unions
(\lean{AclGeom.ClosedIF.coe\_sup}), and ranks of closed subfields are transcendence
degrees, so for tuples $\rk(pq)=\rk(\acl(p)\vee\acl(q))$.  We use the submodularity
$\rk(E\wedge F)+\rk(E\vee F)\le\rk(E)+\rk(F)$ of the rank of closed subfields, which is
Mathlib's \lean{Matroid.eRk\_inter\_add\_eRk\_union\_le} for
\lean{AlgebraicIndependent.matroid k K} transported through \lean{AclGeom.rankEq\_iff\_eRk}
and \lean{AclGeom.ClosedIF.coe\_sup} (the project's \lean{AclGeom/Config/PsiLattice.lean}
contains this derivation as a private lemma), and the exchange property for points,
\lean{AclGeom.pointCl\_exchange} / \lean{AclGeom.racl\_exchange}.

\ref{GC:group}: by (i), $A\vee B\le A\vee B\vee C$ have equal rank $4$, so they are equal;
likewise $A\vee C=B\vee C=A\vee B\vee C$.  Hence $\acl(ab)=\acl(ac)=\acl(bc)$.

\ref{GC:space}: $X\le A\vee Y$ by (ii), and $X\not\le A$ by (iii), so by exchange in the
lattice (if $P\le E\vee Q$ and $P\not\le E$ then $Q\le E\vee P$, for points $P,Q$ and closed
$E$ of finite rank; \lean{AclGeom.pointCl\_exchange}) $Y\le A\vee X$; thus
$\acl(ax)=A\vee X=A\vee Y=\acl(ay)$.  Similarly $Z\le B\vee Y$, $Z\not\le B$ give
$Y\le B\vee Z$, and $Z\le C\vee X$, $Z\not\le C$ give $X\le C\vee Z$.

Distinctness: $X\ne Y$, since $X=Y$ would give $X\le A'\vee Y$ for any atom $A'\le A$,
contradicting (iv) (a rank-two closed field has atoms).  Likewise $Z\ne Y$ (from the
$B$-clause) and $Z\ne X$ (from the $C$-clause).

\ref{GC:indep}: ranks are $\rk(a)=\rk(b)=\rk(c)=2$ (from (i): $\rk(A)+\rk(B)\ge4$ and each
$\rk\le2$) and $\rk(x)=\rk(y)=\rk(z)=1$.  We verify the $16$ non-collinear triples.
\emph{Two group points and a space point} ($9$ triples): $A\vee B=A\vee C=B\vee C$ has
rank $4$ and does not contain $X,Y,Z$ by (iii), so each of these joins has rank $5$.
\emph{One group point and two space points not on its line} ($6$ triples):
$(a,y,z)$: $Z\le A\vee Y$ would give $Z\le(A\vee Y)\wedge(B\vee Y)$, and
$\rk((A\vee Y)\wedge(B\vee Y))\le3+3-\rk(A\vee B\vee Y)=1$ with $Y$ below it, so the meet
is $Y$ and $Z=Y$, a contradiction; hence $\rk(ayz)=4$.  $(a,x,z)$: $A\vee X=A\vee Y$, same.
$(b,x,y)$: $X\le B\vee Y$ would give $X\le(A\vee Y)\wedge(B\vee Y)=Y$; so $\rk(bxy)=4$;
$(b,x,z)$: $B\vee Z=B\vee Y$.  $(c,x,y)$: $Y\le C\vee X$ would give
$Y\le(A\vee X)\wedge(C\vee X)=X$ (rank $\le3+3-5=1$); $(c,y,z)$: $C\vee Z=C\vee X$.
\emph{The three space points}: $X\ne Y$ gives $\rk(xy)=2$; if $Z\le X\vee Y$ then
$Z\le(X\vee Y)\wedge(B\vee Y)$, whose rank is $\le2+3-\rk(B\vee X\vee Y)=2+3-4=1$ (using
$\rk(bxy)=4$), so the meet is $Y$ and $Z=Y$; hence $\rk(xyz)=3$.

\ref{GC:cb}: Let $E_a=\acl(\Cb(xy/a))\subseteq A$, a closed subfield of rank $\le2$.  By
\ref{Cb:ind} and \ref{R:acl}, $xy\ind{E_a}A$, so $\rk(x/E_ay)=\rk(x/Ay)=0$, i.e.\
$X\le E_a\vee Y$.  If $\rk(E_a)=0$ then $E_a=k$ and $X\le Y$, so $X=Y$, which was excluded;
if $\rk(E_a)=1$ then $E_a$ is an atom $A'\le A$ with $X\le A'\vee Y$, contradicting (iv).
Hence $\rk(E_a)=2$, $E_a=A$ and $a\subseteq\acl(\Cb(xy/a))$.  The same argument with the
$B$- and $C$-clauses of (iv) gives the conditions for $b$ (germ $Y\to Z$, closed field
$E_b$) and $c$ (germ $X\to Z$, closed field $E_c$).  (The Lean theorem
\lean{AclGeom.QWitness.Psi.exists\_incidence\_germs\_with\_full\_coefficients}, commit
\lean{48dd789}, is this argument carried out for the canonical plane-curve equations of
the three germs over the fields generated by the displayed generator pairs; the theorem
\lean{AclGeom.QWitness.Psi.exists\_composing\_closed\_incidence\_germs}, commit
\lean{8d01b03}, assuming \lean{IsAlgClosed K}, produces the three germs over the single
closed rank-four field $A\vee B$ with coefficient closures $A$, $B$, $C$ respectively,
together with their composition.  Either can serve as the formal counterpart once $\Cb$
is identified with the coefficient field of the canonical equation, see the Lean note
after \cref{prop:cb}.)

\ref{GC:rank}: $n=2$.
\end{proof}

\begin{remark}[on the hypotheses]\label{rem:psi-hyp}
(a) The clauses (v)--(vii) of $\Psi$ (the partial quadrangle $(S,T,U)$ and the seven
meets) are not used; they enter only the subsequent chart extraction.
(b) Clause (iv) is exactly \ref{GC:cb} for $n=2$: a proper subfield of definition of a
rank-two plane is contained in an atom.  Without (iv) the theorem's faithful conclusion
fails: for independent $s,s',t,u,u'$ let $a=(s,u)$, $b=(s',u')$, $c=(s'/s,u'/u)$,
$y=t$, $x=st$, $z=s't$ (the group $\mathbb G_m\times\mathbb G_m$ acting on the line
through its first factor, with $b=c\ast a$ and $z=c\ast x=(s'/s)(st)$).  Then (i)--(iii)
hold, but (iv) fails for the atom $A'=\acl(s)\le A$ since $X=\acl(st)\le A'\vee Y$; and no
faithful rank-two action on a curve can realise this tuple, because the conclusion of
\cref{thm:main} forces \ref{GC:cb} for the original tuple (\cref{prop:converse} for the
configuration built from $(G,S)$, transferred to the interalgebraic tuple over $k$ by
\cref{lem:gc-invariance}(3), \ref{Cb:basechange} and \ref{R:basechange}; see
\cref{rem:main-props}).
(c) If $K$ is not algebraically closed or has finite transcendence degree, pass to an
algebraically closed extension $K'\supseteq K$ of infinite transcendence degree
(\lean{AlgebraicClosure} of a rational function field in countably many variables over
$K$).  All clauses of $\Psi$ are preserved in $\mathcal G(K'/k)$ for the same tuples (the
closed subfields are replaced by their closures in $K'$; ranks and the relations
$\le,\wedge,\vee$ among closures of finitely generated subfields are unchanged; in Lean
this is \lean{AclGeom.QWitness.Psi.lift} and \lean{AclGeom.QGeom.lift} in
\lean{AclGeom/Transfer/LiftConfig.lean}), and the conclusions below are then about tuples
of $K'$.  Returning the existential witnesses produced in $K'$ to $K$ is a separate
descent step that this document does not prove.  For $K$ algebraically closed it is
model completeness of the theory of algebraically closed fields: a conclusion that is
first-order in the tuples of $K$ (with fixed finite polynomial or quasi-rational
witness formulas, their coefficients and nonvanishing conditions retained, and all
further ambient-field witnesses existentially quantified, as in
\cref{rem:descent-limits}) holds in $K$ as soon as it
holds in $K'$.  For an arbitrary field $K$ there is no such automatic descent (the
reconstruction blueprint's arbitrary-field counterexamples, issue \#25, show that
first-order consequences of $\Psi$ can fail in $K$ while holding in $K'$), so the
present results are results about algebraically closed ambient fields only.  The
library theorem \lean{AclGeom.QWitness.Psi.hasAffineGridCoordinates\_of\_chart}
(\lean{AclGeom/Config/AffineGridNormalization.lean}) is \emph{not} a general
witness-descent statement: it assumes that the initial independent points and the
initial chart rows already lie in $K$ and only descends the remaining rows and
base-field shifts.  \Cref{sec:descent} is the other half of the problem: it brings the
adjoined parameters $e,t$ into $k$ for first-order consequences with independent
parameters.  If $k$ is not algebraically closed, replace it by its
relative algebraic closure $\acl_k(\emptyset)$ in $K$: for $K$ algebraically closed this
changes neither the closed subfields, nor the points, nor the ranks of $\mathcal G(K/k)$, so
the assumption that $k$ is algebraically closed costs nothing.
\end{remark}

\subsection{What the theorem gives}

\begin{corollary}[group and curve from $\Psi$]\label{cor:psi}
Let $w$ satisfy (i)--(iv) as above, with $k,K$ algebraically closed and $\trdeg(K/k)$
infinite.  Then there is a tuple $e\in K^{<\omega}$ with $e\ind{}abcxyz$ such that over
$k'=\acl(k,e)$:
\begin{itemize}[nosep]
\item there are an irreducible $k'$-curve $V$ ($\dim V=1$) and a group $G$ of germs of
  rank $2$, connected in the sense of \cref{rem:constr-connected}, acting faithfully and
  transitively on a set $S\supseteq V_{\Gen}$; for a generic point $y_0$ of $V$ over
  $\acl(k',abcxyz)$ and a tuple $t$ of fresh generic points independent over
  $\acl(k',abcxyz,y_0)$ (so $(e,t)\ind{}abcxyz$), $G$ is a constructible group over
  $k'(t)$ and $S\cong G/H$ with $H$ the stabiliser of $\iota(y_0)$, a constructible
  subgroup over $k'(t,y_0)$ (\cref{thm:constructible}); and
\item generic elements $g_a,g_b\in G$ and a generic point $y'\in V(K)$, independent over
  $k'$, with
  \[
    A\cdot k'\interalg\code{g_a},\quad B\cdot k'\interalg\code{g_b},\quad
    C\cdot k'\interalg\code{g_bg_a^{-1}},\quad X\cdot k'\interalg g_a(y'),\quad
    Y\cdot k'\interalg y',\quad Z\cdot k'\interalg g_b(y'),
  \]
  where $E\cdot k'$ denotes $\acl(k'\cup E)$ and $\interalg$ is interalgebraicity over $k'$.
\end{itemize}
\end{corollary}
\begin{proof}
\Cref{prop:psi-gc}, \cref{thm:main}, \cref{thm:constructible}.
\end{proof}

\begin{remark}[the remaining steps towards the affine chart]\label{rem:remaining}
To obtain the coordinates of blueprint Lemma \texttt{affine-grid-extraction} one still
has to (1) classify the pair $(G,S)$: a connected constructible group of rank $2$ acting
faithfully and transitively on a set of rank $1$ containing the generic points of a curve
is, up to definable isomorphism, the affine group $\{t\mapsto\alpha t+\beta\}$ acting on
the affine line (blueprint Lemma \texttt{affine-action}; Evans--Hrushovski 1995 cite
Cherlin's classification of groups of Morley rank $2$ and Poizat for the identification
of definable groups with algebraic groups; a direct proof will need either Weil's
theorem followed by the classification of two-dimensional algebraic groups acting on
curves, or a self-contained argument for constructible groups of rank two); (2) read off
coordinates $(\alpha,\beta)$ for $g_a$ and $g_b$ and a coordinate on the line;
(3) identify the partial-quadrangle points $S,T,U$ with the multiplier coordinates
(the proof of blueprint Lemma \texttt{affine-grid-extraction}, equation (8.7)); and
(4) descend the parameters $e,t$ into $k$ by
\cref{thm:descent} after expressing the nine chart equalities with explicit polynomial
witnesses, and, if the ambient field was enlarged (\cref{rem:psi-hyp}(c)), descend the
existential witnesses from $K'$ to $K$ by model completeness of algebraically closed
fields.  None of (1)--(4) is part of this document.
\end{remark}

\section{Lean architecture}
\label{sec:lean}

\subsection{Library and module layout}

The formalization is a new library \lean{GroupConfig} in the same Lake package
(\lean{lakefile.toml}: a further \lean{[[lean\_lib]]} entry), importing \lean{AclGeom}
for the relative closure, locus, relocation and parameter-field API already proved
there.  Proposed modules, in dependency order (each corresponding to a part of this
document):

\begin{center}
\begin{tabular}{@{}ll@{}}
\toprule
module & content \\
\midrule
\lean{GroupConfig/Setting.lean} & \cref{conv:ud}: the class \lean{UniversalDomain k K} \\
\lean{GroupConfig/Rank.lean} & \cref{def:closures}, \cref{lem:rank} (rank calculus) \\
\lean{GroupConfig/Automorphisms.lean} & \cref{lem:fresh,lem:conj,lem:equiv,lem:orbit} \\
\lean{GroupConfig/Locus.lean} & \cref{def:locus}--\cref{lem:fresh-conj} (loci, generic points, dimension) \\
\lean{GroupConfig/Rref.lean} & \cref{def:rref,lem:rref} (reduced row echelon bases) \\
\lean{GroupConfig/FieldOfDefinition.lean} & \cref{def:fod-subspace}--\cref{thm:weil-fod} \\
\lean{GroupConfig/CanonicalBase.lean} & \cref{def:cb,prop:cb} \\
\lean{GroupConfig/Dcl.lean} & \cref{lem:dcl-galois,cor:dcl-props,def:qr,lem:qr-dcl} \\
\lean{GroupConfig/Codes.lean} & \cref{def:code,lem:code,def:conjset,lem:conjcode} \\
\lean{GroupConfig/Germ/Basic.lean} & \cref{def:variety}--\cref{def:dom} \\
\lean{GroupConfig/Germ/Composition.lean} & \cref{lem:comp,lem:inverse,lem:comp-codes,lem:eval-dcl} \\
\lean{GroupConfig/Germ/Triple.lean} & \cref{lem:triple}, \cref{rem:code-param} \\
\lean{GroupConfig/HrushovskiWeil.lean} & \cref{def:genset}--\cref{thm:hw} \\
\lean{GroupConfig/Bijections.lean} & \cref{def:diff}--\cref{thm:hw-bij} \\
\lean{GroupConfig/Configuration.lean} & \cref{def:gc}--\cref{prop:converse} \\
\lean{GroupConfig/DefinabilityStep.lean} & \cref{lem:defstep,prop:stepI} \\
\lean{GroupConfig/Independence.lean} & \cref{prop:stepII}, \cref{lem:hw-basechange} \\
\lean{GroupConfig/Main.lean} & \cref{thm:main} \\
\lean{GroupConfig/Constructible/Basic.lean} & \cref{def:constructible}, \cref{thm:chevalley}, \cref{cor:constructible-closure} \\
\lean{GroupConfig/Constructible/Generic.lean} & \cref{lem:gen-constructible,lem:genquant} \\
\lean{GroupConfig/Constructible/Group.lean} & \cref{lem:adm-constructible}, \cref{def:constructible-group}, \cref{thm:constructible} \\
\lean{GroupConfig/Descent.lean} & \cref{lem:kpoints-dense,thm:descent} \\
\lean{GroupConfig/Psi.lean} & \cref{prop:psi-gc,cor:psi} (bridge to \lean{AclGeom.QWitness.Psi}) \\
\bottomrule
\end{tabular}
\end{center}

\subsection{Design decisions}

\begin{enumerate}
\item \textbf{The base field is a variable.}  All definitions and theorems are stated for
  an arbitrary algebraically closed base \lean{k} with
  \lean{[Field k] [IsAlgClosed k] [Algebra k K] [IsAlgClosed K]} and the hypothesis
  \lean{Cardinal.aleph0 ≤ Algebra.trdeg k K}.  ``Base change to $\acl(k,e)$'' is then
  instantiation at \lean{k' := ↥(AclGeom.racl k (Set.range e))} with the induced
  algebra structure; \cref{lem:gc-invariance}(2) and \cref{lem:hw-basechange} are the
  lemmas that transport hypotheses across this instantiation.  Warning: this is exactly
  the situation in which the project has met kernel blow-ups (two presentations of the
  same subfield, \lean{F} versus \lean{F.restrictScalars k}); state facts about the new
  base through \lean{AclGeom.mem\_racl\_adjoin\_base\_iff} and
  \lean{AclGeom.racl\_subset\_racl\_base} rather than by rewriting instances.
\item \textbf{Tuples} are \lean{Fin n → K}; concatenation is \lean{Fin.append} (or
  \lean{Sum.elim} with \lean{Fin n ⊕ Fin m}); ``$a\subseteq\acl(B)$'' is
  \lean{∀ i, a i ∈ AclGeom.racl k B}.  Rank is the matroid rank of
  \lean{AlgebraicIndependent.matroid k K} relative to a set, and \cref{lem:rank} is
  proved once from Mathlib's \lean{Matroid} API (\lean{Matroid.eRk}, submodularity,
  closure) together with \lean{AclGeom.rankEq\_iff\_eRk}.
\item \textbf{Loci} use \lean{AclGeom.idealOf} (which is
  \lean{MvPolynomial.vanishingIdeal L \{a\}}) and \lean{Ideal.map (MvPolynomial.map (algebraMap L K))}
  for $P_L(a)$; primality is \lean{AclGeom.isPrime\_map\_of\_isAlgClosed}; dimension facts
  come from \lean{AclGeom.trdeg\_adjoin\_of\_isGeneric} and \lean{AclGeom.isGeneric\_baseChange}.
  \Cref{lem:dimdrop} is new.
\item \textbf{Fields of definition} are defined via RREF of the finite-dimensional
  subspaces $I_{\le d}$ (\lean{MvPolynomial.restrictTotalDegree} gives the subspace of
  polynomials of degree $\le d$; a basis of monomials is \lean{MvPolynomial.basisMonomials}).
  \Cref{lem:rref} (existence/uniqueness of RREF bases) must be proved; Mathlib only has
  the predicate \lean{Matrix.IsReducedRowEchelon}.  Everything else in \cref{sec:fod} is
  linear algebra plus \lean{Module.Flat} for field extensions.
\item \textbf{Germs} are a quotient type of representatives
  (\lean{structure Representative} with a small algebraically closed field, carried as a
  subfield \lean{M : IntermediateField k K} with \lean{IsAlgClosed M} and finite
  \lean{Algebra.trdeg k M}, and a quasi-rational expression), with the relation of
  \cref{def:germ}; \lean{Germ.graph : Ideal (MvPolynomial (Fin n ⊕ Fin n') K)} and
  \lean{Germ.cb : IntermediateField k K} are defined on the quotient via
  \cref{lem:graph}(1).  Evaluation \lean{Germ.eval g y} is defined for \lean{y} with a proof
  that \lean{y} is generic over \lean{racl k g.cb}, using \lean{Classical.choose} and
  \cref{lem:eval}(1),(2).  Composition and inverse follow \cref{lem:comp,lem:inverse}.
\item \textbf{Quasi-rational expressions} are pairs \lean{(e : ℕ) × (Fin n' → RatFunc-like)}; since
  Mathlib's \lean{RatFunc} is univariate, use pairs of polynomials in
  \lean{MvPolynomial (Fin n) L} with a nonvanishing side condition, or
  \lean{FractionRing (MvPolynomial (Fin n) L)}.  In characteristic $0$ the exponent is
  fixed to $0$; the general case uses \lean{frobenius} and \lean{IsPurelyInseparable}.
\item \textbf{Automorphisms} are \lean{K ≃ₐ[k] K}; \cref{lem:conj} is
  \lean{AclGeom.exists\_ambientAutomorphism\_of\_idealOf\_eq} instantiated at the base
  subfield \lean{L} (as a \lean{Field} with \lean{Algebra L K}).  The action of
  \lean{K ≃ₐ[k] K} on ideals, subfields, germs and codes is defined explicitly
  (\lean{Ideal.map}, \lean{IntermediateField.map}).
\item \textbf{Codes} of finite sets of tuples (\cref{def:code}) use
  \lean{Polynomial (MvPolynomial (Fin m) K)} and \lean{Polynomial.roots}; conjugate sets are
  \lean{Finset (Fin m → K)} obtained as orbits via \lean{MulAction.orbit} of
  \lean{K ≃ₐ[L] K}, finite by \cref{lem:orbit}.
\item \textbf{Constructible sets} (\cref{sec:constructibility}) are a new inductive or
  \lean{BooleanSubalgebra.closure} of zero sets in \lean{Set (Fin N → K)}; Chevalley is
  derived from \lean{PrimeSpectrum.isConstructible\_comap\_image} as in the proof of
  \cref{thm:chevalley}.  This part is independent of the germ development and can be
  formalized in parallel.
\item \textbf{The configuration} is a structure \lean{GroupConfiguration k a b c x y z}
  with fields \ref{GC:group}--\ref{GC:rank}; the bridge from \lean{AclGeom.QWitness.Psi}
  is \cref{prop:psi-gc}, which only needs submodularity of the rank of closed
  subfields (Mathlib's \lean{Matroid.eRk\_inter\_add\_eRk\_union\_le} through
  \lean{AclGeom.rankEq\_iff\_eRk}) and the exchange lemma for points
  \lean{AclGeom.pointCl\_exchange} (\lean{AclGeom/Geometry/Equivalence.lean}).
\end{enumerate}

\subsection{Order of work and milestones}

\begin{enumerate}[nosep]
\item Rank calculus, automorphisms, loci and dimension (\cref{sec:ud}); milestone:
  \cref{lem:dim}.
\item RREF, fields of definition, canonical bases (\cref{sec:fod}); milestone:
  \cref{prop:cb}.
\item $\dcl$, quasi-rational expressions, codes (\cref{sec:codes}); milestone:
  \cref{lem:conjcode}.
\item Germs with evaluation, composition, inverses, triples (\cref{sec:germs});
  milestone: \cref{lem:triple}.
\item Hrushovski--Weil (\cref{sec:hw,sec:bij}); milestone: \cref{thm:hw-bij}.
\item Configurations, the definability step, Steps II--III (\cref{sec:gc,sec:proof});
  milestone: \cref{thm:main}.
\item Constructibility and descent (\cref{sec:constructibility,sec:descent}), in
  parallel with 4--6.
\item The $\Psi$ bridge (\cref{sec:application}); milestone: \cref{cor:psi}.
\end{enumerate}

\begin{leannote}[estimated sizes]
Each milestone is a few thousand lines at the project's current density.  The
largest new foundational items, not available in Mathlib or \lean{AclGeom}, are:
RREF existence/uniqueness; the Galois characterisation of $\dcl$ (\cref{lem:dcl-galois});
\cref{lem:dimdrop}; the quotient construction of germs with evaluation
(\cref{lem:eval}); the $K$-points form of Chevalley (\cref{thm:chevalley}).  Everything
else is rank calculus and bookkeeping of base changes.
\end{leannote}


\section*{References}
\begin{enumerate}[label={[\arabic*]},leftmargin=8mm]
\item M.~Bays, \emph{Geometric stability theory}, lecture notes (M\"unster, 2015; revised 2024), \S5--6, and the companion note ``The Hrushovski--Weil theorem''.
\item E.~Bouscaren, \emph{The group configuration --- after E.~Hrushovski}, in: The Model Theory of Groups (Nesin, Pillay, eds.), Notre Dame Math.\ Lectures 11, 1989.
\item D.~Evans, E.~Hrushovski, \emph{Projective planes in algebraically closed fields}, Proc.\ London Math.\ Soc.\ (3) 62 (1991), 1--24.
\item D.~Evans, E.~Hrushovski, \emph{The automorphism group of the combinatorial geometry of an algebraically closed field}, J.\ London Math.\ Soc.\ (2) 52 (1995), 209--225.
\item E.~Hrushovski, \emph{Contributions to stable model theory}, Ph.D.\ thesis, Berkeley 1986, \S3.1.
\item E.~Hrushovski, S.~Rideau-Kikuchi, \emph{Valued fields, metastable groups}, Selecta Math.\ 25 (2019), \S3.4.
\item A.~Pillay, \emph{Geometric stability theory}, Oxford 1996, Chapter 5, \S\S1, 4.
\item B.~Poizat, \emph{Une th\'eorie de Galois imaginaire}, J.\ Symbolic Logic 48 (1983), 1151--1170.
\item A.~Weil, \emph{On algebraic groups of transformations}, Amer.\ J.\ Math.\ 77 (1955), 355--391; \emph{On algebraic groups and homogeneous spaces}, ibid.\ 493--512.
\end{enumerate}

\end{document}
